\documentclass[11pt]{article}
\usepackage[utf8]{inputenc}
\usepackage[T1]{fontenc}
\usepackage[margin=1in]{geometry}
\usepackage{amsmath,amssymb}
\usepackage{graphicx}
\usepackage{hyperref}
\setlength{\parskip}{0.5em}
\setlength{\parindent}{0pt}

\title{第 8 章}
\author{}
\date{}
\begin{document}
\maketitle

\section{数字逻辑电路基础}

8.1 数制与BCD码 8.2 逻辑代数基础
8.1.1 常用数制 8.2.1 基本逻辑运算
8.1.2 几种简单的编码 8.2.2 复合逻辑运算
8.2.3 逻辑电平与正、负逻辑
8.2.4 基本定律和规则
8.2.5 逻辑函数的标准形式
8.2.6 逻辑函数的化简

\begin{figure}[htbp]
  \centering
  \includegraphics[width=\linewidth]{assets/figure_p1_1.png}
\end{figure}

\begin{figure}[htbp]
  \centering
  \includegraphics[width=\linewidth]{assets/figure_p1_2.png}
\end{figure}

\begin{figure}[htbp]
  \centering
  \includegraphics[width=\linewidth]{assets/figure_p1_3.png}
\end{figure}

\begin{figure}[htbp]
  \centering
  \includegraphics[width=\linewidth]{assets/figure_p1_4.png}
\end{figure}

\begin{figure}[htbp]
  \centering
  \includegraphics[width=\linewidth]{assets/figure_p1_5.png}
\end{figure}

\begin{figure}[htbp]
  \centering
  \includegraphics[width=\linewidth]{assets/figure_p1_6.png}
\end{figure}

\begin{figure}[htbp]
  \centering
  \includegraphics[width=\linewidth]{assets/figure_p1_7.png}
\end{figure}

\begin{figure}[htbp]
  \centering
  \includegraphics[width=\linewidth]{assets/figure_p1_8.png}
\end{figure}

\begin{figure}[htbp]
  \centering
  \includegraphics[width=\linewidth]{assets/figure_p1_9.png}
\end{figure}

\begin{figure}[htbp]
  \centering
  \includegraphics[width=\linewidth]{assets/figure_p1_10.png}
\end{figure}

\begin{figure}[htbp]
  \centering
  \includegraphics[width=\linewidth]{assets/figure_p1_11.png}
\end{figure}

\begin{figure}[htbp]
  \centering
  \includegraphics[width=\linewidth]{assets/figure_p1_12.png}
\end{figure}

\begin{figure}[htbp]
  \centering
  \includegraphics[width=\linewidth]{assets/figure_p1_13.png}
\end{figure}

\begin{figure}[htbp]
  \centering
  \includegraphics[width=\linewidth]{assets/figure_p1_14.png}
\end{figure}

\begin{figure}[htbp]
  \centering
  \includegraphics[width=\linewidth]{assets/figure_p1_15.png}
\end{figure}

\begin{figure}[htbp]
  \centering
  \includegraphics[width=\linewidth]{assets/figure_p1_16.png}
\end{figure}

\begin{figure}[htbp]
  \centering
  \includegraphics[width=\linewidth]{assets/figure_p1_17.png}
\end{figure}

\begin{figure}[htbp]
  \centering
  \includegraphics[width=\linewidth]{assets/figure_p1_18.png}
\end{figure}

\begin{figure}[htbp]
  \centering
  \includegraphics[width=\linewidth]{assets/figure_p1_19.png}
\end{figure}

\section{什么是数字逻辑电路？}

两类信号: 模拟信号;数字信号.

在时间上和幅值上均连续的信号称为模拟信号;

在时间上和幅值上均离散的信号称为数字信号.

处理数字信号的电路称为数字电路.

\begin{figure}[htbp]
  \centering
  \includegraphics[width=\linewidth]{assets/figure_p2_1.png}
\end{figure}

\begin{figure}[htbp]
  \centering
  \includegraphics[width=\linewidth]{assets/figure_p2_2.png}
\end{figure}

\begin{figure}[htbp]
  \centering
  \includegraphics[width=\linewidth]{assets/figure_p2_3.png}
\end{figure}

\begin{figure}[htbp]
  \centering
  \includegraphics[width=\linewidth]{assets/figure_p2_4.png}
\end{figure}

\begin{figure}[htbp]
  \centering
  \includegraphics[width=\linewidth]{assets/figure_p2_5.png}
\end{figure}

\begin{figure}[htbp]
  \centering
  \includegraphics[width=\linewidth]{assets/figure_p2_6.png}
\end{figure}

\begin{figure}[htbp]
  \centering
  \includegraphics[width=\linewidth]{assets/figure_p2_7.png}
\end{figure}

\begin{figure}[htbp]
  \centering
  \includegraphics[width=\linewidth]{assets/figure_p2_8.png}
\end{figure}

\begin{figure}[htbp]
  \centering
  \includegraphics[width=\linewidth]{assets/figure_p2_9.png}
\end{figure}

\begin{figure}[htbp]
  \centering
  \includegraphics[width=\linewidth]{assets/figure_p2_10.png}
\end{figure}

\begin{figure}[htbp]
  \centering
  \includegraphics[width=\linewidth]{assets/figure_p2_11.png}
\end{figure}

\begin{figure}[htbp]
  \centering
  \includegraphics[width=\linewidth]{assets/figure_p2_12.png}
\end{figure}

\begin{figure}[htbp]
  \centering
  \includegraphics[width=\linewidth]{assets/figure_p2_13.png}
\end{figure}

\begin{figure}[htbp]
  \centering
  \includegraphics[width=\linewidth]{assets/figure_p2_14.png}
\end{figure}

\begin{figure}[htbp]
  \centering
  \includegraphics[width=\linewidth]{assets/figure_p2_15.png}
\end{figure}

\begin{figure}[htbp]
  \centering
  \includegraphics[width=\linewidth]{assets/figure_p2_16.png}
\end{figure}

\begin{figure}[htbp]
  \centering
  \includegraphics[width=\linewidth]{assets/figure_p2_17.png}
\end{figure}

数字电路特点:

1) 工作信号是二进制表示的二值信号(具有“0”和“1”
两种取值);

2) 电路中器件工作于“开”和“关”两种状态,电路的输

出和输入为逻辑关系;

3) 电路既能进行“代数”运算,也能进行“逻辑”运算;

4) 电路工作可靠,精度高,抗干扰性好.

\begin{figure}[htbp]
  \centering
  \includegraphics[width=\linewidth]{assets/figure_p3_1.png}
\end{figure}

\begin{figure}[htbp]
  \centering
  \includegraphics[width=\linewidth]{assets/figure_p3_2.png}
\end{figure}

\begin{figure}[htbp]
  \centering
  \includegraphics[width=\linewidth]{assets/figure_p3_3.png}
\end{figure}

\begin{figure}[htbp]
  \centering
  \includegraphics[width=\linewidth]{assets/figure_p3_4.png}
\end{figure}

\begin{figure}[htbp]
  \centering
  \includegraphics[width=\linewidth]{assets/figure_p3_5.png}
\end{figure}

\begin{figure}[htbp]
  \centering
  \includegraphics[width=\linewidth]{assets/figure_p3_6.png}
\end{figure}

\begin{figure}[htbp]
  \centering
  \includegraphics[width=\linewidth]{assets/figure_p3_7.png}
\end{figure}

\begin{figure}[htbp]
  \centering
  \includegraphics[width=\linewidth]{assets/figure_p3_8.png}
\end{figure}

\begin{figure}[htbp]
  \centering
  \includegraphics[width=\linewidth]{assets/figure_p3_9.png}
\end{figure}

\begin{figure}[htbp]
  \centering
  \includegraphics[width=\linewidth]{assets/figure_p3_10.png}
\end{figure}

\begin{figure}[htbp]
  \centering
  \includegraphics[width=\linewidth]{assets/figure_p3_11.png}
\end{figure}

\begin{figure}[htbp]
  \centering
  \includegraphics[width=\linewidth]{assets/figure_p3_12.png}
\end{figure}

\begin{figure}[htbp]
  \centering
  \includegraphics[width=\linewidth]{assets/figure_p3_13.png}
\end{figure}

\begin{figure}[htbp]
  \centering
  \includegraphics[width=\linewidth]{assets/figure_p3_14.png}
\end{figure}

\begin{figure}[htbp]
  \centering
  \includegraphics[width=\linewidth]{assets/figure_p3_15.png}
\end{figure}

\begin{figure}[htbp]
  \centering
  \includegraphics[width=\linewidth]{assets/figure_p3_16.png}
\end{figure}

\begin{figure}[htbp]
  \centering
  \includegraphics[width=\linewidth]{assets/figure_p3_17.png}
\end{figure}

8.1 数制与\textbf{BCD}码

所谓“数制”，指进位计数制，即用进位的方法来计
数.

数制包括计数符号（数码）和进位规则两个方面。

常用数制有十进制、十二进制、十六进制、六十进
制等。

\begin{figure}[htbp]
  \centering
  \includegraphics[width=\linewidth]{assets/figure_p4_1.png}
\end{figure}

\begin{figure}[htbp]
  \centering
  \includegraphics[width=\linewidth]{assets/figure_p4_2.png}
\end{figure}

\begin{figure}[htbp]
  \centering
  \includegraphics[width=\linewidth]{assets/figure_p4_3.png}
\end{figure}

\begin{figure}[htbp]
  \centering
  \includegraphics[width=\linewidth]{assets/figure_p4_4.png}
\end{figure}

\begin{figure}[htbp]
  \centering
  \includegraphics[width=\linewidth]{assets/figure_p4_5.png}
\end{figure}

\begin{figure}[htbp]
  \centering
  \includegraphics[width=\linewidth]{assets/figure_p4_6.png}
\end{figure}

\begin{figure}[htbp]
  \centering
  \includegraphics[width=\linewidth]{assets/figure_p4_7.png}
\end{figure}

\begin{figure}[htbp]
  \centering
  \includegraphics[width=\linewidth]{assets/figure_p4_8.png}
\end{figure}

\begin{figure}[htbp]
  \centering
  \includegraphics[width=\linewidth]{assets/figure_p4_9.png}
\end{figure}

\begin{figure}[htbp]
  \centering
  \includegraphics[width=\linewidth]{assets/figure_p4_10.png}
\end{figure}

\begin{figure}[htbp]
  \centering
  \includegraphics[width=\linewidth]{assets/figure_p4_11.png}
\end{figure}

\begin{figure}[htbp]
  \centering
  \includegraphics[width=\linewidth]{assets/figure_p4_12.png}
\end{figure}

\begin{figure}[htbp]
  \centering
  \includegraphics[width=\linewidth]{assets/figure_p4_13.png}
\end{figure}

\begin{figure}[htbp]
  \centering
  \includegraphics[width=\linewidth]{assets/figure_p4_14.png}
\end{figure}

\begin{figure}[htbp]
  \centering
  \includegraphics[width=\linewidth]{assets/figure_p4_15.png}
\end{figure}

\begin{figure}[htbp]
  \centering
  \includegraphics[width=\linewidth]{assets/figure_p4_16.png}
\end{figure}

\begin{figure}[htbp]
  \centering
  \includegraphics[width=\linewidth]{assets/figure_p4_17.png}
\end{figure}

8.1.1 常用数制

1. 十进制

(1) 计数符号: 0, 1, 2, 3, 4, 5, 6, 7, 8, 9.

(2)进位规则: 逢十进一.

例: $\textbf{1987.45=1}\times \textbf{10}\textsuperscript{\textbf{3}}\textbf{+9}\times \textbf{10}\textsuperscript{\textbf{2}}$ \textbf{+ $8}\times \textbf{10}\textsuperscript{\textbf{1}}\textbf{+ 7}\times \textbf{10}\textsuperscript{\textbf{0}}$
$\textbf{+4}\times \textbf{10}\textsuperscript{\textbf{-1}}\textbf{+5}\times \textbf{10}\textsuperscript{\textbf{-2}}$

(3) 十进制数按权展开式

\begin{figure}[htbp]
  \centering
  \includegraphics[width=\linewidth]{assets/figure_p5_1.png}
\end{figure}

\begin{figure}[htbp]
  \centering
  \includegraphics[width=\linewidth]{assets/figure_p5_2.png}
\end{figure}

\begin{figure}[htbp]
  \centering
  \includegraphics[width=\linewidth]{assets/figure_p5_3.png}
\end{figure}

\begin{figure}[htbp]
  \centering
  \includegraphics[width=\linewidth]{assets/figure_p5_4.png}
\end{figure}

\begin{figure}[htbp]
  \centering
  \includegraphics[width=\linewidth]{assets/figure_p5_5.png}
\end{figure}

\begin{figure}[htbp]
  \centering
  \includegraphics[width=\linewidth]{assets/figure_p5_6.png}
\end{figure}

\begin{figure}[htbp]
  \centering
  \includegraphics[width=\linewidth]{assets/figure_p5_7.png}
\end{figure}

\begin{figure}[htbp]
  \centering
  \includegraphics[width=\linewidth]{assets/figure_p5_8.png}
\end{figure}

\begin{figure}[htbp]
  \centering
  \includegraphics[width=\linewidth]{assets/figure_p5_9.png}
\end{figure}

\begin{figure}[htbp]
  \centering
  \includegraphics[width=\linewidth]{assets/figure_p5_10.png}
\end{figure}

\begin{figure}[htbp]
  \centering
  \includegraphics[width=\linewidth]{assets/figure_p5_11.png}
\end{figure}

\begin{figure}[htbp]
  \centering
  \includegraphics[width=\linewidth]{assets/figure_p5_12.png}
\end{figure}

\begin{figure}[htbp]
  \centering
  \includegraphics[width=\linewidth]{assets/figure_p5_13.png}
\end{figure}

\begin{figure}[htbp]
  \centering
  \includegraphics[width=\linewidth]{assets/figure_p5_14.png}
\end{figure}

\begin{figure}[htbp]
  \centering
  \includegraphics[width=\linewidth]{assets/figure_p5_15.png}
\end{figure}

\begin{figure}[htbp]
  \centering
  \includegraphics[width=\linewidth]{assets/figure_p5_16.png}
\end{figure}

\begin{figure}[htbp]
  \centering
  \includegraphics[width=\linewidth]{assets/figure_p5_17.png}
\end{figure}

\textbf{n}\textbf{1}
\textbf{i}

\textbf{(N)}  \textbf{a} \textbf{10}

\textbf{10 i}
\textbf{i}\textbf{m}

系数 权

2. 二进制

(1) 计数符号: 0, 1 .

(2)进位规则: 逢二进一.

n 1
i

(3) 二进制数按权展开式 (N )   a  2

2 i
i  m

\begin{figure}[htbp]
  \centering
  \includegraphics[width=\linewidth]{assets/figure_p6_1.png}
\end{figure}

\begin{figure}[htbp]
  \centering
  \includegraphics[width=\linewidth]{assets/figure_p6_2.png}
\end{figure}

\begin{figure}[htbp]
  \centering
  \includegraphics[width=\linewidth]{assets/figure_p6_3.png}
\end{figure}

\begin{figure}[htbp]
  \centering
  \includegraphics[width=\linewidth]{assets/figure_p6_4.png}
\end{figure}

\begin{figure}[htbp]
  \centering
  \includegraphics[width=\linewidth]{assets/figure_p6_5.png}
\end{figure}

\begin{figure}[htbp]
  \centering
  \includegraphics[width=\linewidth]{assets/figure_p6_6.png}
\end{figure}

\begin{figure}[htbp]
  \centering
  \includegraphics[width=\linewidth]{assets/figure_p6_7.png}
\end{figure}

\begin{figure}[htbp]
  \centering
  \includegraphics[width=\linewidth]{assets/figure_p6_8.png}
\end{figure}

\begin{figure}[htbp]
  \centering
  \includegraphics[width=\linewidth]{assets/figure_p6_9.png}
\end{figure}

\begin{figure}[htbp]
  \centering
  \includegraphics[width=\linewidth]{assets/figure_p6_10.png}
\end{figure}

\begin{figure}[htbp]
  \centering
  \includegraphics[width=\linewidth]{assets/figure_p6_11.png}
\end{figure}

\begin{figure}[htbp]
  \centering
  \includegraphics[width=\linewidth]{assets/figure_p6_12.png}
\end{figure}

\begin{figure}[htbp]
  \centering
  \includegraphics[width=\linewidth]{assets/figure_p6_13.png}
\end{figure}

\begin{figure}[htbp]
  \centering
  \includegraphics[width=\linewidth]{assets/figure_p6_14.png}
\end{figure}

\begin{figure}[htbp]
  \centering
  \includegraphics[width=\linewidth]{assets/figure_p6_15.png}
\end{figure}

\begin{figure}[htbp]
  \centering
  \includegraphics[width=\linewidth]{assets/figure_p6_16.png}
\end{figure}

\begin{figure}[htbp]
  \centering
  \includegraphics[width=\linewidth]{assets/figure_p6_17.png}
\end{figure}

数字电路中采用二进制的原因：

1）数字装置简单可靠；

2）二进制数运算规则简单；

3）数字电路既可以进行算术运算，也可以进行逻辑运算.

3.十六进制和八进制

十六进制数计数符号: 0,1, .,9,\textbf{A,B,C,D,E,F}.

十六进制数进位规则: 逢十六进一.

n1
i

按权展开式： (N)  a 16

16 i
im

\begin{figure}[htbp]
  \centering
  \includegraphics[width=\linewidth]{assets/figure_p7_1.png}
\end{figure}

\begin{figure}[htbp]
  \centering
  \includegraphics[width=\linewidth]{assets/figure_p7_2.png}
\end{figure}

\begin{figure}[htbp]
  \centering
  \includegraphics[width=\linewidth]{assets/figure_p7_3.png}
\end{figure}

\begin{figure}[htbp]
  \centering
  \includegraphics[width=\linewidth]{assets/figure_p7_4.png}
\end{figure}

\begin{figure}[htbp]
  \centering
  \includegraphics[width=\linewidth]{assets/figure_p7_5.png}
\end{figure}

\begin{figure}[htbp]
  \centering
  \includegraphics[width=\linewidth]{assets/figure_p7_6.png}
\end{figure}

\begin{figure}[htbp]
  \centering
  \includegraphics[width=\linewidth]{assets/figure_p7_7.png}
\end{figure}

\begin{figure}[htbp]
  \centering
  \includegraphics[width=\linewidth]{assets/figure_p7_8.png}
\end{figure}

\begin{figure}[htbp]
  \centering
  \includegraphics[width=\linewidth]{assets/figure_p7_9.png}
\end{figure}

\begin{figure}[htbp]
  \centering
  \includegraphics[width=\linewidth]{assets/figure_p7_10.png}
\end{figure}

\begin{figure}[htbp]
  \centering
  \includegraphics[width=\linewidth]{assets/figure_p7_11.png}
\end{figure}

\begin{figure}[htbp]
  \centering
  \includegraphics[width=\linewidth]{assets/figure_p7_12.png}
\end{figure}

\begin{figure}[htbp]
  \centering
  \includegraphics[width=\linewidth]{assets/figure_p7_13.png}
\end{figure}

\begin{figure}[htbp]
  \centering
  \includegraphics[width=\linewidth]{assets/figure_p7_14.png}
\end{figure}

\begin{figure}[htbp]
  \centering
  \includegraphics[width=\linewidth]{assets/figure_p7_15.png}
\end{figure}

\begin{figure}[htbp]
  \centering
  \includegraphics[width=\linewidth]{assets/figure_p7_16.png}
\end{figure}

\begin{figure}[htbp]
  \centering
  \includegraphics[width=\linewidth]{assets/figure_p7_17.png}
\end{figure}

例:

1 0 1 2

(6D.4B)  616  D16  416  B16

16
1 0 1 2

 616 1316 416 1116

八进制数计数符号: 0,1, . . .6,7.
八进制数进位规则: 逢八进一.

\textbf{n}\textbf{1}
\textbf{i}

按权展开式： \textbf{(N)}  \textbf{a} \textbf{8}

\textbf{8 i}
\textbf{i}\textbf{m}

\begin{figure}[htbp]
  \centering
  \includegraphics[width=\linewidth]{assets/figure_p8_1.png}
\end{figure}

\begin{figure}[htbp]
  \centering
  \includegraphics[width=\linewidth]{assets/figure_p8_2.png}
\end{figure}

\begin{figure}[htbp]
  \centering
  \includegraphics[width=\linewidth]{assets/figure_p8_3.png}
\end{figure}

\begin{figure}[htbp]
  \centering
  \includegraphics[width=\linewidth]{assets/figure_p8_4.png}
\end{figure}

\begin{figure}[htbp]
  \centering
  \includegraphics[width=\linewidth]{assets/figure_p8_5.png}
\end{figure}

\begin{figure}[htbp]
  \centering
  \includegraphics[width=\linewidth]{assets/figure_p8_6.png}
\end{figure}

\begin{figure}[htbp]
  \centering
  \includegraphics[width=\linewidth]{assets/figure_p8_7.png}
\end{figure}

\begin{figure}[htbp]
  \centering
  \includegraphics[width=\linewidth]{assets/figure_p8_8.png}
\end{figure}

\begin{figure}[htbp]
  \centering
  \includegraphics[width=\linewidth]{assets/figure_p8_9.png}
\end{figure}

\begin{figure}[htbp]
  \centering
  \includegraphics[width=\linewidth]{assets/figure_p8_10.png}
\end{figure}

\begin{figure}[htbp]
  \centering
  \includegraphics[width=\linewidth]{assets/figure_p8_11.png}
\end{figure}

\begin{figure}[htbp]
  \centering
  \includegraphics[width=\linewidth]{assets/figure_p8_12.png}
\end{figure}

\begin{figure}[htbp]
  \centering
  \includegraphics[width=\linewidth]{assets/figure_p8_13.png}
\end{figure}

\begin{figure}[htbp]
  \centering
  \includegraphics[width=\linewidth]{assets/figure_p8_14.png}
\end{figure}

\begin{figure}[htbp]
  \centering
  \includegraphics[width=\linewidth]{assets/figure_p8_15.png}
\end{figure}

\begin{figure}[htbp]
  \centering
  \includegraphics[width=\linewidth]{assets/figure_p8_16.png}
\end{figure}

\begin{figure}[htbp]
  \centering
  \includegraphics[width=\linewidth]{assets/figure_p8_17.png}
\end{figure}

1 0 1 2

\subsection{例: ( 63 . 45 )  6  8  3  8  4  8  5  8}

8

4. 二进制数与十进制数之间的转换

(1)二进制数转换为十进制数(按权展开法)

\textbf{3 1 0} \textbf{1} \textbf{3}

例： \textbf{(1011.101)}  \textbf{1}\textbf{2} \textbf{1}\textbf{2} \textbf{1}\textbf{2} \textbf{1}\textbf{2} \textbf{1}\textbf{2}

\textbf{2}

 \textbf{8} \textbf{2}\textbf{1} \textbf{0.5} \textbf{0.125}

\begin{figure}[htbp]
  \centering
  \includegraphics[width=\linewidth]{assets/figure_p9_1.png}
\end{figure}

\begin{figure}[htbp]
  \centering
  \includegraphics[width=\linewidth]{assets/figure_p9_2.png}
\end{figure}

\begin{figure}[htbp]
  \centering
  \includegraphics[width=\linewidth]{assets/figure_p9_3.png}
\end{figure}

\begin{figure}[htbp]
  \centering
  \includegraphics[width=\linewidth]{assets/figure_p9_4.png}
\end{figure}

\begin{figure}[htbp]
  \centering
  \includegraphics[width=\linewidth]{assets/figure_p9_5.png}
\end{figure}

\begin{figure}[htbp]
  \centering
  \includegraphics[width=\linewidth]{assets/figure_p9_6.png}
\end{figure}

\begin{figure}[htbp]
  \centering
  \includegraphics[width=\linewidth]{assets/figure_p9_7.png}
\end{figure}

\begin{figure}[htbp]
  \centering
  \includegraphics[width=\linewidth]{assets/figure_p9_8.png}
\end{figure}

\begin{figure}[htbp]
  \centering
  \includegraphics[width=\linewidth]{assets/figure_p9_9.png}
\end{figure}

\begin{figure}[htbp]
  \centering
  \includegraphics[width=\linewidth]{assets/figure_p9_10.png}
\end{figure}

\begin{figure}[htbp]
  \centering
  \includegraphics[width=\linewidth]{assets/figure_p9_11.png}
\end{figure}

\begin{figure}[htbp]
  \centering
  \includegraphics[width=\linewidth]{assets/figure_p9_12.png}
\end{figure}

\begin{figure}[htbp]
  \centering
  \includegraphics[width=\linewidth]{assets/figure_p9_13.png}
\end{figure}

\begin{figure}[htbp]
  \centering
  \includegraphics[width=\linewidth]{assets/figure_p9_14.png}
\end{figure}

\begin{figure}[htbp]
  \centering
  \includegraphics[width=\linewidth]{assets/figure_p9_15.png}
\end{figure}

\begin{figure}[htbp]
  \centering
  \includegraphics[width=\linewidth]{assets/figure_p9_16.png}
\end{figure}

\begin{figure}[htbp]
  \centering
  \includegraphics[width=\linewidth]{assets/figure_p9_17.png}
\end{figure}

（2）十进制数转换为二进制数(提取2的幂法)

例： (45.5) 328410.5

10

5 4 3 2 1 0 -1

12 02 12 12 02 12 12

 (101101.1)

2

\textbf{•} 数制转换还可以采用基数连乘、连除等方法.

\begin{figure}[htbp]
  \centering
  \includegraphics[width=\linewidth]{assets/figure_p10_1.png}
\end{figure}

\begin{figure}[htbp]
  \centering
  \includegraphics[width=\linewidth]{assets/figure_p10_2.png}
\end{figure}

\begin{figure}[htbp]
  \centering
  \includegraphics[width=\linewidth]{assets/figure_p10_3.png}
\end{figure}

\begin{figure}[htbp]
  \centering
  \includegraphics[width=\linewidth]{assets/figure_p10_4.png}
\end{figure}

\begin{figure}[htbp]
  \centering
  \includegraphics[width=\linewidth]{assets/figure_p10_5.png}
\end{figure}

\begin{figure}[htbp]
  \centering
  \includegraphics[width=\linewidth]{assets/figure_p10_6.png}
\end{figure}

\begin{figure}[htbp]
  \centering
  \includegraphics[width=\linewidth]{assets/figure_p10_7.png}
\end{figure}

\begin{figure}[htbp]
  \centering
  \includegraphics[width=\linewidth]{assets/figure_p10_8.png}
\end{figure}

\begin{figure}[htbp]
  \centering
  \includegraphics[width=\linewidth]{assets/figure_p10_9.png}
\end{figure}

\begin{figure}[htbp]
  \centering
  \includegraphics[width=\linewidth]{assets/figure_p10_10.png}
\end{figure}

\begin{figure}[htbp]
  \centering
  \includegraphics[width=\linewidth]{assets/figure_p10_11.png}
\end{figure}

\begin{figure}[htbp]
  \centering
  \includegraphics[width=\linewidth]{assets/figure_p10_12.png}
\end{figure}

\begin{figure}[htbp]
  \centering
  \includegraphics[width=\linewidth]{assets/figure_p10_13.png}
\end{figure}

\begin{figure}[htbp]
  \centering
  \includegraphics[width=\linewidth]{assets/figure_p10_14.png}
\end{figure}

\begin{figure}[htbp]
  \centering
  \includegraphics[width=\linewidth]{assets/figure_p10_15.png}
\end{figure}

\begin{figure}[htbp]
  \centering
  \includegraphics[width=\linewidth]{assets/figure_p10_16.png}
\end{figure}

\begin{figure}[htbp]
  \centering
  \includegraphics[width=\linewidth]{assets/figure_p10_17.png}
\end{figure}

5. 二进制算术运算

（1）二进制算术运算特点

算术运算： 1：和十进制算数运算的规则相同
2：逢二进一
特 点： 加、减、乘、除全部可以用移位和
相加这两种操作实现。

电路结构简化
所以数字电路中普遍采用二进制算数运算

\begin{figure}[htbp]
  \centering
  \includegraphics[width=\linewidth]{assets/figure_p11_1.png}
\end{figure}

\begin{figure}[htbp]
  \centering
  \includegraphics[width=\linewidth]{assets/figure_p11_2.png}
\end{figure}

\begin{figure}[htbp]
  \centering
  \includegraphics[width=\linewidth]{assets/figure_p11_3.png}
\end{figure}

\begin{figure}[htbp]
  \centering
  \includegraphics[width=\linewidth]{assets/figure_p11_4.png}
\end{figure}

\begin{figure}[htbp]
  \centering
  \includegraphics[width=\linewidth]{assets/figure_p11_5.png}
\end{figure}

\begin{figure}[htbp]
  \centering
  \includegraphics[width=\linewidth]{assets/figure_p11_6.png}
\end{figure}

\begin{figure}[htbp]
  \centering
  \includegraphics[width=\linewidth]{assets/figure_p11_7.png}
\end{figure}

\begin{figure}[htbp]
  \centering
  \includegraphics[width=\linewidth]{assets/figure_p11_8.png}
\end{figure}

\begin{figure}[htbp]
  \centering
  \includegraphics[width=\linewidth]{assets/figure_p11_9.png}
\end{figure}

\begin{figure}[htbp]
  \centering
  \includegraphics[width=\linewidth]{assets/figure_p11_10.png}
\end{figure}

\begin{figure}[htbp]
  \centering
  \includegraphics[width=\linewidth]{assets/figure_p11_11.png}
\end{figure}

\begin{figure}[htbp]
  \centering
  \includegraphics[width=\linewidth]{assets/figure_p11_12.png}
\end{figure}

\begin{figure}[htbp]
  \centering
  \includegraphics[width=\linewidth]{assets/figure_p11_13.png}
\end{figure}

\begin{figure}[htbp]
  \centering
  \includegraphics[width=\linewidth]{assets/figure_p11_14.png}
\end{figure}

\begin{figure}[htbp]
  \centering
  \includegraphics[width=\linewidth]{assets/figure_p11_15.png}
\end{figure}

\begin{figure}[htbp]
  \centering
  \includegraphics[width=\linewidth]{assets/figure_p11_16.png}
\end{figure}

\begin{figure}[htbp]
  \centering
  \includegraphics[width=\linewidth]{assets/figure_p11_17.png}
\end{figure}

（2）反码、补码和补码运算

——带符号二进制数的表示及运算

计算机中的符号数可表示为：

符号位\textbf{+}真值

机器数

“\textbf{0}” 表示正
“\textbf{1}” 表示负

\begin{figure}[htbp]
  \centering
  \includegraphics[width=\linewidth]{assets/figure_p12_1.png}
\end{figure}

\begin{figure}[htbp]
  \centering
  \includegraphics[width=\linewidth]{assets/figure_p12_2.png}
\end{figure}

\begin{figure}[htbp]
  \centering
  \includegraphics[width=\linewidth]{assets/figure_p12_3.png}
\end{figure}

\begin{figure}[htbp]
  \centering
  \includegraphics[width=\linewidth]{assets/figure_p12_4.png}
\end{figure}

\begin{figure}[htbp]
  \centering
  \includegraphics[width=\linewidth]{assets/figure_p12_5.png}
\end{figure}

\begin{figure}[htbp]
  \centering
  \includegraphics[width=\linewidth]{assets/figure_p12_6.png}
\end{figure}

\begin{figure}[htbp]
  \centering
  \includegraphics[width=\linewidth]{assets/figure_p12_7.png}
\end{figure}

\begin{figure}[htbp]
  \centering
  \includegraphics[width=\linewidth]{assets/figure_p12_8.png}
\end{figure}

\begin{figure}[htbp]
  \centering
  \includegraphics[width=\linewidth]{assets/figure_p12_9.png}
\end{figure}

\begin{figure}[htbp]
  \centering
  \includegraphics[width=\linewidth]{assets/figure_p12_10.png}
\end{figure}

\begin{figure}[htbp]
  \centering
  \includegraphics[width=\linewidth]{assets/figure_p12_11.png}
\end{figure}

\begin{figure}[htbp]
  \centering
  \includegraphics[width=\linewidth]{assets/figure_p12_12.png}
\end{figure}

\begin{figure}[htbp]
  \centering
  \includegraphics[width=\linewidth]{assets/figure_p12_13.png}
\end{figure}

\begin{figure}[htbp]
  \centering
  \includegraphics[width=\linewidth]{assets/figure_p12_14.png}
\end{figure}

\begin{figure}[htbp]
  \centering
  \includegraphics[width=\linewidth]{assets/figure_p12_15.png}
\end{figure}

\begin{figure}[htbp]
  \centering
  \includegraphics[width=\linewidth]{assets/figure_p12_16.png}
\end{figure}

\begin{figure}[htbp]
  \centering
  \includegraphics[width=\linewidth]{assets/figure_p12_17.png}
\end{figure}

• 原码:真值X的原码记为[X]\textsubscript{真，}在原码表示法中不论
数的正负，数值部分均保持原真值不变。
• \textbf{True}
• 反码:真值X的反码记为[X]\textsubscript{反。}正数的反码同原码。
负数的反码的数值部分为真值的各位按位取反。
• \textbf{1’s complement}
• 补码：真值X的补码记为[X]\textsubscript{补。}正数的补码同原码。
负数的补码的数值部分为真值的各位按位取反加1。
• \textbf{2’s complement}

\begin{figure}[htbp]
  \centering
  \includegraphics[width=\linewidth]{assets/figure_p13_1.png}
\end{figure}

\begin{figure}[htbp]
  \centering
  \includegraphics[width=\linewidth]{assets/figure_p13_2.png}
\end{figure}

\begin{figure}[htbp]
  \centering
  \includegraphics[width=\linewidth]{assets/figure_p13_3.png}
\end{figure}

\begin{figure}[htbp]
  \centering
  \includegraphics[width=\linewidth]{assets/figure_p13_4.png}
\end{figure}

\begin{figure}[htbp]
  \centering
  \includegraphics[width=\linewidth]{assets/figure_p13_5.png}
\end{figure}

\begin{figure}[htbp]
  \centering
  \includegraphics[width=\linewidth]{assets/figure_p13_6.png}
\end{figure}

\begin{figure}[htbp]
  \centering
  \includegraphics[width=\linewidth]{assets/figure_p13_7.png}
\end{figure}

\begin{figure}[htbp]
  \centering
  \includegraphics[width=\linewidth]{assets/figure_p13_8.png}
\end{figure}

\begin{figure}[htbp]
  \centering
  \includegraphics[width=\linewidth]{assets/figure_p13_9.png}
\end{figure}

\begin{figure}[htbp]
  \centering
  \includegraphics[width=\linewidth]{assets/figure_p13_10.png}
\end{figure}

\begin{figure}[htbp]
  \centering
  \includegraphics[width=\linewidth]{assets/figure_p13_11.png}
\end{figure}

\begin{figure}[htbp]
  \centering
  \includegraphics[width=\linewidth]{assets/figure_p13_12.png}
\end{figure}

\begin{figure}[htbp]
  \centering
  \includegraphics[width=\linewidth]{assets/figure_p13_13.png}
\end{figure}

\begin{figure}[htbp]
  \centering
  \includegraphics[width=\linewidth]{assets/figure_p13_14.png}
\end{figure}

\begin{figure}[htbp]
  \centering
  \includegraphics[width=\linewidth]{assets/figure_p13_15.png}
\end{figure}

\begin{figure}[htbp]
  \centering
  \includegraphics[width=\linewidth]{assets/figure_p13_16.png}
\end{figure}

\begin{figure}[htbp]
  \centering
  \includegraphics[width=\linewidth]{assets/figure_p13_17.png}
\end{figure}

\begin{figure}[htbp]
  \centering
  \includegraphics[width=\linewidth]{assets/figure_p13_18.png}
\end{figure}

\section{原码：}

• 最高位为符号位，用“0”表示正，用“1”表
示负；其余为真值部分

• 优点： 真值和其原码表示之间的对应关系简
单，容易理解
• 缺点： 计算机中用原码进行加减运算比较困
难，0的表示不唯一

\textbf{8}位数\textbf{0}的原码：\textbf{+0=0 0000000}
\textbf{-0=1 0000000}

\begin{figure}[htbp]
  \centering
  \includegraphics[width=\linewidth]{assets/figure_p14_1.png}
\end{figure}

\begin{figure}[htbp]
  \centering
  \includegraphics[width=\linewidth]{assets/figure_p14_2.png}
\end{figure}

\begin{figure}[htbp]
  \centering
  \includegraphics[width=\linewidth]{assets/figure_p14_3.png}
\end{figure}

\begin{figure}[htbp]
  \centering
  \includegraphics[width=\linewidth]{assets/figure_p14_4.png}
\end{figure}

\begin{figure}[htbp]
  \centering
  \includegraphics[width=\linewidth]{assets/figure_p14_5.png}
\end{figure}

\begin{figure}[htbp]
  \centering
  \includegraphics[width=\linewidth]{assets/figure_p14_6.png}
\end{figure}

\begin{figure}[htbp]
  \centering
  \includegraphics[width=\linewidth]{assets/figure_p14_7.png}
\end{figure}

\begin{figure}[htbp]
  \centering
  \includegraphics[width=\linewidth]{assets/figure_p14_8.png}
\end{figure}

\begin{figure}[htbp]
  \centering
  \includegraphics[width=\linewidth]{assets/figure_p14_9.png}
\end{figure}

\begin{figure}[htbp]
  \centering
  \includegraphics[width=\linewidth]{assets/figure_p14_10.png}
\end{figure}

\begin{figure}[htbp]
  \centering
  \includegraphics[width=\linewidth]{assets/figure_p14_11.png}
\end{figure}

\begin{figure}[htbp]
  \centering
  \includegraphics[width=\linewidth]{assets/figure_p14_12.png}
\end{figure}

\begin{figure}[htbp]
  \centering
  \includegraphics[width=\linewidth]{assets/figure_p14_13.png}
\end{figure}

\begin{figure}[htbp]
  \centering
  \includegraphics[width=\linewidth]{assets/figure_p14_14.png}
\end{figure}

\begin{figure}[htbp]
  \centering
  \includegraphics[width=\linewidth]{assets/figure_p14_15.png}
\end{figure}

\begin{figure}[htbp]
  \centering
  \includegraphics[width=\linewidth]{assets/figure_p14_16.png}
\end{figure}

\begin{figure}[htbp]
  \centering
  \includegraphics[width=\linewidth]{assets/figure_p14_17.png}
\end{figure}

\section{反码：}

对一个机器数\textbf{X}：

• 若\textbf{X>0} ，则 \textbf{[X]}\textsubscript{反}\textbf{=[X]}\textsubscript{原}
• 若\textbf{X<0}， 则 \textbf{[X]}\textsubscript{反}\textbf{=} 对应原码的符号位不变，数值部分
按位求反
\textbf{[}例\textbf{]}：

\textbf{[+0]}\textsubscript{反}\textbf{=00000000}

\textbf{X= -52 = -0110100}

\textbf{[-0]}\textsubscript{反}\textbf{=11111111}

\textbf{[X]}\textsubscript{原}\textbf{=1 0110100}

\textbf{[X] =1 1001011} 即：数\textbf{0}的反码也不唯一

反

\begin{figure}[htbp]
  \centering
  \includegraphics[width=\linewidth]{assets/figure_p15_1.png}
\end{figure}

\begin{figure}[htbp]
  \centering
  \includegraphics[width=\linewidth]{assets/figure_p15_2.png}
\end{figure}

\begin{figure}[htbp]
  \centering
  \includegraphics[width=\linewidth]{assets/figure_p15_3.png}
\end{figure}

\begin{figure}[htbp]
  \centering
  \includegraphics[width=\linewidth]{assets/figure_p15_4.png}
\end{figure}

\begin{figure}[htbp]
  \centering
  \includegraphics[width=\linewidth]{assets/figure_p15_5.png}
\end{figure}

\begin{figure}[htbp]
  \centering
  \includegraphics[width=\linewidth]{assets/figure_p15_6.png}
\end{figure}

\begin{figure}[htbp]
  \centering
  \includegraphics[width=\linewidth]{assets/figure_p15_7.png}
\end{figure}

\begin{figure}[htbp]
  \centering
  \includegraphics[width=\linewidth]{assets/figure_p15_8.png}
\end{figure}

\begin{figure}[htbp]
  \centering
  \includegraphics[width=\linewidth]{assets/figure_p15_9.png}
\end{figure}

\begin{figure}[htbp]
  \centering
  \includegraphics[width=\linewidth]{assets/figure_p15_10.png}
\end{figure}

\begin{figure}[htbp]
  \centering
  \includegraphics[width=\linewidth]{assets/figure_p15_11.png}
\end{figure}

\begin{figure}[htbp]
  \centering
  \includegraphics[width=\linewidth]{assets/figure_p15_12.png}
\end{figure}

\begin{figure}[htbp]
  \centering
  \includegraphics[width=\linewidth]{assets/figure_p15_13.png}
\end{figure}

\begin{figure}[htbp]
  \centering
  \includegraphics[width=\linewidth]{assets/figure_p15_14.png}
\end{figure}

\begin{figure}[htbp]
  \centering
  \includegraphics[width=\linewidth]{assets/figure_p15_15.png}
\end{figure}

\begin{figure}[htbp]
  \centering
  \includegraphics[width=\linewidth]{assets/figure_p15_16.png}
\end{figure}

\begin{figure}[htbp]
  \centering
  \includegraphics[width=\linewidth]{assets/figure_p15_17.png}
\end{figure}

\section{补码：}

若\textbf{X>0}， 则\textbf{[X]}\textsubscript{补}\textbf{= [X]}\textsubscript{反}\textbf{= [X]}\textsubscript{原}

若\textbf{X<0}， 则\textbf{[X]}\textsubscript{补}\textbf{= [X]}\textsubscript{反}\textbf{+1}

\textbf{[}例\textbf{]}：
\textbf{X= –52= – 0110100}
\textbf{[X]}原\textbf{=10110100}
\textbf{[X]}反\textbf{=11001011}
\textbf{[X]}补\textbf{= [X]}反\textbf{+1=11001100}

\begin{figure}[htbp]
  \centering
  \includegraphics[width=\linewidth]{assets/figure_p16_1.png}
\end{figure}

\begin{figure}[htbp]
  \centering
  \includegraphics[width=\linewidth]{assets/figure_p16_2.png}
\end{figure}

\begin{figure}[htbp]
  \centering
  \includegraphics[width=\linewidth]{assets/figure_p16_3.png}
\end{figure}

\begin{figure}[htbp]
  \centering
  \includegraphics[width=\linewidth]{assets/figure_p16_4.png}
\end{figure}

\begin{figure}[htbp]
  \centering
  \includegraphics[width=\linewidth]{assets/figure_p16_5.png}
\end{figure}

\begin{figure}[htbp]
  \centering
  \includegraphics[width=\linewidth]{assets/figure_p16_6.png}
\end{figure}

\begin{figure}[htbp]
  \centering
  \includegraphics[width=\linewidth]{assets/figure_p16_7.png}
\end{figure}

\begin{figure}[htbp]
  \centering
  \includegraphics[width=\linewidth]{assets/figure_p16_8.png}
\end{figure}

\begin{figure}[htbp]
  \centering
  \includegraphics[width=\linewidth]{assets/figure_p16_9.png}
\end{figure}

\begin{figure}[htbp]
  \centering
  \includegraphics[width=\linewidth]{assets/figure_p16_10.png}
\end{figure}

\begin{figure}[htbp]
  \centering
  \includegraphics[width=\linewidth]{assets/figure_p16_11.png}
\end{figure}

\begin{figure}[htbp]
  \centering
  \includegraphics[width=\linewidth]{assets/figure_p16_12.png}
\end{figure}

\begin{figure}[htbp]
  \centering
  \includegraphics[width=\linewidth]{assets/figure_p16_13.png}
\end{figure}

\begin{figure}[htbp]
  \centering
  \includegraphics[width=\linewidth]{assets/figure_p16_14.png}
\end{figure}

\begin{figure}[htbp]
  \centering
  \includegraphics[width=\linewidth]{assets/figure_p16_15.png}
\end{figure}

\begin{figure}[htbp]
  \centering
  \includegraphics[width=\linewidth]{assets/figure_p16_16.png}
\end{figure}

\begin{figure}[htbp]
  \centering
  \includegraphics[width=\linewidth]{assets/figure_p16_17.png}
\end{figure}

\section{0 的补码：表示唯一}

[+0]\textsubscript{补}= [+0]\textsubscript{原}=00000000
[-0]\textsubscript{补}= [-0]\textsubscript{反}+1=11111111+1
=1 00000000

对8位字长，进位被舍掉

\begin{figure}[htbp]
  \centering
  \includegraphics[width=\linewidth]{assets/figure_p17_1.png}
\end{figure}

\begin{figure}[htbp]
  \centering
  \includegraphics[width=\linewidth]{assets/figure_p17_2.png}
\end{figure}

\begin{figure}[htbp]
  \centering
  \includegraphics[width=\linewidth]{assets/figure_p17_3.png}
\end{figure}

\begin{figure}[htbp]
  \centering
  \includegraphics[width=\linewidth]{assets/figure_p17_4.png}
\end{figure}

\begin{figure}[htbp]
  \centering
  \includegraphics[width=\linewidth]{assets/figure_p17_5.png}
\end{figure}

\begin{figure}[htbp]
  \centering
  \includegraphics[width=\linewidth]{assets/figure_p17_6.png}
\end{figure}

\begin{figure}[htbp]
  \centering
  \includegraphics[width=\linewidth]{assets/figure_p17_7.png}
\end{figure}

\begin{figure}[htbp]
  \centering
  \includegraphics[width=\linewidth]{assets/figure_p17_8.png}
\end{figure}

\begin{figure}[htbp]
  \centering
  \includegraphics[width=\linewidth]{assets/figure_p17_9.png}
\end{figure}

\begin{figure}[htbp]
  \centering
  \includegraphics[width=\linewidth]{assets/figure_p17_10.png}
\end{figure}

\begin{figure}[htbp]
  \centering
  \includegraphics[width=\linewidth]{assets/figure_p17_11.png}
\end{figure}

\begin{figure}[htbp]
  \centering
  \includegraphics[width=\linewidth]{assets/figure_p17_12.png}
\end{figure}

\begin{figure}[htbp]
  \centering
  \includegraphics[width=\linewidth]{assets/figure_p17_13.png}
\end{figure}

\begin{figure}[htbp]
  \centering
  \includegraphics[width=\linewidth]{assets/figure_p17_14.png}
\end{figure}

\begin{figure}[htbp]
  \centering
  \includegraphics[width=\linewidth]{assets/figure_p17_15.png}
\end{figure}

\begin{figure}[htbp]
  \centering
  \includegraphics[width=\linewidth]{assets/figure_p17_16.png}
\end{figure}

\begin{figure}[htbp]
  \centering
  \includegraphics[width=\linewidth]{assets/figure_p17_17.png}
\end{figure}

\section{特殊数 10000000}

• 该数在原码中定义为： -0
• 在反码中定义为： -127
• 在补码中定义为： -128
• 对无符号数，（10000000）\textsubscript{B}=128

\begin{figure}[htbp]
  \centering
  \includegraphics[width=\linewidth]{assets/figure_p18_1.png}
\end{figure}

\begin{figure}[htbp]
  \centering
  \includegraphics[width=\linewidth]{assets/figure_p18_2.png}
\end{figure}

\begin{figure}[htbp]
  \centering
  \includegraphics[width=\linewidth]{assets/figure_p18_3.png}
\end{figure}

\begin{figure}[htbp]
  \centering
  \includegraphics[width=\linewidth]{assets/figure_p18_4.png}
\end{figure}

\begin{figure}[htbp]
  \centering
  \includegraphics[width=\linewidth]{assets/figure_p18_5.png}
\end{figure}

\begin{figure}[htbp]
  \centering
  \includegraphics[width=\linewidth]{assets/figure_p18_6.png}
\end{figure}

\begin{figure}[htbp]
  \centering
  \includegraphics[width=\linewidth]{assets/figure_p18_7.png}
\end{figure}

\begin{figure}[htbp]
  \centering
  \includegraphics[width=\linewidth]{assets/figure_p18_8.png}
\end{figure}

\begin{figure}[htbp]
  \centering
  \includegraphics[width=\linewidth]{assets/figure_p18_9.png}
\end{figure}

\begin{figure}[htbp]
  \centering
  \includegraphics[width=\linewidth]{assets/figure_p18_10.png}
\end{figure}

\begin{figure}[htbp]
  \centering
  \includegraphics[width=\linewidth]{assets/figure_p18_11.png}
\end{figure}

\begin{figure}[htbp]
  \centering
  \includegraphics[width=\linewidth]{assets/figure_p18_12.png}
\end{figure}

\begin{figure}[htbp]
  \centering
  \includegraphics[width=\linewidth]{assets/figure_p18_13.png}
\end{figure}

\begin{figure}[htbp]
  \centering
  \includegraphics[width=\linewidth]{assets/figure_p18_14.png}
\end{figure}

\begin{figure}[htbp]
  \centering
  \includegraphics[width=\linewidth]{assets/figure_p18_15.png}
\end{figure}

\begin{figure}[htbp]
  \centering
  \includegraphics[width=\linewidth]{assets/figure_p18_16.png}
\end{figure}

\begin{figure}[htbp]
  \centering
  \includegraphics[width=\linewidth]{assets/figure_p18_17.png}
\end{figure}

\section{负数三种编码之间的转换关系}

\begin{figure}[htbp]
  \centering
  \includegraphics[width=\linewidth]{assets/figure_p19_1.png}
\end{figure}

\begin{figure}[htbp]
  \centering
  \includegraphics[width=\linewidth]{assets/figure_p19_2.png}
\end{figure}

\begin{figure}[htbp]
  \centering
  \includegraphics[width=\linewidth]{assets/figure_p19_3.png}
\end{figure}

\begin{figure}[htbp]
  \centering
  \includegraphics[width=\linewidth]{assets/figure_p19_4.png}
\end{figure}

\begin{figure}[htbp]
  \centering
  \includegraphics[width=\linewidth]{assets/figure_p19_5.png}
\end{figure}

\begin{figure}[htbp]
  \centering
  \includegraphics[width=\linewidth]{assets/figure_p19_6.png}
\end{figure}

\begin{figure}[htbp]
  \centering
  \includegraphics[width=\linewidth]{assets/figure_p19_7.png}
\end{figure}

\begin{figure}[htbp]
  \centering
  \includegraphics[width=\linewidth]{assets/figure_p19_8.png}
\end{figure}

\begin{figure}[htbp]
  \centering
  \includegraphics[width=\linewidth]{assets/figure_p19_9.png}
\end{figure}

\begin{figure}[htbp]
  \centering
  \includegraphics[width=\linewidth]{assets/figure_p19_10.png}
\end{figure}

\begin{figure}[htbp]
  \centering
  \includegraphics[width=\linewidth]{assets/figure_p19_11.png}
\end{figure}

\begin{figure}[htbp]
  \centering
  \includegraphics[width=\linewidth]{assets/figure_p19_12.png}
\end{figure}

\begin{figure}[htbp]
  \centering
  \includegraphics[width=\linewidth]{assets/figure_p19_13.png}
\end{figure}

\begin{figure}[htbp]
  \centering
  \includegraphics[width=\linewidth]{assets/figure_p19_14.png}
\end{figure}

\begin{figure}[htbp]
  \centering
  \includegraphics[width=\linewidth]{assets/figure_p19_15.png}
\end{figure}

\begin{figure}[htbp]
  \centering
  \includegraphics[width=\linewidth]{assets/figure_p19_16.png}
\end{figure}

\begin{figure}[htbp]
  \centering
  \includegraphics[width=\linewidth]{assets/figure_p19_17.png}
\end{figure}

\begin{figure}[htbp]
  \centering
  \includegraphics[width=\linewidth]{assets/figure_p19_18.png}
\end{figure}

\section{无符号数的表示范围}

$0 \le X \le$ 2\textsuperscript{n}-1

若运算结果超出这个范围，则产生溢出。

无符号数的溢出判断准则：运算时，当最高位向更
高位有进位（或借位）时则产生溢出。

\begin{figure}[htbp]
  \centering
  \includegraphics[width=\linewidth]{assets/figure_p20_1.png}
\end{figure}

\begin{figure}[htbp]
  \centering
  \includegraphics[width=\linewidth]{assets/figure_p20_2.png}
\end{figure}

\begin{figure}[htbp]
  \centering
  \includegraphics[width=\linewidth]{assets/figure_p20_3.png}
\end{figure}

\begin{figure}[htbp]
  \centering
  \includegraphics[width=\linewidth]{assets/figure_p20_4.png}
\end{figure}

\begin{figure}[htbp]
  \centering
  \includegraphics[width=\linewidth]{assets/figure_p20_5.png}
\end{figure}

\begin{figure}[htbp]
  \centering
  \includegraphics[width=\linewidth]{assets/figure_p20_6.png}
\end{figure}

\begin{figure}[htbp]
  \centering
  \includegraphics[width=\linewidth]{assets/figure_p20_7.png}
\end{figure}

\begin{figure}[htbp]
  \centering
  \includegraphics[width=\linewidth]{assets/figure_p20_8.png}
\end{figure}

\begin{figure}[htbp]
  \centering
  \includegraphics[width=\linewidth]{assets/figure_p20_9.png}
\end{figure}

\begin{figure}[htbp]
  \centering
  \includegraphics[width=\linewidth]{assets/figure_p20_10.png}
\end{figure}

\begin{figure}[htbp]
  \centering
  \includegraphics[width=\linewidth]{assets/figure_p20_11.png}
\end{figure}

\begin{figure}[htbp]
  \centering
  \includegraphics[width=\linewidth]{assets/figure_p20_12.png}
\end{figure}

\begin{figure}[htbp]
  \centering
  \includegraphics[width=\linewidth]{assets/figure_p20_13.png}
\end{figure}

\begin{figure}[htbp]
  \centering
  \includegraphics[width=\linewidth]{assets/figure_p20_14.png}
\end{figure}

\begin{figure}[htbp]
  \centering
  \includegraphics[width=\linewidth]{assets/figure_p20_15.png}
\end{figure}

\begin{figure}[htbp]
  \centering
  \includegraphics[width=\linewidth]{assets/figure_p20_16.png}
\end{figure}

\begin{figure}[htbp]
  \centering
  \includegraphics[width=\linewidth]{assets/figure_p20_17.png}
\end{figure}

\section{无符号二进制数的运算}

算术的四种基本运算：加、减、乘、除

试计算011与010之和

结论：两个二进制数相加是通过逐位相加来实现的。

返回

\begin{figure}[htbp]
  \centering
  \includegraphics[width=\linewidth]{assets/figure_p21_1.png}
\end{figure}

\begin{figure}[htbp]
  \centering
  \includegraphics[width=\linewidth]{assets/figure_p21_2.png}
\end{figure}

\begin{figure}[htbp]
  \centering
  \includegraphics[width=\linewidth]{assets/figure_p21_3.png}
\end{figure}

\begin{figure}[htbp]
  \centering
  \includegraphics[width=\linewidth]{assets/figure_p21_4.png}
\end{figure}

\begin{figure}[htbp]
  \centering
  \includegraphics[width=\linewidth]{assets/figure_p21_5.png}
\end{figure}

\begin{figure}[htbp]
  \centering
  \includegraphics[width=\linewidth]{assets/figure_p21_6.png}
\end{figure}

\begin{figure}[htbp]
  \centering
  \includegraphics[width=\linewidth]{assets/figure_p21_7.png}
\end{figure}

\begin{figure}[htbp]
  \centering
  \includegraphics[width=\linewidth]{assets/figure_p21_8.png}
\end{figure}

\begin{figure}[htbp]
  \centering
  \includegraphics[width=\linewidth]{assets/figure_p21_9.png}
\end{figure}

\begin{figure}[htbp]
  \centering
  \includegraphics[width=\linewidth]{assets/figure_p21_10.png}
\end{figure}

\begin{figure}[htbp]
  \centering
  \includegraphics[width=\linewidth]{assets/figure_p21_11.png}
\end{figure}

\begin{figure}[htbp]
  \centering
  \includegraphics[width=\linewidth]{assets/figure_p21_12.png}
\end{figure}

\begin{figure}[htbp]
  \centering
  \includegraphics[width=\linewidth]{assets/figure_p21_13.png}
\end{figure}

\begin{figure}[htbp]
  \centering
  \includegraphics[width=\linewidth]{assets/figure_p21_14.png}
\end{figure}

\begin{figure}[htbp]
  \centering
  \includegraphics[width=\linewidth]{assets/figure_p21_15.png}
\end{figure}

\begin{figure}[htbp]
  \centering
  \includegraphics[width=\linewidth]{assets/figure_p21_16.png}
\end{figure}

\begin{figure}[htbp]
  \centering
  \includegraphics[width=\linewidth]{assets/figure_p21_17.png}
\end{figure}

\begin{figure}[htbp]
  \centering
  \includegraphics[width=\linewidth]{assets/figure_p21_18.png}
\end{figure}

\section{加法运算}

推广:

设两个二进制数分别为\textbf{A=A}\textsubscript{\textbf{3}}\textbf{A}\textsubscript{\textbf{2}}\textbf{A}\textsubscript{\textbf{1}}\textbf{A}\textsubscript{\textbf{0}}\textsubscript{，}\textbf{B=B}\textsubscript{\textbf{3}}\textbf{B}\textsubscript{\textbf{2}}\textbf{B}\textsubscript{\textbf{1}}\textbf{B}\textsubscript{\textbf{0}}两
数之和为\textbf{S=S}\textsubscript{\textbf{3}}\textbf{S}\textsubscript{\textbf{2}}\textbf{S}\textsubscript{\textbf{1}}\textbf{S}\textsubscript{\textbf{0}}

\textbf{S}\textsubscript{\textbf{0}}\textbf{=A}\textsubscript{\textbf{0}}\textbf{+B}\textsubscript{\textbf{0}} 进位\textbf{C}\textsubscript{\textbf{1}}

\[ S 1 =A 1 +B 1 +C 1 进位 C 2 \]

\[ S 2 =A 2 +B 2 +C 2 进位 C 3 \]

\[ S 3 =A 3 +B 3 +C 3 进位 C 4 \]

\textbf{A+B=C}\textsubscript{\textbf{4}}\textbf{S}\textsubscript{\textbf{3}}\textbf{S}\textsubscript{\textbf{2}}\textbf{S}\textsubscript{\textbf{1}}\textbf{S}\textsubscript{\textbf{0}}
返回

\begin{figure}[htbp]
  \centering
  \includegraphics[width=\linewidth]{assets/figure_p22_1.png}
\end{figure}

\begin{figure}[htbp]
  \centering
  \includegraphics[width=\linewidth]{assets/figure_p22_2.png}
\end{figure}

\begin{figure}[htbp]
  \centering
  \includegraphics[width=\linewidth]{assets/figure_p22_3.png}
\end{figure}

\begin{figure}[htbp]
  \centering
  \includegraphics[width=\linewidth]{assets/figure_p22_4.png}
\end{figure}

\begin{figure}[htbp]
  \centering
  \includegraphics[width=\linewidth]{assets/figure_p22_5.png}
\end{figure}

\begin{figure}[htbp]
  \centering
  \includegraphics[width=\linewidth]{assets/figure_p22_6.png}
\end{figure}

\begin{figure}[htbp]
  \centering
  \includegraphics[width=\linewidth]{assets/figure_p22_7.png}
\end{figure}

\begin{figure}[htbp]
  \centering
  \includegraphics[width=\linewidth]{assets/figure_p22_8.png}
\end{figure}

\begin{figure}[htbp]
  \centering
  \includegraphics[width=\linewidth]{assets/figure_p22_9.png}
\end{figure}

\begin{figure}[htbp]
  \centering
  \includegraphics[width=\linewidth]{assets/figure_p22_10.png}
\end{figure}

\begin{figure}[htbp]
  \centering
  \includegraphics[width=\linewidth]{assets/figure_p22_11.png}
\end{figure}

\begin{figure}[htbp]
  \centering
  \includegraphics[width=\linewidth]{assets/figure_p22_12.png}
\end{figure}

\begin{figure}[htbp]
  \centering
  \includegraphics[width=\linewidth]{assets/figure_p22_13.png}
\end{figure}

\begin{figure}[htbp]
  \centering
  \includegraphics[width=\linewidth]{assets/figure_p22_14.png}
\end{figure}

\begin{figure}[htbp]
  \centering
  \includegraphics[width=\linewidth]{assets/figure_p22_15.png}
\end{figure}

\begin{figure}[htbp]
  \centering
  \includegraphics[width=\linewidth]{assets/figure_p22_16.png}
\end{figure}

\begin{figure}[htbp]
  \centering
  \includegraphics[width=\linewidth]{assets/figure_p22_17.png}
\end{figure}

\section{减法运算}

原理：将减数Ｂ变成补码后，再与被减数Ａ相加，其和
（如有进位的话，则舍去进位）就是两个数之差。

什么是补码？

补码＝反码（原码取反）＋１

Ｙ＝１０００－０１００
＝１０００＋（１０１１＋１）
＝１０００＋１１００
＝１０１００ １是进位，舍去

\begin{figure}[htbp]
  \centering
  \includegraphics[width=\linewidth]{assets/figure_p23_1.png}
\end{figure}

\begin{figure}[htbp]
  \centering
  \includegraphics[width=\linewidth]{assets/figure_p23_2.png}
\end{figure}

\begin{figure}[htbp]
  \centering
  \includegraphics[width=\linewidth]{assets/figure_p23_3.png}
\end{figure}

\begin{figure}[htbp]
  \centering
  \includegraphics[width=\linewidth]{assets/figure_p23_4.png}
\end{figure}

\begin{figure}[htbp]
  \centering
  \includegraphics[width=\linewidth]{assets/figure_p23_5.png}
\end{figure}

\begin{figure}[htbp]
  \centering
  \includegraphics[width=\linewidth]{assets/figure_p23_6.png}
\end{figure}

\begin{figure}[htbp]
  \centering
  \includegraphics[width=\linewidth]{assets/figure_p23_7.png}
\end{figure}

\begin{figure}[htbp]
  \centering
  \includegraphics[width=\linewidth]{assets/figure_p23_8.png}
\end{figure}

\begin{figure}[htbp]
  \centering
  \includegraphics[width=\linewidth]{assets/figure_p23_9.png}
\end{figure}

\begin{figure}[htbp]
  \centering
  \includegraphics[width=\linewidth]{assets/figure_p23_10.png}
\end{figure}

\begin{figure}[htbp]
  \centering
  \includegraphics[width=\linewidth]{assets/figure_p23_11.png}
\end{figure}

\begin{figure}[htbp]
  \centering
  \includegraphics[width=\linewidth]{assets/figure_p23_12.png}
\end{figure}

\begin{figure}[htbp]
  \centering
  \includegraphics[width=\linewidth]{assets/figure_p23_13.png}
\end{figure}

\begin{figure}[htbp]
  \centering
  \includegraphics[width=\linewidth]{assets/figure_p23_14.png}
\end{figure}

\begin{figure}[htbp]
  \centering
  \includegraphics[width=\linewidth]{assets/figure_p23_15.png}
\end{figure}

\begin{figure}[htbp]
  \centering
  \includegraphics[width=\linewidth]{assets/figure_p23_16.png}
\end{figure}

\begin{figure}[htbp]
  \centering
  \includegraphics[width=\linewidth]{assets/figure_p23_17.png}
\end{figure}

\section{乘法运算}

• 对二进制数，乘以\textbf{2}相当于左移一位

$\textbf{00001011}\times \textbf{010=0010110B}$

方法：1 按照十进制的乘法过程
2 采用移位加的方法

\begin{figure}[htbp]
  \centering
  \includegraphics[width=\linewidth]{assets/figure_p24_1.png}
\end{figure}

\begin{figure}[htbp]
  \centering
  \includegraphics[width=\linewidth]{assets/figure_p24_2.png}
\end{figure}

\begin{figure}[htbp]
  \centering
  \includegraphics[width=\linewidth]{assets/figure_p24_3.png}
\end{figure}

\begin{figure}[htbp]
  \centering
  \includegraphics[width=\linewidth]{assets/figure_p24_4.png}
\end{figure}

\begin{figure}[htbp]
  \centering
  \includegraphics[width=\linewidth]{assets/figure_p24_5.png}
\end{figure}

\begin{figure}[htbp]
  \centering
  \includegraphics[width=\linewidth]{assets/figure_p24_6.png}
\end{figure}

\begin{figure}[htbp]
  \centering
  \includegraphics[width=\linewidth]{assets/figure_p24_7.png}
\end{figure}

\begin{figure}[htbp]
  \centering
  \includegraphics[width=\linewidth]{assets/figure_p24_8.png}
\end{figure}

\begin{figure}[htbp]
  \centering
  \includegraphics[width=\linewidth]{assets/figure_p24_9.png}
\end{figure}

\begin{figure}[htbp]
  \centering
  \includegraphics[width=\linewidth]{assets/figure_p24_10.png}
\end{figure}

\begin{figure}[htbp]
  \centering
  \includegraphics[width=\linewidth]{assets/figure_p24_11.png}
\end{figure}

\begin{figure}[htbp]
  \centering
  \includegraphics[width=\linewidth]{assets/figure_p24_12.png}
\end{figure}

\begin{figure}[htbp]
  \centering
  \includegraphics[width=\linewidth]{assets/figure_p24_13.png}
\end{figure}

\begin{figure}[htbp]
  \centering
  \includegraphics[width=\linewidth]{assets/figure_p24_14.png}
\end{figure}

\begin{figure}[htbp]
  \centering
  \includegraphics[width=\linewidth]{assets/figure_p24_15.png}
\end{figure}

\begin{figure}[htbp]
  \centering
  \includegraphics[width=\linewidth]{assets/figure_p24_16.png}
\end{figure}

\begin{figure}[htbp]
  \centering
  \includegraphics[width=\linewidth]{assets/figure_p24_17.png}
\end{figure}

\begin{figure}[htbp]
  \centering
  \includegraphics[width=\linewidth]{assets/figure_p24_18.png}
\end{figure}

\begin{figure}[htbp]
  \centering
  \includegraphics[width=\linewidth]{assets/figure_p24_19.png}
\end{figure}

\begin{figure}[htbp]
  \centering
  \includegraphics[width=\linewidth]{assets/figure_p24_20.png}
\end{figure}

\section{除法运算}

• 对二进制数，除以\textbf{2}则相当于右移\textbf{1}位

$\textbf{00001011}\div \textbf{010=00000010B}$
即：商\textbf{=00000010B}
余数\textbf{=11B}

\begin{figure}[htbp]
  \centering
  \includegraphics[width=\linewidth]{assets/figure_p25_1.png}
\end{figure}

\begin{figure}[htbp]
  \centering
  \includegraphics[width=\linewidth]{assets/figure_p25_2.png}
\end{figure}

\begin{figure}[htbp]
  \centering
  \includegraphics[width=\linewidth]{assets/figure_p25_3.png}
\end{figure}

\begin{figure}[htbp]
  \centering
  \includegraphics[width=\linewidth]{assets/figure_p25_4.png}
\end{figure}

\begin{figure}[htbp]
  \centering
  \includegraphics[width=\linewidth]{assets/figure_p25_5.png}
\end{figure}

\begin{figure}[htbp]
  \centering
  \includegraphics[width=\linewidth]{assets/figure_p25_6.png}
\end{figure}

\begin{figure}[htbp]
  \centering
  \includegraphics[width=\linewidth]{assets/figure_p25_7.png}
\end{figure}

\begin{figure}[htbp]
  \centering
  \includegraphics[width=\linewidth]{assets/figure_p25_8.png}
\end{figure}

\begin{figure}[htbp]
  \centering
  \includegraphics[width=\linewidth]{assets/figure_p25_9.png}
\end{figure}

\begin{figure}[htbp]
  \centering
  \includegraphics[width=\linewidth]{assets/figure_p25_10.png}
\end{figure}

\begin{figure}[htbp]
  \centering
  \includegraphics[width=\linewidth]{assets/figure_p25_11.png}
\end{figure}

\begin{figure}[htbp]
  \centering
  \includegraphics[width=\linewidth]{assets/figure_p25_12.png}
\end{figure}

\begin{figure}[htbp]
  \centering
  \includegraphics[width=\linewidth]{assets/figure_p25_13.png}
\end{figure}

\begin{figure}[htbp]
  \centering
  \includegraphics[width=\linewidth]{assets/figure_p25_14.png}
\end{figure}

\begin{figure}[htbp]
  \centering
  \includegraphics[width=\linewidth]{assets/figure_p25_15.png}
\end{figure}

\begin{figure}[htbp]
  \centering
  \includegraphics[width=\linewidth]{assets/figure_p25_16.png}
\end{figure}

\begin{figure}[htbp]
  \centering
  \includegraphics[width=\linewidth]{assets/figure_p25_17.png}
\end{figure}

\begin{figure}[htbp]
  \centering
  \includegraphics[width=\linewidth]{assets/figure_p25_18.png}
\end{figure}

\begin{figure}[htbp]
  \centering
  \includegraphics[width=\linewidth]{assets/figure_p25_19.png}
\end{figure}

\begin{figure}[htbp]
  \centering
  \includegraphics[width=\linewidth]{assets/figure_p25_20.png}
\end{figure}

\section{带符号数的表示范围：}

对\textbf{8}位二进制数：
• 原码： \textbf{-127 \textasciitilde{} +127}
• 反码： \textbf{-127 \textasciitilde{} +127}
• 补码： \textbf{-128 \textasciitilde{} +127}

对用补码表示的二进制数转换成十进制：

\textbf{1}）求出真值
\textbf{2}）进行二\textbf{——}十转换

\begin{figure}[htbp]
  \centering
  \includegraphics[width=\linewidth]{assets/figure_p26_1.png}
\end{figure}

\begin{figure}[htbp]
  \centering
  \includegraphics[width=\linewidth]{assets/figure_p26_2.png}
\end{figure}

\begin{figure}[htbp]
  \centering
  \includegraphics[width=\linewidth]{assets/figure_p26_3.png}
\end{figure}

\begin{figure}[htbp]
  \centering
  \includegraphics[width=\linewidth]{assets/figure_p26_4.png}
\end{figure}

\begin{figure}[htbp]
  \centering
  \includegraphics[width=\linewidth]{assets/figure_p26_5.png}
\end{figure}

\begin{figure}[htbp]
  \centering
  \includegraphics[width=\linewidth]{assets/figure_p26_6.png}
\end{figure}

\begin{figure}[htbp]
  \centering
  \includegraphics[width=\linewidth]{assets/figure_p26_7.png}
\end{figure}

\begin{figure}[htbp]
  \centering
  \includegraphics[width=\linewidth]{assets/figure_p26_8.png}
\end{figure}

\begin{figure}[htbp]
  \centering
  \includegraphics[width=\linewidth]{assets/figure_p26_9.png}
\end{figure}

\begin{figure}[htbp]
  \centering
  \includegraphics[width=\linewidth]{assets/figure_p26_10.png}
\end{figure}

\begin{figure}[htbp]
  \centering
  \includegraphics[width=\linewidth]{assets/figure_p26_11.png}
\end{figure}

\begin{figure}[htbp]
  \centering
  \includegraphics[width=\linewidth]{assets/figure_p26_12.png}
\end{figure}

\begin{figure}[htbp]
  \centering
  \includegraphics[width=\linewidth]{assets/figure_p26_13.png}
\end{figure}

\begin{figure}[htbp]
  \centering
  \includegraphics[width=\linewidth]{assets/figure_p26_14.png}
\end{figure}

\begin{figure}[htbp]
  \centering
  \includegraphics[width=\linewidth]{assets/figure_p26_15.png}
\end{figure}

\begin{figure}[htbp]
  \centering
  \includegraphics[width=\linewidth]{assets/figure_p26_16.png}
\end{figure}

\begin{figure}[htbp]
  \centering
  \includegraphics[width=\linewidth]{assets/figure_p26_17.png}
\end{figure}

\begin{figure}[htbp]
  \centering
  \includegraphics[width=\linewidth]{assets/figure_p26_18.png}
\end{figure}

\section{符号数的算术运算}

• 通过引进补码，可将减法运算转换为加法运算

即：\textbf{[X+Y]}\textsubscript{补}\textbf{=[X]}\textsubscript{补}\textbf{+[Y]}\textsubscript{补}
\textbf{[X-Y]}\textsubscript{补}\textbf{=[X+(-Y)]}\textsubscript{补}
\textbf{=[X]}\textsubscript{补}\textbf{+[-Y]}\textsubscript{补}

\begin{figure}[htbp]
  \centering
  \includegraphics[width=\linewidth]{assets/figure_p27_1.png}
\end{figure}

\begin{figure}[htbp]
  \centering
  \includegraphics[width=\linewidth]{assets/figure_p27_2.png}
\end{figure}

\begin{figure}[htbp]
  \centering
  \includegraphics[width=\linewidth]{assets/figure_p27_3.png}
\end{figure}

\begin{figure}[htbp]
  \centering
  \includegraphics[width=\linewidth]{assets/figure_p27_4.png}
\end{figure}

\begin{figure}[htbp]
  \centering
  \includegraphics[width=\linewidth]{assets/figure_p27_5.png}
\end{figure}

\begin{figure}[htbp]
  \centering
  \includegraphics[width=\linewidth]{assets/figure_p27_6.png}
\end{figure}

\begin{figure}[htbp]
  \centering
  \includegraphics[width=\linewidth]{assets/figure_p27_7.png}
\end{figure}

\begin{figure}[htbp]
  \centering
  \includegraphics[width=\linewidth]{assets/figure_p27_8.png}
\end{figure}

\begin{figure}[htbp]
  \centering
  \includegraphics[width=\linewidth]{assets/figure_p27_9.png}
\end{figure}

\begin{figure}[htbp]
  \centering
  \includegraphics[width=\linewidth]{assets/figure_p27_10.png}
\end{figure}

\begin{figure}[htbp]
  \centering
  \includegraphics[width=\linewidth]{assets/figure_p27_11.png}
\end{figure}

\begin{figure}[htbp]
  \centering
  \includegraphics[width=\linewidth]{assets/figure_p27_12.png}
\end{figure}

\begin{figure}[htbp]
  \centering
  \includegraphics[width=\linewidth]{assets/figure_p27_13.png}
\end{figure}

\begin{figure}[htbp]
  \centering
  \includegraphics[width=\linewidth]{assets/figure_p27_14.png}
\end{figure}

\begin{figure}[htbp]
  \centering
  \includegraphics[width=\linewidth]{assets/figure_p27_15.png}
\end{figure}

\begin{figure}[htbp]
  \centering
  \includegraphics[width=\linewidth]{assets/figure_p27_16.png}
\end{figure}

\begin{figure}[htbp]
  \centering
  \includegraphics[width=\linewidth]{assets/figure_p27_17.png}
\end{figure}

\section{[ 例 ] ：}

\[ X=-0110100 ， Y=+1110100 ，求 X+Y= ？ \]

• \textbf{[X] =10110100}\textsuperscript{无论正负，真值不变}

原

\[ • [X] 补 = [X] 反 +1=11001100 \]

• \textbf{[Y]}\textsubscript{补}\textbf{= [Y]}\textsubscript{原}\textbf{=01110100}

\[ • 所以： [X+Y] 补 = [X] 补 + [Y] 补 \]

\textbf{=11001100+01110100}

\textbf{=01000000}

\begin{figure}[htbp]
  \centering
  \includegraphics[width=\linewidth]{assets/figure_p28_1.png}
\end{figure}

\begin{figure}[htbp]
  \centering
  \includegraphics[width=\linewidth]{assets/figure_p28_2.png}
\end{figure}

\begin{figure}[htbp]
  \centering
  \includegraphics[width=\linewidth]{assets/figure_p28_3.png}
\end{figure}

\begin{figure}[htbp]
  \centering
  \includegraphics[width=\linewidth]{assets/figure_p28_4.png}
\end{figure}

\begin{figure}[htbp]
  \centering
  \includegraphics[width=\linewidth]{assets/figure_p28_5.png}
\end{figure}

\begin{figure}[htbp]
  \centering
  \includegraphics[width=\linewidth]{assets/figure_p28_6.png}
\end{figure}

\begin{figure}[htbp]
  \centering
  \includegraphics[width=\linewidth]{assets/figure_p28_7.png}
\end{figure}

\begin{figure}[htbp]
  \centering
  \includegraphics[width=\linewidth]{assets/figure_p28_8.png}
\end{figure}

\begin{figure}[htbp]
  \centering
  \includegraphics[width=\linewidth]{assets/figure_p28_9.png}
\end{figure}

\begin{figure}[htbp]
  \centering
  \includegraphics[width=\linewidth]{assets/figure_p28_10.png}
\end{figure}

\begin{figure}[htbp]
  \centering
  \includegraphics[width=\linewidth]{assets/figure_p28_11.png}
\end{figure}

\begin{figure}[htbp]
  \centering
  \includegraphics[width=\linewidth]{assets/figure_p28_12.png}
\end{figure}

\begin{figure}[htbp]
  \centering
  \includegraphics[width=\linewidth]{assets/figure_p28_13.png}
\end{figure}

\begin{figure}[htbp]
  \centering
  \includegraphics[width=\linewidth]{assets/figure_p28_14.png}
\end{figure}

\begin{figure}[htbp]
  \centering
  \includegraphics[width=\linewidth]{assets/figure_p28_15.png}
\end{figure}

\begin{figure}[htbp]
  \centering
  \includegraphics[width=\linewidth]{assets/figure_p28_16.png}
\end{figure}

\begin{figure}[htbp]
  \centering
  \includegraphics[width=\linewidth]{assets/figure_p28_17.png}
\end{figure}

两个补码表示的二进制数相加时的符号位讨论
例：用二进制补码运算求出
13＋10

解：  13 0 01101  13 0 01101
 10 0 01010  10 1 10110
 23 0 10111  3 0 00011

 13 1 10011  13 1 10011
 10 0 01010  10 1 10110
 3 1 11101  23 1 01001

\begin{figure}[htbp]
  \centering
  \includegraphics[width=\linewidth]{assets/figure_p29_1.png}
\end{figure}

\begin{figure}[htbp]
  \centering
  \includegraphics[width=\linewidth]{assets/figure_p29_2.png}
\end{figure}

\begin{figure}[htbp]
  \centering
  \includegraphics[width=\linewidth]{assets/figure_p29_3.png}
\end{figure}

\begin{figure}[htbp]
  \centering
  \includegraphics[width=\linewidth]{assets/figure_p29_4.png}
\end{figure}

\begin{figure}[htbp]
  \centering
  \includegraphics[width=\linewidth]{assets/figure_p29_5.png}
\end{figure}

\begin{figure}[htbp]
  \centering
  \includegraphics[width=\linewidth]{assets/figure_p29_6.png}
\end{figure}

\begin{figure}[htbp]
  \centering
  \includegraphics[width=\linewidth]{assets/figure_p29_7.png}
\end{figure}

\begin{figure}[htbp]
  \centering
  \includegraphics[width=\linewidth]{assets/figure_p29_8.png}
\end{figure}

\begin{figure}[htbp]
  \centering
  \includegraphics[width=\linewidth]{assets/figure_p29_9.png}
\end{figure}

\begin{figure}[htbp]
  \centering
  \includegraphics[width=\linewidth]{assets/figure_p29_10.png}
\end{figure}

\begin{figure}[htbp]
  \centering
  \includegraphics[width=\linewidth]{assets/figure_p29_11.png}
\end{figure}

\begin{figure}[htbp]
  \centering
  \includegraphics[width=\linewidth]{assets/figure_p29_12.png}
\end{figure}

\begin{figure}[htbp]
  \centering
  \includegraphics[width=\linewidth]{assets/figure_p29_13.png}
\end{figure}

\begin{figure}[htbp]
  \centering
  \includegraphics[width=\linewidth]{assets/figure_p29_14.png}
\end{figure}

\begin{figure}[htbp]
  \centering
  \includegraphics[width=\linewidth]{assets/figure_p29_15.png}
\end{figure}

\begin{figure}[htbp]
  \centering
  \includegraphics[width=\linewidth]{assets/figure_p29_16.png}
\end{figure}

\begin{figure}[htbp]
  \centering
  \includegraphics[width=\linewidth]{assets/figure_p29_17.png}
\end{figure}

\begin{figure}[htbp]
  \centering
  \includegraphics[width=\linewidth]{assets/figure_p29_18.png}
\end{figure}

\begin{figure}[htbp]
  \centering
  \includegraphics[width=\linewidth]{assets/figure_p29_19.png}
\end{figure}

\begin{figure}[htbp]
  \centering
  \includegraphics[width=\linewidth]{assets/figure_p29_20.png}
\end{figure}

\begin{figure}[htbp]
  \centering
  \includegraphics[width=\linewidth]{assets/figure_p29_21.png}
\end{figure}

\begin{figure}[htbp]
  \centering
  \includegraphics[width=\linewidth]{assets/figure_p29_22.png}
\end{figure}

\begin{figure}[htbp]
  \centering
  \includegraphics[width=\linewidth]{assets/figure_p29_23.png}
\end{figure}

\begin{figure}[htbp]
  \centering
  \includegraphics[width=\linewidth]{assets/figure_p29_24.png}
\end{figure}

\begin{figure}[htbp]
  \centering
  \includegraphics[width=\linewidth]{assets/figure_p29_25.png}
\end{figure}

\begin{figure}[htbp]
  \centering
  \includegraphics[width=\linewidth]{assets/figure_p29_26.png}
\end{figure}

\begin{figure}[htbp]
  \centering
  \includegraphics[width=\linewidth]{assets/figure_p29_27.png}
\end{figure}

8.1.2 几种简单的编码
1. 二 - 十进制码 (BCD码)\textbf{( Binary Coded Decimal codes)}

用四位二进制代码来表示一位十进制数码,这样的代
码称为二-十进制码,或BCD码.

四位二进制有16种不同的组合,可以在这16种代码中

任选10种表示十进制数的10个不同符号,选择方法很多.选
择方法不同,就能得到不同的编码形式.

常见的\textbf{BCD}码有8421码、5421码、2421码、余3码等。

\begin{figure}[htbp]
  \centering
  \includegraphics[width=\linewidth]{assets/figure_p30_1.png}
\end{figure}

\begin{figure}[htbp]
  \centering
  \includegraphics[width=\linewidth]{assets/figure_p30_2.png}
\end{figure}

\begin{figure}[htbp]
  \centering
  \includegraphics[width=\linewidth]{assets/figure_p30_3.png}
\end{figure}

\begin{figure}[htbp]
  \centering
  \includegraphics[width=\linewidth]{assets/figure_p30_4.png}
\end{figure}

\begin{figure}[htbp]
  \centering
  \includegraphics[width=\linewidth]{assets/figure_p30_5.png}
\end{figure}

\begin{figure}[htbp]
  \centering
  \includegraphics[width=\linewidth]{assets/figure_p30_6.png}
\end{figure}

\begin{figure}[htbp]
  \centering
  \includegraphics[width=\linewidth]{assets/figure_p30_7.png}
\end{figure}

\begin{figure}[htbp]
  \centering
  \includegraphics[width=\linewidth]{assets/figure_p30_8.png}
\end{figure}

\begin{figure}[htbp]
  \centering
  \includegraphics[width=\linewidth]{assets/figure_p30_9.png}
\end{figure}

\begin{figure}[htbp]
  \centering
  \includegraphics[width=\linewidth]{assets/figure_p30_10.png}
\end{figure}

\begin{figure}[htbp]
  \centering
  \includegraphics[width=\linewidth]{assets/figure_p30_11.png}
\end{figure}

\begin{figure}[htbp]
  \centering
  \includegraphics[width=\linewidth]{assets/figure_p30_12.png}
\end{figure}

\begin{figure}[htbp]
  \centering
  \includegraphics[width=\linewidth]{assets/figure_p30_13.png}
\end{figure}

\begin{figure}[htbp]
  \centering
  \includegraphics[width=\linewidth]{assets/figure_p30_14.png}
\end{figure}

\begin{figure}[htbp]
  \centering
  \includegraphics[width=\linewidth]{assets/figure_p30_15.png}
\end{figure}

\begin{figure}[htbp]
  \centering
  \includegraphics[width=\linewidth]{assets/figure_p30_16.png}
\end{figure}

\begin{figure}[htbp]
  \centering
  \includegraphics[width=\linewidth]{assets/figure_p30_17.png}
\end{figure}

常用BCD码

十进制数 \textbf{8421}码 \textbf{5421}码 \textbf{2421}码 余\textbf{3}码
\textbf{0 0000 0000 0000 0011}
\textbf{1 0001 0001 0001 0100}
\textbf{2 0010 0010 0010 0101}
\textbf{3 0011 0011 0011 0110}
\textbf{4 0100 0100 0100 0111}
\textbf{5 0101 1000 1011 1000}
\textbf{6 0110 1001 1100 1001}
\textbf{7 0111 1010 1101 1010}
\textbf{8 1000 1011 1110 1011}
\textbf{9 1001 1100 1111 1100}

\begin{figure}[htbp]
  \centering
  \includegraphics[width=\linewidth]{assets/figure_p31_1.png}
\end{figure}

\begin{figure}[htbp]
  \centering
  \includegraphics[width=\linewidth]{assets/figure_p31_2.png}
\end{figure}

\begin{figure}[htbp]
  \centering
  \includegraphics[width=\linewidth]{assets/figure_p31_3.png}
\end{figure}

\begin{figure}[htbp]
  \centering
  \includegraphics[width=\linewidth]{assets/figure_p31_4.png}
\end{figure}

\begin{figure}[htbp]
  \centering
  \includegraphics[width=\linewidth]{assets/figure_p31_5.png}
\end{figure}

\begin{figure}[htbp]
  \centering
  \includegraphics[width=\linewidth]{assets/figure_p31_6.png}
\end{figure}

\begin{figure}[htbp]
  \centering
  \includegraphics[width=\linewidth]{assets/figure_p31_7.png}
\end{figure}

\begin{figure}[htbp]
  \centering
  \includegraphics[width=\linewidth]{assets/figure_p31_8.png}
\end{figure}

\begin{figure}[htbp]
  \centering
  \includegraphics[width=\linewidth]{assets/figure_p31_9.png}
\end{figure}

\begin{figure}[htbp]
  \centering
  \includegraphics[width=\linewidth]{assets/figure_p31_10.png}
\end{figure}

\begin{figure}[htbp]
  \centering
  \includegraphics[width=\linewidth]{assets/figure_p31_11.png}
\end{figure}

\begin{figure}[htbp]
  \centering
  \includegraphics[width=\linewidth]{assets/figure_p31_12.png}
\end{figure}

\begin{figure}[htbp]
  \centering
  \includegraphics[width=\linewidth]{assets/figure_p31_13.png}
\end{figure}

\begin{figure}[htbp]
  \centering
  \includegraphics[width=\linewidth]{assets/figure_p31_14.png}
\end{figure}

\begin{figure}[htbp]
  \centering
  \includegraphics[width=\linewidth]{assets/figure_p31_15.png}
\end{figure}

\begin{figure}[htbp]
  \centering
  \includegraphics[width=\linewidth]{assets/figure_p31_16.png}
\end{figure}

\begin{figure}[htbp]
  \centering
  \includegraphics[width=\linewidth]{assets/figure_p31_17.png}
\end{figure}

(1) 有权\textbf{BCD}码：每位数码都有确定的位权的码，

例如：8421码、5421码、2421码.

如: 5421码1011代表5+0+2+1=8;
2421码1100代表2+4+0+0=6.
* 5421\textbf{BCD}码和2421\textbf{BCD}码不唯一.
例: 2421\textbf{BCD}码0110也可表示6
* 在表中：
① 8421\textbf{BCD}码和代表\textbf{0\textasciitilde{}9}的二进制数一一对应；

\begin{figure}[htbp]
  \centering
  \includegraphics[width=\linewidth]{assets/figure_p32_1.png}
\end{figure}

\begin{figure}[htbp]
  \centering
  \includegraphics[width=\linewidth]{assets/figure_p32_2.png}
\end{figure}

\begin{figure}[htbp]
  \centering
  \includegraphics[width=\linewidth]{assets/figure_p32_3.png}
\end{figure}

\begin{figure}[htbp]
  \centering
  \includegraphics[width=\linewidth]{assets/figure_p32_4.png}
\end{figure}

\begin{figure}[htbp]
  \centering
  \includegraphics[width=\linewidth]{assets/figure_p32_5.png}
\end{figure}

\begin{figure}[htbp]
  \centering
  \includegraphics[width=\linewidth]{assets/figure_p32_6.png}
\end{figure}

\begin{figure}[htbp]
  \centering
  \includegraphics[width=\linewidth]{assets/figure_p32_7.png}
\end{figure}

\begin{figure}[htbp]
  \centering
  \includegraphics[width=\linewidth]{assets/figure_p32_8.png}
\end{figure}

\begin{figure}[htbp]
  \centering
  \includegraphics[width=\linewidth]{assets/figure_p32_9.png}
\end{figure}

\begin{figure}[htbp]
  \centering
  \includegraphics[width=\linewidth]{assets/figure_p32_10.png}
\end{figure}

\begin{figure}[htbp]
  \centering
  \includegraphics[width=\linewidth]{assets/figure_p32_11.png}
\end{figure}

\begin{figure}[htbp]
  \centering
  \includegraphics[width=\linewidth]{assets/figure_p32_12.png}
\end{figure}

\begin{figure}[htbp]
  \centering
  \includegraphics[width=\linewidth]{assets/figure_p32_13.png}
\end{figure}

\begin{figure}[htbp]
  \centering
  \includegraphics[width=\linewidth]{assets/figure_p32_14.png}
\end{figure}

\begin{figure}[htbp]
  \centering
  \includegraphics[width=\linewidth]{assets/figure_p32_15.png}
\end{figure}

\begin{figure}[htbp]
  \centering
  \includegraphics[width=\linewidth]{assets/figure_p32_16.png}
\end{figure}

\begin{figure}[htbp]
  \centering
  \includegraphics[width=\linewidth]{assets/figure_p32_17.png}
\end{figure}

② 5421\textbf{BCD}码的前5个码和8421\textbf{BCD}码相同，后5个码在

前5个码的基础上加1000构成，这样的码，前5个码和后5

个码一一对应相同，仅高位不同；

③ 2421\textbf{BCD}码的前5个码和8421\textbf{BCD}码相同，后5个码以

中心对称取反,这样的码称为自反代码.

例： $\textbf{4\to 0100 5\to 1011}$
$\textbf{0\to 0000 9\to 1111}$

\begin{figure}[htbp]
  \centering
  \includegraphics[width=\linewidth]{assets/figure_p33_1.png}
\end{figure}

\begin{figure}[htbp]
  \centering
  \includegraphics[width=\linewidth]{assets/figure_p33_2.png}
\end{figure}

\begin{figure}[htbp]
  \centering
  \includegraphics[width=\linewidth]{assets/figure_p33_3.png}
\end{figure}

\begin{figure}[htbp]
  \centering
  \includegraphics[width=\linewidth]{assets/figure_p33_4.png}
\end{figure}

\begin{figure}[htbp]
  \centering
  \includegraphics[width=\linewidth]{assets/figure_p33_5.png}
\end{figure}

\begin{figure}[htbp]
  \centering
  \includegraphics[width=\linewidth]{assets/figure_p33_6.png}
\end{figure}

\begin{figure}[htbp]
  \centering
  \includegraphics[width=\linewidth]{assets/figure_p33_7.png}
\end{figure}

\begin{figure}[htbp]
  \centering
  \includegraphics[width=\linewidth]{assets/figure_p33_8.png}
\end{figure}

\begin{figure}[htbp]
  \centering
  \includegraphics[width=\linewidth]{assets/figure_p33_9.png}
\end{figure}

\begin{figure}[htbp]
  \centering
  \includegraphics[width=\linewidth]{assets/figure_p33_10.png}
\end{figure}

\begin{figure}[htbp]
  \centering
  \includegraphics[width=\linewidth]{assets/figure_p33_11.png}
\end{figure}

\begin{figure}[htbp]
  \centering
  \includegraphics[width=\linewidth]{assets/figure_p33_12.png}
\end{figure}

\begin{figure}[htbp]
  \centering
  \includegraphics[width=\linewidth]{assets/figure_p33_13.png}
\end{figure}

\begin{figure}[htbp]
  \centering
  \includegraphics[width=\linewidth]{assets/figure_p33_14.png}
\end{figure}

\begin{figure}[htbp]
  \centering
  \includegraphics[width=\linewidth]{assets/figure_p33_15.png}
\end{figure}

\begin{figure}[htbp]
  \centering
  \includegraphics[width=\linewidth]{assets/figure_p33_16.png}
\end{figure}

\begin{figure}[htbp]
  \centering
  \includegraphics[width=\linewidth]{assets/figure_p33_17.png}
\end{figure}

(2) 无权\textbf{BCD}码：每位数码无确定的位权，例如：余3码.

余3码的编码规律为: 在8421\textbf{BCD}码上加0011,
例 6的余3码为: 0110+0011=1001

2. 格雷码(\textbf{Gray}码)
格雷码为无权码,特点为：相邻两个代码之间仅有一

位不同,其余各位均相同.具有这种特点的代码称为循环码,
格雷码是循环码.

\begin{figure}[htbp]
  \centering
  \includegraphics[width=\linewidth]{assets/figure_p34_1.png}
\end{figure}

\begin{figure}[htbp]
  \centering
  \includegraphics[width=\linewidth]{assets/figure_p34_2.png}
\end{figure}

\begin{figure}[htbp]
  \centering
  \includegraphics[width=\linewidth]{assets/figure_p34_3.png}
\end{figure}

\begin{figure}[htbp]
  \centering
  \includegraphics[width=\linewidth]{assets/figure_p34_4.png}
\end{figure}

\begin{figure}[htbp]
  \centering
  \includegraphics[width=\linewidth]{assets/figure_p34_5.png}
\end{figure}

\begin{figure}[htbp]
  \centering
  \includegraphics[width=\linewidth]{assets/figure_p34_6.png}
\end{figure}

\begin{figure}[htbp]
  \centering
  \includegraphics[width=\linewidth]{assets/figure_p34_7.png}
\end{figure}

\begin{figure}[htbp]
  \centering
  \includegraphics[width=\linewidth]{assets/figure_p34_8.png}
\end{figure}

\begin{figure}[htbp]
  \centering
  \includegraphics[width=\linewidth]{assets/figure_p34_9.png}
\end{figure}

\begin{figure}[htbp]
  \centering
  \includegraphics[width=\linewidth]{assets/figure_p34_10.png}
\end{figure}

\begin{figure}[htbp]
  \centering
  \includegraphics[width=\linewidth]{assets/figure_p34_11.png}
\end{figure}

\begin{figure}[htbp]
  \centering
  \includegraphics[width=\linewidth]{assets/figure_p34_12.png}
\end{figure}

\begin{figure}[htbp]
  \centering
  \includegraphics[width=\linewidth]{assets/figure_p34_13.png}
\end{figure}

\begin{figure}[htbp]
  \centering
  \includegraphics[width=\linewidth]{assets/figure_p34_14.png}
\end{figure}

\begin{figure}[htbp]
  \centering
  \includegraphics[width=\linewidth]{assets/figure_p34_15.png}
\end{figure}

\begin{figure}[htbp]
  \centering
  \includegraphics[width=\linewidth]{assets/figure_p34_16.png}
\end{figure}

\begin{figure}[htbp]
  \centering
  \includegraphics[width=\linewidth]{assets/figure_p34_17.png}
\end{figure}

格雷码和四位二进制码之间的关系:

设四位二进制码为\textbf{B}\textsubscript{\textbf{3}}\textbf{B}\textsubscript{\textbf{2}}\textbf{B}\textsubscript{\textbf{1}}\textbf{B}\textsubscript{\textbf{0}},格雷码为\textbf{R}\textsubscript{\textbf{3}}\textbf{R}\textsubscript{\textbf{2}}\textbf{R}\textsubscript{\textbf{1}}\textbf{R}\textsubscript{\textbf{0}},

则

\textbf{R}\textsubscript{\textbf{3}}\textbf{=B}\textsubscript{\textbf{3}}\textbf{,} 其中, 为异或运算符,其运算



\textbf{R}\textsubscript{\textbf{2}}\textbf{=B}\textsubscript{\textbf{3}}\textbf{B}\textsubscript{\textbf{2}} 规则为:若两运算数相同,结果
\textbf{R}\textsubscript{\textbf{1}}\textbf{=B}\textsubscript{\textbf{2}}\textbf{B}\textsubscript{\textbf{1}} 为“0”;两运算数不同,结果为
\textbf{R =B} \textbf{B} “1”.

\textbf{0 1 0}

\begin{figure}[htbp]
  \centering
  \includegraphics[width=\linewidth]{assets/figure_p35_1.png}
\end{figure}

\begin{figure}[htbp]
  \centering
  \includegraphics[width=\linewidth]{assets/figure_p35_2.png}
\end{figure}

\begin{figure}[htbp]
  \centering
  \includegraphics[width=\linewidth]{assets/figure_p35_3.png}
\end{figure}

\begin{figure}[htbp]
  \centering
  \includegraphics[width=\linewidth]{assets/figure_p35_4.png}
\end{figure}

\begin{figure}[htbp]
  \centering
  \includegraphics[width=\linewidth]{assets/figure_p35_5.png}
\end{figure}

\begin{figure}[htbp]
  \centering
  \includegraphics[width=\linewidth]{assets/figure_p35_6.png}
\end{figure}

\begin{figure}[htbp]
  \centering
  \includegraphics[width=\linewidth]{assets/figure_p35_7.png}
\end{figure}

\begin{figure}[htbp]
  \centering
  \includegraphics[width=\linewidth]{assets/figure_p35_8.png}
\end{figure}

\begin{figure}[htbp]
  \centering
  \includegraphics[width=\linewidth]{assets/figure_p35_9.png}
\end{figure}

\begin{figure}[htbp]
  \centering
  \includegraphics[width=\linewidth]{assets/figure_p35_10.png}
\end{figure}

\begin{figure}[htbp]
  \centering
  \includegraphics[width=\linewidth]{assets/figure_p35_11.png}
\end{figure}

\begin{figure}[htbp]
  \centering
  \includegraphics[width=\linewidth]{assets/figure_p35_12.png}
\end{figure}

\begin{figure}[htbp]
  \centering
  \includegraphics[width=\linewidth]{assets/figure_p35_13.png}
\end{figure}

\begin{figure}[htbp]
  \centering
  \includegraphics[width=\linewidth]{assets/figure_p35_14.png}
\end{figure}

\begin{figure}[htbp]
  \centering
  \includegraphics[width=\linewidth]{assets/figure_p35_15.png}
\end{figure}

\begin{figure}[htbp]
  \centering
  \includegraphics[width=\linewidth]{assets/figure_p35_16.png}
\end{figure}

\begin{figure}[htbp]
  \centering
  \includegraphics[width=\linewidth]{assets/figure_p35_17.png}
\end{figure}

\section{计算机字符集}

• 各种文字和符号的总称，包括各国家文字、标点符
号、图形符号、数字等。
• 常见字符集名称

⊿ ASCII字符集
⊿ GB2312字符集
⊿ BIG5字符集
⊿ GB18030字符集
⊿ Unicode字符集

\begin{figure}[htbp]
  \centering
  \includegraphics[width=\linewidth]{assets/figure_p36_1.png}
\end{figure}

\begin{figure}[htbp]
  \centering
  \includegraphics[width=\linewidth]{assets/figure_p36_2.png}
\end{figure}

\begin{figure}[htbp]
  \centering
  \includegraphics[width=\linewidth]{assets/figure_p36_3.png}
\end{figure}

\begin{figure}[htbp]
  \centering
  \includegraphics[width=\linewidth]{assets/figure_p36_4.png}
\end{figure}

\begin{figure}[htbp]
  \centering
  \includegraphics[width=\linewidth]{assets/figure_p36_5.png}
\end{figure}

\begin{figure}[htbp]
  \centering
  \includegraphics[width=\linewidth]{assets/figure_p36_6.png}
\end{figure}

\begin{figure}[htbp]
  \centering
  \includegraphics[width=\linewidth]{assets/figure_p36_7.png}
\end{figure}

\begin{figure}[htbp]
  \centering
  \includegraphics[width=\linewidth]{assets/figure_p36_8.png}
\end{figure}

\begin{figure}[htbp]
  \centering
  \includegraphics[width=\linewidth]{assets/figure_p36_9.png}
\end{figure}

\begin{figure}[htbp]
  \centering
  \includegraphics[width=\linewidth]{assets/figure_p36_10.png}
\end{figure}

\begin{figure}[htbp]
  \centering
  \includegraphics[width=\linewidth]{assets/figure_p36_11.png}
\end{figure}

\begin{figure}[htbp]
  \centering
  \includegraphics[width=\linewidth]{assets/figure_p36_12.png}
\end{figure}

\begin{figure}[htbp]
  \centering
  \includegraphics[width=\linewidth]{assets/figure_p36_13.png}
\end{figure}

\begin{figure}[htbp]
  \centering
  \includegraphics[width=\linewidth]{assets/figure_p36_14.png}
\end{figure}

\begin{figure}[htbp]
  \centering
  \includegraphics[width=\linewidth]{assets/figure_p36_15.png}
\end{figure}

\begin{figure}[htbp]
  \centering
  \includegraphics[width=\linewidth]{assets/figure_p36_16.png}
\end{figure}

\begin{figure}[htbp]
  \centering
  \includegraphics[width=\linewidth]{assets/figure_p36_17.png}
\end{figure}

8.2 逻辑代数基础

研究数字电路的基础为逻辑代数，由英国数学家

\textbf{George Boole}在1847年提出的，逻辑代数也称布尔代数.

\begin{figure}[htbp]
  \centering
  \includegraphics[width=\linewidth]{assets/figure_p38_1.png}
\end{figure}

\begin{figure}[htbp]
  \centering
  \includegraphics[width=\linewidth]{assets/figure_p38_2.png}
\end{figure}

\begin{figure}[htbp]
  \centering
  \includegraphics[width=\linewidth]{assets/figure_p38_3.png}
\end{figure}

\begin{figure}[htbp]
  \centering
  \includegraphics[width=\linewidth]{assets/figure_p38_4.png}
\end{figure}

\begin{figure}[htbp]
  \centering
  \includegraphics[width=\linewidth]{assets/figure_p38_5.png}
\end{figure}

\begin{figure}[htbp]
  \centering
  \includegraphics[width=\linewidth]{assets/figure_p38_6.png}
\end{figure}

\begin{figure}[htbp]
  \centering
  \includegraphics[width=\linewidth]{assets/figure_p38_7.png}
\end{figure}

\begin{figure}[htbp]
  \centering
  \includegraphics[width=\linewidth]{assets/figure_p38_8.png}
\end{figure}

\begin{figure}[htbp]
  \centering
  \includegraphics[width=\linewidth]{assets/figure_p38_9.png}
\end{figure}

\begin{figure}[htbp]
  \centering
  \includegraphics[width=\linewidth]{assets/figure_p38_10.png}
\end{figure}

\begin{figure}[htbp]
  \centering
  \includegraphics[width=\linewidth]{assets/figure_p38_11.png}
\end{figure}

\begin{figure}[htbp]
  \centering
  \includegraphics[width=\linewidth]{assets/figure_p38_12.png}
\end{figure}

\begin{figure}[htbp]
  \centering
  \includegraphics[width=\linewidth]{assets/figure_p38_13.png}
\end{figure}

\begin{figure}[htbp]
  \centering
  \includegraphics[width=\linewidth]{assets/figure_p38_14.png}
\end{figure}

\begin{figure}[htbp]
  \centering
  \includegraphics[width=\linewidth]{assets/figure_p38_15.png}
\end{figure}

\begin{figure}[htbp]
  \centering
  \includegraphics[width=\linewidth]{assets/figure_p38_16.png}
\end{figure}

\begin{figure}[htbp]
  \centering
  \includegraphics[width=\linewidth]{assets/figure_p38_17.png}
\end{figure}

8.2.1 基本逻辑运算

在逻辑代数中,变量常用字母\textbf{A,B,C,……Y,Z, a,b,}

\textbf{c,……x.y.z}等表示，变量的取值只能是“0”或“1”.

逻辑代数中只有三种基本逻辑运算,即“与”、

“或”、“非”。

\begin{figure}[htbp]
  \centering
  \includegraphics[width=\linewidth]{assets/figure_p39_1.png}
\end{figure}

\begin{figure}[htbp]
  \centering
  \includegraphics[width=\linewidth]{assets/figure_p39_2.png}
\end{figure}

\begin{figure}[htbp]
  \centering
  \includegraphics[width=\linewidth]{assets/figure_p39_3.png}
\end{figure}

\begin{figure}[htbp]
  \centering
  \includegraphics[width=\linewidth]{assets/figure_p39_4.png}
\end{figure}

\begin{figure}[htbp]
  \centering
  \includegraphics[width=\linewidth]{assets/figure_p39_5.png}
\end{figure}

\begin{figure}[htbp]
  \centering
  \includegraphics[width=\linewidth]{assets/figure_p39_6.png}
\end{figure}

\begin{figure}[htbp]
  \centering
  \includegraphics[width=\linewidth]{assets/figure_p39_7.png}
\end{figure}

\begin{figure}[htbp]
  \centering
  \includegraphics[width=\linewidth]{assets/figure_p39_8.png}
\end{figure}

\begin{figure}[htbp]
  \centering
  \includegraphics[width=\linewidth]{assets/figure_p39_9.png}
\end{figure}

\begin{figure}[htbp]
  \centering
  \includegraphics[width=\linewidth]{assets/figure_p39_10.png}
\end{figure}

\begin{figure}[htbp]
  \centering
  \includegraphics[width=\linewidth]{assets/figure_p39_11.png}
\end{figure}

\begin{figure}[htbp]
  \centering
  \includegraphics[width=\linewidth]{assets/figure_p39_12.png}
\end{figure}

\begin{figure}[htbp]
  \centering
  \includegraphics[width=\linewidth]{assets/figure_p39_13.png}
\end{figure}

\begin{figure}[htbp]
  \centering
  \includegraphics[width=\linewidth]{assets/figure_p39_14.png}
\end{figure}

\begin{figure}[htbp]
  \centering
  \includegraphics[width=\linewidth]{assets/figure_p39_15.png}
\end{figure}

\begin{figure}[htbp]
  \centering
  \includegraphics[width=\linewidth]{assets/figure_p39_16.png}
\end{figure}

\begin{figure}[htbp]
  \centering
  \includegraphics[width=\linewidth]{assets/figure_p39_17.png}
\end{figure}

1. 与逻辑运算
定义：只有决定一事件的全部条件都具备时，这件

事才成立；如果有一个或一个以上条件不具备，则这件事
就不成立。这样的因果关系称为“与”逻辑关系。

与逻辑电路状态表

A B

开关\textbf{A}状态 开关 \textbf{B}状态 灯\textbf{F}状态
断 断 灭
断 合 灭
E F 合 断 灭
合 合 亮

与逻辑电路

\begin{figure}[htbp]
  \centering
  \includegraphics[width=\linewidth]{assets/figure_p40_1.png}
\end{figure}

\begin{figure}[htbp]
  \centering
  \includegraphics[width=\linewidth]{assets/figure_p40_2.png}
\end{figure}

\begin{figure}[htbp]
  \centering
  \includegraphics[width=\linewidth]{assets/figure_p40_3.png}
\end{figure}

\begin{figure}[htbp]
  \centering
  \includegraphics[width=\linewidth]{assets/figure_p40_4.png}
\end{figure}

\begin{figure}[htbp]
  \centering
  \includegraphics[width=\linewidth]{assets/figure_p40_5.png}
\end{figure}

\begin{figure}[htbp]
  \centering
  \includegraphics[width=\linewidth]{assets/figure_p40_6.png}
\end{figure}

\begin{figure}[htbp]
  \centering
  \includegraphics[width=\linewidth]{assets/figure_p40_7.png}
\end{figure}

\begin{figure}[htbp]
  \centering
  \includegraphics[width=\linewidth]{assets/figure_p40_8.png}
\end{figure}

\begin{figure}[htbp]
  \centering
  \includegraphics[width=\linewidth]{assets/figure_p40_9.png}
\end{figure}

\begin{figure}[htbp]
  \centering
  \includegraphics[width=\linewidth]{assets/figure_p40_10.png}
\end{figure}

\begin{figure}[htbp]
  \centering
  \includegraphics[width=\linewidth]{assets/figure_p40_11.png}
\end{figure}

\begin{figure}[htbp]
  \centering
  \includegraphics[width=\linewidth]{assets/figure_p40_12.png}
\end{figure}

\begin{figure}[htbp]
  \centering
  \includegraphics[width=\linewidth]{assets/figure_p40_13.png}
\end{figure}

\begin{figure}[htbp]
  \centering
  \includegraphics[width=\linewidth]{assets/figure_p40_14.png}
\end{figure}

\begin{figure}[htbp]
  \centering
  \includegraphics[width=\linewidth]{assets/figure_p40_15.png}
\end{figure}

\begin{figure}[htbp]
  \centering
  \includegraphics[width=\linewidth]{assets/figure_p40_16.png}
\end{figure}

\begin{figure}[htbp]
  \centering
  \includegraphics[width=\linewidth]{assets/figure_p40_17.png}
\end{figure}

若将开关断开和灯的熄灭状态用逻辑量“0”表示;将开关
合上和灯亮的状态用逻辑量“1”表示,则上述状态表可表

示为:

与逻辑真值表
\textbf{A B F=A $\cdot$ B}

\textbf{\&}

\textbf{A}

\textbf{0 0 0 F=AB} 与门的逻辑功能概括：
\textbf{0 1 0 B} 1）有“0”出“0”；
\textbf{1 0 0}

2）全“1”出“1”。

\textbf{1 1 1} 与门逻辑符号

\begin{figure}[htbp]
  \centering
  \includegraphics[width=\linewidth]{assets/figure_p41_1.png}
\end{figure}

\begin{figure}[htbp]
  \centering
  \includegraphics[width=\linewidth]{assets/figure_p41_2.png}
\end{figure}

\begin{figure}[htbp]
  \centering
  \includegraphics[width=\linewidth]{assets/figure_p41_3.png}
\end{figure}

\begin{figure}[htbp]
  \centering
  \includegraphics[width=\linewidth]{assets/figure_p41_4.png}
\end{figure}

\begin{figure}[htbp]
  \centering
  \includegraphics[width=\linewidth]{assets/figure_p41_5.png}
\end{figure}

\begin{figure}[htbp]
  \centering
  \includegraphics[width=\linewidth]{assets/figure_p41_6.png}
\end{figure}

\begin{figure}[htbp]
  \centering
  \includegraphics[width=\linewidth]{assets/figure_p41_7.png}
\end{figure}

\begin{figure}[htbp]
  \centering
  \includegraphics[width=\linewidth]{assets/figure_p41_8.png}
\end{figure}

\begin{figure}[htbp]
  \centering
  \includegraphics[width=\linewidth]{assets/figure_p41_9.png}
\end{figure}

\begin{figure}[htbp]
  \centering
  \includegraphics[width=\linewidth]{assets/figure_p41_10.png}
\end{figure}

\begin{figure}[htbp]
  \centering
  \includegraphics[width=\linewidth]{assets/figure_p41_11.png}
\end{figure}

\begin{figure}[htbp]
  \centering
  \includegraphics[width=\linewidth]{assets/figure_p41_12.png}
\end{figure}

\begin{figure}[htbp]
  \centering
  \includegraphics[width=\linewidth]{assets/figure_p41_13.png}
\end{figure}

\begin{figure}[htbp]
  \centering
  \includegraphics[width=\linewidth]{assets/figure_p41_14.png}
\end{figure}

\begin{figure}[htbp]
  \centering
  \includegraphics[width=\linewidth]{assets/figure_p41_15.png}
\end{figure}

\begin{figure}[htbp]
  \centering
  \includegraphics[width=\linewidth]{assets/figure_p41_16.png}
\end{figure}

\begin{figure}[htbp]
  \centering
  \includegraphics[width=\linewidth]{assets/figure_p41_17.png}
\end{figure}

2. 或逻辑运算

定义：在决定一事件的各种条件中,只要有一个或一

个以上条件具备时，这件事就成立;只有所有的条件都不

具备时,这件事就不成立.这样的因果关系称为“或”逻辑
关系。 A\textsubscript{或逻辑真值表}

\textbf{A B F=A+ B}

\textbf{0 0 0}
B \textbf{0 1 1}

E F

\textbf{1 0 1}
\textbf{1 1 1}

或逻辑电路

\begin{figure}[htbp]
  \centering
  \includegraphics[width=\linewidth]{assets/figure_p42_1.png}
\end{figure}

\begin{figure}[htbp]
  \centering
  \includegraphics[width=\linewidth]{assets/figure_p42_2.png}
\end{figure}

\begin{figure}[htbp]
  \centering
  \includegraphics[width=\linewidth]{assets/figure_p42_3.png}
\end{figure}

\begin{figure}[htbp]
  \centering
  \includegraphics[width=\linewidth]{assets/figure_p42_4.png}
\end{figure}

\begin{figure}[htbp]
  \centering
  \includegraphics[width=\linewidth]{assets/figure_p42_5.png}
\end{figure}

\begin{figure}[htbp]
  \centering
  \includegraphics[width=\linewidth]{assets/figure_p42_6.png}
\end{figure}

\begin{figure}[htbp]
  \centering
  \includegraphics[width=\linewidth]{assets/figure_p42_7.png}
\end{figure}

\begin{figure}[htbp]
  \centering
  \includegraphics[width=\linewidth]{assets/figure_p42_8.png}
\end{figure}

\begin{figure}[htbp]
  \centering
  \includegraphics[width=\linewidth]{assets/figure_p42_9.png}
\end{figure}

\begin{figure}[htbp]
  \centering
  \includegraphics[width=\linewidth]{assets/figure_p42_10.png}
\end{figure}

\begin{figure}[htbp]
  \centering
  \includegraphics[width=\linewidth]{assets/figure_p42_11.png}
\end{figure}

\begin{figure}[htbp]
  \centering
  \includegraphics[width=\linewidth]{assets/figure_p42_12.png}
\end{figure}

\begin{figure}[htbp]
  \centering
  \includegraphics[width=\linewidth]{assets/figure_p42_13.png}
\end{figure}

\begin{figure}[htbp]
  \centering
  \includegraphics[width=\linewidth]{assets/figure_p42_14.png}
\end{figure}

\begin{figure}[htbp]
  \centering
  \includegraphics[width=\linewidth]{assets/figure_p42_15.png}
\end{figure}

\begin{figure}[htbp]
  \centering
  \includegraphics[width=\linewidth]{assets/figure_p42_16.png}
\end{figure}

\begin{figure}[htbp]
  \centering
  \includegraphics[width=\linewidth]{assets/figure_p42_17.png}
\end{figure}

$\textbf{\ge 1}$

\textbf{A}
\textbf{F=A+B} 或门的逻辑功能概括为:
\textbf{B} 1) 有“1”出“1”;

2) 全“0” 出“0”.
或门逻辑符号

3. 非逻辑运算

定义:假定事件\textbf{F}成立与否同条件\textbf{A}的具备与否有关,

若\textbf{A}具备,则\textbf{F}不成立;若\textbf{A}不具备,则\textbf{F}成立.\textbf{F}和\textbf{A}之间的这

种因果关系称为“非”逻辑关系.

\begin{figure}[htbp]
  \centering
  \includegraphics[width=\linewidth]{assets/figure_p43_1.png}
\end{figure}

\begin{figure}[htbp]
  \centering
  \includegraphics[width=\linewidth]{assets/figure_p43_2.png}
\end{figure}

\begin{figure}[htbp]
  \centering
  \includegraphics[width=\linewidth]{assets/figure_p43_3.png}
\end{figure}

\begin{figure}[htbp]
  \centering
  \includegraphics[width=\linewidth]{assets/figure_p43_4.png}
\end{figure}

\begin{figure}[htbp]
  \centering
  \includegraphics[width=\linewidth]{assets/figure_p43_5.png}
\end{figure}

\begin{figure}[htbp]
  \centering
  \includegraphics[width=\linewidth]{assets/figure_p43_6.png}
\end{figure}

\begin{figure}[htbp]
  \centering
  \includegraphics[width=\linewidth]{assets/figure_p43_7.png}
\end{figure}

\begin{figure}[htbp]
  \centering
  \includegraphics[width=\linewidth]{assets/figure_p43_8.png}
\end{figure}

\begin{figure}[htbp]
  \centering
  \includegraphics[width=\linewidth]{assets/figure_p43_9.png}
\end{figure}

\begin{figure}[htbp]
  \centering
  \includegraphics[width=\linewidth]{assets/figure_p43_10.png}
\end{figure}

\begin{figure}[htbp]
  \centering
  \includegraphics[width=\linewidth]{assets/figure_p43_11.png}
\end{figure}

\begin{figure}[htbp]
  \centering
  \includegraphics[width=\linewidth]{assets/figure_p43_12.png}
\end{figure}

\begin{figure}[htbp]
  \centering
  \includegraphics[width=\linewidth]{assets/figure_p43_13.png}
\end{figure}

\begin{figure}[htbp]
  \centering
  \includegraphics[width=\linewidth]{assets/figure_p43_14.png}
\end{figure}

\begin{figure}[htbp]
  \centering
  \includegraphics[width=\linewidth]{assets/figure_p43_15.png}
\end{figure}

\begin{figure}[htbp]
  \centering
  \includegraphics[width=\linewidth]{assets/figure_p43_16.png}
\end{figure}

\begin{figure}[htbp]
  \centering
  \includegraphics[width=\linewidth]{assets/figure_p43_17.png}
\end{figure}

非逻辑真值表

\textbf{1}

\textbf{A F=A}

E A F \textbf{A F=A}

\textbf{0 1}
\textbf{1 0}
非门逻辑符号

非逻辑电路

\textbf{•} 与门和或门均可以有多个输入端\textbf{.}

\begin{figure}[htbp]
  \centering
  \includegraphics[width=\linewidth]{assets/figure_p44_1.png}
\end{figure}

\begin{figure}[htbp]
  \centering
  \includegraphics[width=\linewidth]{assets/figure_p44_2.png}
\end{figure}

\begin{figure}[htbp]
  \centering
  \includegraphics[width=\linewidth]{assets/figure_p44_3.png}
\end{figure}

\begin{figure}[htbp]
  \centering
  \includegraphics[width=\linewidth]{assets/figure_p44_4.png}
\end{figure}

\begin{figure}[htbp]
  \centering
  \includegraphics[width=\linewidth]{assets/figure_p44_5.png}
\end{figure}

\begin{figure}[htbp]
  \centering
  \includegraphics[width=\linewidth]{assets/figure_p44_6.png}
\end{figure}

\begin{figure}[htbp]
  \centering
  \includegraphics[width=\linewidth]{assets/figure_p44_7.png}
\end{figure}

\begin{figure}[htbp]
  \centering
  \includegraphics[width=\linewidth]{assets/figure_p44_8.png}
\end{figure}

\begin{figure}[htbp]
  \centering
  \includegraphics[width=\linewidth]{assets/figure_p44_9.png}
\end{figure}

\begin{figure}[htbp]
  \centering
  \includegraphics[width=\linewidth]{assets/figure_p44_10.png}
\end{figure}

\begin{figure}[htbp]
  \centering
  \includegraphics[width=\linewidth]{assets/figure_p44_11.png}
\end{figure}

\begin{figure}[htbp]
  \centering
  \includegraphics[width=\linewidth]{assets/figure_p44_12.png}
\end{figure}

\begin{figure}[htbp]
  \centering
  \includegraphics[width=\linewidth]{assets/figure_p44_13.png}
\end{figure}

\begin{figure}[htbp]
  \centering
  \includegraphics[width=\linewidth]{assets/figure_p44_14.png}
\end{figure}

\begin{figure}[htbp]
  \centering
  \includegraphics[width=\linewidth]{assets/figure_p44_15.png}
\end{figure}

\begin{figure}[htbp]
  \centering
  \includegraphics[width=\linewidth]{assets/figure_p44_16.png}
\end{figure}

\begin{figure}[htbp]
  \centering
  \includegraphics[width=\linewidth]{assets/figure_p44_17.png}
\end{figure}

8.2.2 复合逻辑运算

1. 与非逻辑 (将与逻辑和非逻辑组合而成)

与非逻辑真值表
\textbf{A B F=A $\cdot$ B}\textsuperscript{\textbf{\&}}

\textbf{A}
\textbf{F=AB}

\textbf{0 0 1 B}
\textbf{0 1 1}
\textbf{1 0 1} 与非门逻辑符号
\textbf{1 1 0}

\begin{figure}[htbp]
  \centering
  \includegraphics[width=\linewidth]{assets/figure_p45_1.png}
\end{figure}

\begin{figure}[htbp]
  \centering
  \includegraphics[width=\linewidth]{assets/figure_p45_2.png}
\end{figure}

\begin{figure}[htbp]
  \centering
  \includegraphics[width=\linewidth]{assets/figure_p45_3.png}
\end{figure}

\begin{figure}[htbp]
  \centering
  \includegraphics[width=\linewidth]{assets/figure_p45_4.png}
\end{figure}

\begin{figure}[htbp]
  \centering
  \includegraphics[width=\linewidth]{assets/figure_p45_5.png}
\end{figure}

\begin{figure}[htbp]
  \centering
  \includegraphics[width=\linewidth]{assets/figure_p45_6.png}
\end{figure}

\begin{figure}[htbp]
  \centering
  \includegraphics[width=\linewidth]{assets/figure_p45_7.png}
\end{figure}

\begin{figure}[htbp]
  \centering
  \includegraphics[width=\linewidth]{assets/figure_p45_8.png}
\end{figure}

\begin{figure}[htbp]
  \centering
  \includegraphics[width=\linewidth]{assets/figure_p45_9.png}
\end{figure}

\begin{figure}[htbp]
  \centering
  \includegraphics[width=\linewidth]{assets/figure_p45_10.png}
\end{figure}

\begin{figure}[htbp]
  \centering
  \includegraphics[width=\linewidth]{assets/figure_p45_11.png}
\end{figure}

\begin{figure}[htbp]
  \centering
  \includegraphics[width=\linewidth]{assets/figure_p45_12.png}
\end{figure}

\begin{figure}[htbp]
  \centering
  \includegraphics[width=\linewidth]{assets/figure_p45_13.png}
\end{figure}

\begin{figure}[htbp]
  \centering
  \includegraphics[width=\linewidth]{assets/figure_p45_14.png}
\end{figure}

\begin{figure}[htbp]
  \centering
  \includegraphics[width=\linewidth]{assets/figure_p45_15.png}
\end{figure}

\begin{figure}[htbp]
  \centering
  \includegraphics[width=\linewidth]{assets/figure_p45_16.png}
\end{figure}

\begin{figure}[htbp]
  \centering
  \includegraphics[width=\linewidth]{assets/figure_p45_17.png}
\end{figure}

2. 或非逻辑 (将或逻辑和非逻辑组合而成)

或非逻辑真值表
\textbf{A B F=A +B}

$\textbf{\ge 1}$

\textbf{A}

\textbf{0 0 1 F=A+B}
\textbf{0 1 0 B}
\textbf{1 0 0}
\textbf{1 1 0} 或非门逻辑符号

\begin{figure}[htbp]
  \centering
  \includegraphics[width=\linewidth]{assets/figure_p46_1.png}
\end{figure}

\begin{figure}[htbp]
  \centering
  \includegraphics[width=\linewidth]{assets/figure_p46_2.png}
\end{figure}

\begin{figure}[htbp]
  \centering
  \includegraphics[width=\linewidth]{assets/figure_p46_3.png}
\end{figure}

\begin{figure}[htbp]
  \centering
  \includegraphics[width=\linewidth]{assets/figure_p46_4.png}
\end{figure}

\begin{figure}[htbp]
  \centering
  \includegraphics[width=\linewidth]{assets/figure_p46_5.png}
\end{figure}

\begin{figure}[htbp]
  \centering
  \includegraphics[width=\linewidth]{assets/figure_p46_6.png}
\end{figure}

\begin{figure}[htbp]
  \centering
  \includegraphics[width=\linewidth]{assets/figure_p46_7.png}
\end{figure}

\begin{figure}[htbp]
  \centering
  \includegraphics[width=\linewidth]{assets/figure_p46_8.png}
\end{figure}

\begin{figure}[htbp]
  \centering
  \includegraphics[width=\linewidth]{assets/figure_p46_9.png}
\end{figure}

\begin{figure}[htbp]
  \centering
  \includegraphics[width=\linewidth]{assets/figure_p46_10.png}
\end{figure}

\begin{figure}[htbp]
  \centering
  \includegraphics[width=\linewidth]{assets/figure_p46_11.png}
\end{figure}

\begin{figure}[htbp]
  \centering
  \includegraphics[width=\linewidth]{assets/figure_p46_12.png}
\end{figure}

\begin{figure}[htbp]
  \centering
  \includegraphics[width=\linewidth]{assets/figure_p46_13.png}
\end{figure}

\begin{figure}[htbp]
  \centering
  \includegraphics[width=\linewidth]{assets/figure_p46_14.png}
\end{figure}

\begin{figure}[htbp]
  \centering
  \includegraphics[width=\linewidth]{assets/figure_p46_15.png}
\end{figure}

\begin{figure}[htbp]
  \centering
  \includegraphics[width=\linewidth]{assets/figure_p46_16.png}
\end{figure}

\begin{figure}[htbp]
  \centering
  \includegraphics[width=\linewidth]{assets/figure_p46_17.png}
\end{figure}

3.与或非逻辑 (由与、或、非三种逻辑组合而成）

与或非逻辑函数式：
F=AB+CD\textsubscript{\textbf{F=AB+CD}}

$\textbf{\ge 1}$

与或非门的逻辑符号 \textbf{\&}

\textbf{A B C D}

\begin{figure}[htbp]
  \centering
  \includegraphics[width=\linewidth]{assets/figure_p47_1.png}
\end{figure}

\begin{figure}[htbp]
  \centering
  \includegraphics[width=\linewidth]{assets/figure_p47_2.png}
\end{figure}

\begin{figure}[htbp]
  \centering
  \includegraphics[width=\linewidth]{assets/figure_p47_3.png}
\end{figure}

\begin{figure}[htbp]
  \centering
  \includegraphics[width=\linewidth]{assets/figure_p47_4.png}
\end{figure}

\begin{figure}[htbp]
  \centering
  \includegraphics[width=\linewidth]{assets/figure_p47_5.png}
\end{figure}

\begin{figure}[htbp]
  \centering
  \includegraphics[width=\linewidth]{assets/figure_p47_6.png}
\end{figure}

\begin{figure}[htbp]
  \centering
  \includegraphics[width=\linewidth]{assets/figure_p47_7.png}
\end{figure}

\begin{figure}[htbp]
  \centering
  \includegraphics[width=\linewidth]{assets/figure_p47_8.png}
\end{figure}

\begin{figure}[htbp]
  \centering
  \includegraphics[width=\linewidth]{assets/figure_p47_9.png}
\end{figure}

\begin{figure}[htbp]
  \centering
  \includegraphics[width=\linewidth]{assets/figure_p47_10.png}
\end{figure}

\begin{figure}[htbp]
  \centering
  \includegraphics[width=\linewidth]{assets/figure_p47_11.png}
\end{figure}

\begin{figure}[htbp]
  \centering
  \includegraphics[width=\linewidth]{assets/figure_p47_12.png}
\end{figure}

\begin{figure}[htbp]
  \centering
  \includegraphics[width=\linewidth]{assets/figure_p47_13.png}
\end{figure}

\begin{figure}[htbp]
  \centering
  \includegraphics[width=\linewidth]{assets/figure_p47_14.png}
\end{figure}

\begin{figure}[htbp]
  \centering
  \includegraphics[width=\linewidth]{assets/figure_p47_15.png}
\end{figure}

\begin{figure}[htbp]
  \centering
  \includegraphics[width=\linewidth]{assets/figure_p47_16.png}
\end{figure}

\begin{figure}[htbp]
  \centering
  \includegraphics[width=\linewidth]{assets/figure_p47_17.png}
\end{figure}

4.异或逻辑

\[ 异或逻辑的函数式为： F=AB+AB = A  B \]

异或逻辑真值表
\textbf{A B F=A}\textbf{B}

\textbf{=1}

\textbf{A}

\textbf{0 0 0 F=A} \textbf{B}

\textbf{B}

\textbf{0 1 1}
\textbf{1 0 1}

异或门逻辑符号

\textbf{1 1 0}

异或逻辑的功能为: 1) 相同得“0”;
2) 相异得“1”.

\begin{figure}[htbp]
  \centering
  \includegraphics[width=\linewidth]{assets/figure_p48_1.png}
\end{figure}

\begin{figure}[htbp]
  \centering
  \includegraphics[width=\linewidth]{assets/figure_p48_2.png}
\end{figure}

\begin{figure}[htbp]
  \centering
  \includegraphics[width=\linewidth]{assets/figure_p48_3.png}
\end{figure}

\begin{figure}[htbp]
  \centering
  \includegraphics[width=\linewidth]{assets/figure_p48_4.png}
\end{figure}

\begin{figure}[htbp]
  \centering
  \includegraphics[width=\linewidth]{assets/figure_p48_5.png}
\end{figure}

\begin{figure}[htbp]
  \centering
  \includegraphics[width=\linewidth]{assets/figure_p48_6.png}
\end{figure}

\begin{figure}[htbp]
  \centering
  \includegraphics[width=\linewidth]{assets/figure_p48_7.png}
\end{figure}

\begin{figure}[htbp]
  \centering
  \includegraphics[width=\linewidth]{assets/figure_p48_8.png}
\end{figure}

\begin{figure}[htbp]
  \centering
  \includegraphics[width=\linewidth]{assets/figure_p48_9.png}
\end{figure}

\begin{figure}[htbp]
  \centering
  \includegraphics[width=\linewidth]{assets/figure_p48_10.png}
\end{figure}

\begin{figure}[htbp]
  \centering
  \includegraphics[width=\linewidth]{assets/figure_p48_11.png}
\end{figure}

\begin{figure}[htbp]
  \centering
  \includegraphics[width=\linewidth]{assets/figure_p48_12.png}
\end{figure}

\begin{figure}[htbp]
  \centering
  \includegraphics[width=\linewidth]{assets/figure_p48_13.png}
\end{figure}

\begin{figure}[htbp]
  \centering
  \includegraphics[width=\linewidth]{assets/figure_p48_14.png}
\end{figure}

\begin{figure}[htbp]
  \centering
  \includegraphics[width=\linewidth]{assets/figure_p48_15.png}
\end{figure}

\begin{figure}[htbp]
  \centering
  \includegraphics[width=\linewidth]{assets/figure_p48_16.png}
\end{figure}

\begin{figure}[htbp]
  \centering
  \includegraphics[width=\linewidth]{assets/figure_p48_17.png}
\end{figure}

5.同或逻辑

\[ 同或逻辑式为 :F = A B + A B =A . B \]

同或逻辑 真值表
\textbf{A B F=A . B}

\textbf{=}

\textbf{A}

\textbf{F=A . B}
\textbf{0 0 1}

\textbf{B}

\textbf{0 1 0}
\textbf{1 0 0} 同或门逻辑符号
\textbf{1 1 1}

对照异或和同或逻辑真值表,可以发现: 同或和异或互
为反函数,即:
\textbf{A} \textbf{B = A . B}

\begin{figure}[htbp]
  \centering
  \includegraphics[width=\linewidth]{assets/figure_p49_1.png}
\end{figure}

\begin{figure}[htbp]
  \centering
  \includegraphics[width=\linewidth]{assets/figure_p49_2.png}
\end{figure}

\begin{figure}[htbp]
  \centering
  \includegraphics[width=\linewidth]{assets/figure_p49_3.png}
\end{figure}

\begin{figure}[htbp]
  \centering
  \includegraphics[width=\linewidth]{assets/figure_p49_4.png}
\end{figure}

\begin{figure}[htbp]
  \centering
  \includegraphics[width=\linewidth]{assets/figure_p49_5.png}
\end{figure}

\begin{figure}[htbp]
  \centering
  \includegraphics[width=\linewidth]{assets/figure_p49_6.png}
\end{figure}

\begin{figure}[htbp]
  \centering
  \includegraphics[width=\linewidth]{assets/figure_p49_7.png}
\end{figure}

\begin{figure}[htbp]
  \centering
  \includegraphics[width=\linewidth]{assets/figure_p49_8.png}
\end{figure}

\begin{figure}[htbp]
  \centering
  \includegraphics[width=\linewidth]{assets/figure_p49_9.png}
\end{figure}

\begin{figure}[htbp]
  \centering
  \includegraphics[width=\linewidth]{assets/figure_p49_10.png}
\end{figure}

\begin{figure}[htbp]
  \centering
  \includegraphics[width=\linewidth]{assets/figure_p49_11.png}
\end{figure}

\begin{figure}[htbp]
  \centering
  \includegraphics[width=\linewidth]{assets/figure_p49_12.png}
\end{figure}

\begin{figure}[htbp]
  \centering
  \includegraphics[width=\linewidth]{assets/figure_p49_13.png}
\end{figure}

\begin{figure}[htbp]
  \centering
  \includegraphics[width=\linewidth]{assets/figure_p49_14.png}
\end{figure}

\begin{figure}[htbp]
  \centering
  \includegraphics[width=\linewidth]{assets/figure_p49_15.png}
\end{figure}

\begin{figure}[htbp]
  \centering
  \includegraphics[width=\linewidth]{assets/figure_p49_16.png}
\end{figure}

\begin{figure}[htbp]
  \centering
  \includegraphics[width=\linewidth]{assets/figure_p49_17.png}
\end{figure}

• 表\textbf{8.13}给出了门电路的几种表示方法，本
课程中，均采用“国标”。国外流行的电
路符号常见于外文书籍中，特别在我国引
进的一些计算机辅助分析和设计软件中，
常使用这些符号。

\begin{figure}[htbp]
  \centering
  \includegraphics[width=\linewidth]{assets/figure_p50_1.png}
\end{figure}

\begin{figure}[htbp]
  \centering
  \includegraphics[width=\linewidth]{assets/figure_p50_2.png}
\end{figure}

\begin{figure}[htbp]
  \centering
  \includegraphics[width=\linewidth]{assets/figure_p50_3.png}
\end{figure}

\begin{figure}[htbp]
  \centering
  \includegraphics[width=\linewidth]{assets/figure_p50_4.png}
\end{figure}

\begin{figure}[htbp]
  \centering
  \includegraphics[width=\linewidth]{assets/figure_p50_5.png}
\end{figure}

\begin{figure}[htbp]
  \centering
  \includegraphics[width=\linewidth]{assets/figure_p50_6.png}
\end{figure}

\begin{figure}[htbp]
  \centering
  \includegraphics[width=\linewidth]{assets/figure_p50_7.png}
\end{figure}

\begin{figure}[htbp]
  \centering
  \includegraphics[width=\linewidth]{assets/figure_p50_8.png}
\end{figure}

\begin{figure}[htbp]
  \centering
  \includegraphics[width=\linewidth]{assets/figure_p50_9.png}
\end{figure}

\begin{figure}[htbp]
  \centering
  \includegraphics[width=\linewidth]{assets/figure_p50_10.png}
\end{figure}

\begin{figure}[htbp]
  \centering
  \includegraphics[width=\linewidth]{assets/figure_p50_11.png}
\end{figure}

\begin{figure}[htbp]
  \centering
  \includegraphics[width=\linewidth]{assets/figure_p50_12.png}
\end{figure}

\begin{figure}[htbp]
  \centering
  \includegraphics[width=\linewidth]{assets/figure_p50_13.png}
\end{figure}

\begin{figure}[htbp]
  \centering
  \includegraphics[width=\linewidth]{assets/figure_p50_14.png}
\end{figure}

\begin{figure}[htbp]
  \centering
  \includegraphics[width=\linewidth]{assets/figure_p50_15.png}
\end{figure}

\begin{figure}[htbp]
  \centering
  \includegraphics[width=\linewidth]{assets/figure_p50_16.png}
\end{figure}

\begin{figure}[htbp]
  \centering
  \includegraphics[width=\linewidth]{assets/figure_p50_17.png}
\end{figure}

8.2.3 逻辑电平及正、负逻辑

门电路的输入、输出为二值信号,用“0”和“1”表
示.这里的“0”、“1”一般用两个不同电平值来表示.

若用高电平V\textsubscript{H}表示逻辑“1”,用低电平V\textsubscript{L}表示逻辑

“0”,则称为正逻辑约定,简称正逻辑;

若用高电平V\textsubscript{H}表示逻辑“0”,用低电平V\textsubscript{L}表示逻辑

“1”,则称为负逻辑约定,简称负逻辑.

\begin{figure}[htbp]
  \centering
  \includegraphics[width=\linewidth]{assets/figure_p51_1.png}
\end{figure}

\begin{figure}[htbp]
  \centering
  \includegraphics[width=\linewidth]{assets/figure_p51_2.png}
\end{figure}

\begin{figure}[htbp]
  \centering
  \includegraphics[width=\linewidth]{assets/figure_p51_3.png}
\end{figure}

\begin{figure}[htbp]
  \centering
  \includegraphics[width=\linewidth]{assets/figure_p51_4.png}
\end{figure}

\begin{figure}[htbp]
  \centering
  \includegraphics[width=\linewidth]{assets/figure_p51_5.png}
\end{figure}

\begin{figure}[htbp]
  \centering
  \includegraphics[width=\linewidth]{assets/figure_p51_6.png}
\end{figure}

\begin{figure}[htbp]
  \centering
  \includegraphics[width=\linewidth]{assets/figure_p51_7.png}
\end{figure}

\begin{figure}[htbp]
  \centering
  \includegraphics[width=\linewidth]{assets/figure_p51_8.png}
\end{figure}

\begin{figure}[htbp]
  \centering
  \includegraphics[width=\linewidth]{assets/figure_p51_9.png}
\end{figure}

\begin{figure}[htbp]
  \centering
  \includegraphics[width=\linewidth]{assets/figure_p51_10.png}
\end{figure}

\begin{figure}[htbp]
  \centering
  \includegraphics[width=\linewidth]{assets/figure_p51_11.png}
\end{figure}

\begin{figure}[htbp]
  \centering
  \includegraphics[width=\linewidth]{assets/figure_p51_12.png}
\end{figure}

\begin{figure}[htbp]
  \centering
  \includegraphics[width=\linewidth]{assets/figure_p51_13.png}
\end{figure}

\begin{figure}[htbp]
  \centering
  \includegraphics[width=\linewidth]{assets/figure_p51_14.png}
\end{figure}

\begin{figure}[htbp]
  \centering
  \includegraphics[width=\linewidth]{assets/figure_p51_15.png}
\end{figure}

\begin{figure}[htbp]
  \centering
  \includegraphics[width=\linewidth]{assets/figure_p51_16.png}
\end{figure}

\begin{figure}[htbp]
  \centering
  \includegraphics[width=\linewidth]{assets/figure_p51_17.png}
\end{figure}

在本课程中,如不作特殊说明,一般都采用正逻辑表示.

\textbf{V}\textsubscript{\textbf{H}}和\textbf{V}\textsubscript{\textbf{L}}的具体值,由所使用的集成电路品种以及所

加电源电压而定,有两种常用的集成电路:

1) \textbf{TTL}电路,电源电压为\textbf{5}伏,\textbf{V}\textsubscript{\textbf{H}}约为\textbf{3V}左右,\textbf{V}\textsubscript{\textbf{L}}约为

\textbf{0.2}伏左右;

2) \textbf{CMOS}电路,电源电压范围较宽,\textbf{CMOS4000}系列

的电源电压\textbf{V}\textsubscript{\textbf{DD}}为\textbf{3\textasciitilde{}18}伏. \textbf{CMOS}电路的\textbf{V}\textsubscript{\textbf{H}}约为\textbf{0.9 V}\textsubscript{\textbf{DD}},
而\textbf{V}\textsubscript{\textbf{L}}约为\textbf{0}伏左右.

\begin{figure}[htbp]
  \centering
  \includegraphics[width=\linewidth]{assets/figure_p52_1.png}
\end{figure}

\begin{figure}[htbp]
  \centering
  \includegraphics[width=\linewidth]{assets/figure_p52_2.png}
\end{figure}

\begin{figure}[htbp]
  \centering
  \includegraphics[width=\linewidth]{assets/figure_p52_3.png}
\end{figure}

\begin{figure}[htbp]
  \centering
  \includegraphics[width=\linewidth]{assets/figure_p52_4.png}
\end{figure}

\begin{figure}[htbp]
  \centering
  \includegraphics[width=\linewidth]{assets/figure_p52_5.png}
\end{figure}

\begin{figure}[htbp]
  \centering
  \includegraphics[width=\linewidth]{assets/figure_p52_6.png}
\end{figure}

\begin{figure}[htbp]
  \centering
  \includegraphics[width=\linewidth]{assets/figure_p52_7.png}
\end{figure}

\begin{figure}[htbp]
  \centering
  \includegraphics[width=\linewidth]{assets/figure_p52_8.png}
\end{figure}

\begin{figure}[htbp]
  \centering
  \includegraphics[width=\linewidth]{assets/figure_p52_9.png}
\end{figure}

\begin{figure}[htbp]
  \centering
  \includegraphics[width=\linewidth]{assets/figure_p52_10.png}
\end{figure}

\begin{figure}[htbp]
  \centering
  \includegraphics[width=\linewidth]{assets/figure_p52_11.png}
\end{figure}

\begin{figure}[htbp]
  \centering
  \includegraphics[width=\linewidth]{assets/figure_p52_12.png}
\end{figure}

\begin{figure}[htbp]
  \centering
  \includegraphics[width=\linewidth]{assets/figure_p52_13.png}
\end{figure}

\begin{figure}[htbp]
  \centering
  \includegraphics[width=\linewidth]{assets/figure_p52_14.png}
\end{figure}

\begin{figure}[htbp]
  \centering
  \includegraphics[width=\linewidth]{assets/figure_p52_15.png}
\end{figure}

\begin{figure}[htbp]
  \centering
  \includegraphics[width=\linewidth]{assets/figure_p52_16.png}
\end{figure}

\begin{figure}[htbp]
  \centering
  \includegraphics[width=\linewidth]{assets/figure_p52_17.png}
\end{figure}

对一个特定的逻辑门,采用不同的逻辑表示时,其门的

名称也就不同.

正负逻辑转换举例
电平真值表 正逻辑\textbf{(}与非门\textbf{)} 负逻辑\textbf{(}或非门\textbf{)}
\textbf{Vi1 Vi2 Vo A B Y A B Y}
\textbf{VL VL VH 0 0 1 1 1 0}
\textbf{VL VH VH 0 1 1 1 0 0}
\textbf{VH VL VH 1 0 1 0 1 0}
\textbf{VH VH VL 1 1 0 0 0 1}

\begin{figure}[htbp]
  \centering
  \includegraphics[width=\linewidth]{assets/figure_p53_1.png}
\end{figure}

\begin{figure}[htbp]
  \centering
  \includegraphics[width=\linewidth]{assets/figure_p53_2.png}
\end{figure}

\begin{figure}[htbp]
  \centering
  \includegraphics[width=\linewidth]{assets/figure_p53_3.png}
\end{figure}

\begin{figure}[htbp]
  \centering
  \includegraphics[width=\linewidth]{assets/figure_p53_4.png}
\end{figure}

\begin{figure}[htbp]
  \centering
  \includegraphics[width=\linewidth]{assets/figure_p53_5.png}
\end{figure}

\begin{figure}[htbp]
  \centering
  \includegraphics[width=\linewidth]{assets/figure_p53_6.png}
\end{figure}

\begin{figure}[htbp]
  \centering
  \includegraphics[width=\linewidth]{assets/figure_p53_7.png}
\end{figure}

\begin{figure}[htbp]
  \centering
  \includegraphics[width=\linewidth]{assets/figure_p53_8.png}
\end{figure}

\begin{figure}[htbp]
  \centering
  \includegraphics[width=\linewidth]{assets/figure_p53_9.png}
\end{figure}

\begin{figure}[htbp]
  \centering
  \includegraphics[width=\linewidth]{assets/figure_p53_10.png}
\end{figure}

\begin{figure}[htbp]
  \centering
  \includegraphics[width=\linewidth]{assets/figure_p53_11.png}
\end{figure}

\begin{figure}[htbp]
  \centering
  \includegraphics[width=\linewidth]{assets/figure_p53_12.png}
\end{figure}

\begin{figure}[htbp]
  \centering
  \includegraphics[width=\linewidth]{assets/figure_p53_13.png}
\end{figure}

\begin{figure}[htbp]
  \centering
  \includegraphics[width=\linewidth]{assets/figure_p53_14.png}
\end{figure}

\begin{figure}[htbp]
  \centering
  \includegraphics[width=\linewidth]{assets/figure_p53_15.png}
\end{figure}

\begin{figure}[htbp]
  \centering
  \includegraphics[width=\linewidth]{assets/figure_p53_16.png}
\end{figure}

\begin{figure}[htbp]
  \centering
  \includegraphics[width=\linewidth]{assets/figure_p53_17.png}
\end{figure}

8.2.4 基本定律和规则

1. 逻辑函数的相等

设有两个逻辑:\textbf{F}\textsubscript{\textbf{1}}\textbf{=f}\textsubscript{\textbf{1}}\textbf{(A}\textsubscript{\textbf{1}}\textbf{,A}\textsubscript{\textbf{2}}\textbf{,…,A}\textsubscript{\textbf{n}}\textbf{)}

\textbf{F}\textsubscript{\textbf{2}}\textbf{=f}\textsubscript{\textbf{2}}\textbf{(A}\textsubscript{\textbf{1}}\textbf{,A}\textsubscript{\textbf{2}}\textbf{,…,A}\textsubscript{\textbf{n}}\textbf{)}

\textbf{n}

如果对于\textbf{A}\textsubscript{\textbf{1}}\textbf{,A}\textsubscript{\textbf{2}}\textbf{,…,A}\textsubscript{\textbf{n}} 的任何一组取值(共\textbf{2} 组),

\textbf{F}\textsubscript{\textbf{1}} 和 \textbf{F}\textsubscript{\textbf{2}}均相等,则称\textbf{F}\textsubscript{\textbf{1}}和 \textbf{F}\textsubscript{\textbf{2}}相等.

因此,如两个函数的真值表相等,则这两个函数一定相等.

\begin{figure}[htbp]
  \centering
  \includegraphics[width=\linewidth]{assets/figure_p54_1.png}
\end{figure}

\begin{figure}[htbp]
  \centering
  \includegraphics[width=\linewidth]{assets/figure_p54_2.png}
\end{figure}

\begin{figure}[htbp]
  \centering
  \includegraphics[width=\linewidth]{assets/figure_p54_3.png}
\end{figure}

\begin{figure}[htbp]
  \centering
  \includegraphics[width=\linewidth]{assets/figure_p54_4.png}
\end{figure}

\begin{figure}[htbp]
  \centering
  \includegraphics[width=\linewidth]{assets/figure_p54_5.png}
\end{figure}

\begin{figure}[htbp]
  \centering
  \includegraphics[width=\linewidth]{assets/figure_p54_6.png}
\end{figure}

\begin{figure}[htbp]
  \centering
  \includegraphics[width=\linewidth]{assets/figure_p54_7.png}
\end{figure}

\begin{figure}[htbp]
  \centering
  \includegraphics[width=\linewidth]{assets/figure_p54_8.png}
\end{figure}

\begin{figure}[htbp]
  \centering
  \includegraphics[width=\linewidth]{assets/figure_p54_9.png}
\end{figure}

\begin{figure}[htbp]
  \centering
  \includegraphics[width=\linewidth]{assets/figure_p54_10.png}
\end{figure}

\begin{figure}[htbp]
  \centering
  \includegraphics[width=\linewidth]{assets/figure_p54_11.png}
\end{figure}

\begin{figure}[htbp]
  \centering
  \includegraphics[width=\linewidth]{assets/figure_p54_12.png}
\end{figure}

\begin{figure}[htbp]
  \centering
  \includegraphics[width=\linewidth]{assets/figure_p54_13.png}
\end{figure}

\begin{figure}[htbp]
  \centering
  \includegraphics[width=\linewidth]{assets/figure_p54_14.png}
\end{figure}

\begin{figure}[htbp]
  \centering
  \includegraphics[width=\linewidth]{assets/figure_p54_15.png}
\end{figure}

\begin{figure}[htbp]
  \centering
  \includegraphics[width=\linewidth]{assets/figure_p54_16.png}
\end{figure}

\begin{figure}[htbp]
  \centering
  \includegraphics[width=\linewidth]{assets/figure_p54_17.png}
\end{figure}

2. 基本定律

① \textbf{0}－\textbf{1}律 \textbf{A $\cdot 0=0 ;$ A+1=1}
②自等律 \textbf{A $\cdot 1=A ;$ A+0=A}
③重迭律 \textbf{A $\cdot A=A ;$ A+A=A}
④互补律 \textbf{A $\cdot$ A=0 ; A+A=1}
⑤交换律 \textbf{A $\cdot$ B= $B \cdot A ;$ A+B=B+A}
⑥结合律 \textbf{A(BC)=(AB)C ; A+(B+C)=(A+B)+C}
⑦分配律 \textbf{A(B+C)=AB+AC ; A+BC=(A+B)(A+C)}

\[ ⑧反演律 A+B=A\cdot B ; AB=A + B \]

=

⑨还原律 \textbf{A = A}

PLAY

\begin{figure}[htbp]
  \centering
  \includegraphics[width=\linewidth]{assets/figure_p55_1.png}
\end{figure}

\begin{figure}[htbp]
  \centering
  \includegraphics[width=\linewidth]{assets/figure_p55_2.png}
\end{figure}

\begin{figure}[htbp]
  \centering
  \includegraphics[width=\linewidth]{assets/figure_p55_3.png}
\end{figure}

\begin{figure}[htbp]
  \centering
  \includegraphics[width=\linewidth]{assets/figure_p55_4.png}
\end{figure}

\begin{figure}[htbp]
  \centering
  \includegraphics[width=\linewidth]{assets/figure_p55_5.png}
\end{figure}

\begin{figure}[htbp]
  \centering
  \includegraphics[width=\linewidth]{assets/figure_p55_6.png}
\end{figure}

\begin{figure}[htbp]
  \centering
  \includegraphics[width=\linewidth]{assets/figure_p55_7.png}
\end{figure}

\begin{figure}[htbp]
  \centering
  \includegraphics[width=\linewidth]{assets/figure_p55_8.png}
\end{figure}

\begin{figure}[htbp]
  \centering
  \includegraphics[width=\linewidth]{assets/figure_p55_9.png}
\end{figure}

\begin{figure}[htbp]
  \centering
  \includegraphics[width=\linewidth]{assets/figure_p55_10.png}
\end{figure}

\begin{figure}[htbp]
  \centering
  \includegraphics[width=\linewidth]{assets/figure_p55_11.png}
\end{figure}

\begin{figure}[htbp]
  \centering
  \includegraphics[width=\linewidth]{assets/figure_p55_12.png}
\end{figure}

\begin{figure}[htbp]
  \centering
  \includegraphics[width=\linewidth]{assets/figure_p55_13.png}
\end{figure}

\begin{figure}[htbp]
  \centering
  \includegraphics[width=\linewidth]{assets/figure_p55_14.png}
\end{figure}

\begin{figure}[htbp]
  \centering
  \includegraphics[width=\linewidth]{assets/figure_p55_15.png}
\end{figure}

\begin{figure}[htbp]
  \centering
  \includegraphics[width=\linewidth]{assets/figure_p55_16.png}
\end{figure}

\begin{figure}[htbp]
  \centering
  \includegraphics[width=\linewidth]{assets/figure_p55_17.png}
\end{figure}

$反演律也称德\cdot 摩根定理,是一个非常有用的定理.$

3. 逻辑代数的三条规则

(1) 代入规则

任何一个含有变量\textbf{x}的等式,如果将所有出现\textbf{x}的位置,

都用一个逻辑函数式\textbf{F}代替,则等式仍然成立.

\begin{figure}[htbp]
  \centering
  \includegraphics[width=\linewidth]{assets/figure_p56_1.png}
\end{figure}

\begin{figure}[htbp]
  \centering
  \includegraphics[width=\linewidth]{assets/figure_p56_2.png}
\end{figure}

\begin{figure}[htbp]
  \centering
  \includegraphics[width=\linewidth]{assets/figure_p56_3.png}
\end{figure}

\begin{figure}[htbp]
  \centering
  \includegraphics[width=\linewidth]{assets/figure_p56_4.png}
\end{figure}

\begin{figure}[htbp]
  \centering
  \includegraphics[width=\linewidth]{assets/figure_p56_5.png}
\end{figure}

\begin{figure}[htbp]
  \centering
  \includegraphics[width=\linewidth]{assets/figure_p56_6.png}
\end{figure}

\begin{figure}[htbp]
  \centering
  \includegraphics[width=\linewidth]{assets/figure_p56_7.png}
\end{figure}

\begin{figure}[htbp]
  \centering
  \includegraphics[width=\linewidth]{assets/figure_p56_8.png}
\end{figure}

\begin{figure}[htbp]
  \centering
  \includegraphics[width=\linewidth]{assets/figure_p56_9.png}
\end{figure}

\begin{figure}[htbp]
  \centering
  \includegraphics[width=\linewidth]{assets/figure_p56_10.png}
\end{figure}

\begin{figure}[htbp]
  \centering
  \includegraphics[width=\linewidth]{assets/figure_p56_11.png}
\end{figure}

\begin{figure}[htbp]
  \centering
  \includegraphics[width=\linewidth]{assets/figure_p56_12.png}
\end{figure}

\begin{figure}[htbp]
  \centering
  \includegraphics[width=\linewidth]{assets/figure_p56_13.png}
\end{figure}

\begin{figure}[htbp]
  \centering
  \includegraphics[width=\linewidth]{assets/figure_p56_14.png}
\end{figure}

\begin{figure}[htbp]
  \centering
  \includegraphics[width=\linewidth]{assets/figure_p56_15.png}
\end{figure}

\begin{figure}[htbp]
  \centering
  \includegraphics[width=\linewidth]{assets/figure_p56_16.png}
\end{figure}

\begin{figure}[htbp]
  \centering
  \includegraphics[width=\linewidth]{assets/figure_p56_17.png}
\end{figure}

例: 已知等式 \textbf{A+B=A $\cdot B$ ,}有函数式\textbf{F=B+C},则
用\textbf{F}代替等式中的\textbf{B,}
有 \textbf{A+(B+C)=A B+C}
即 \textbf{A+B+C=A B C}

由此可以证明反演定律对n变量仍然成立.

(2) 反演规则

\begin{figure}[htbp]
  \centering
  \includegraphics[width=\linewidth]{assets/figure_p57_1.png}
\end{figure}

\begin{figure}[htbp]
  \centering
  \includegraphics[width=\linewidth]{assets/figure_p57_2.png}
\end{figure}

\begin{figure}[htbp]
  \centering
  \includegraphics[width=\linewidth]{assets/figure_p57_3.png}
\end{figure}

\begin{figure}[htbp]
  \centering
  \includegraphics[width=\linewidth]{assets/figure_p57_4.png}
\end{figure}

\begin{figure}[htbp]
  \centering
  \includegraphics[width=\linewidth]{assets/figure_p57_5.png}
\end{figure}

\begin{figure}[htbp]
  \centering
  \includegraphics[width=\linewidth]{assets/figure_p57_6.png}
\end{figure}

\begin{figure}[htbp]
  \centering
  \includegraphics[width=\linewidth]{assets/figure_p57_7.png}
\end{figure}

\begin{figure}[htbp]
  \centering
  \includegraphics[width=\linewidth]{assets/figure_p57_8.png}
\end{figure}

\begin{figure}[htbp]
  \centering
  \includegraphics[width=\linewidth]{assets/figure_p57_9.png}
\end{figure}

\begin{figure}[htbp]
  \centering
  \includegraphics[width=\linewidth]{assets/figure_p57_10.png}
\end{figure}

\begin{figure}[htbp]
  \centering
  \includegraphics[width=\linewidth]{assets/figure_p57_11.png}
\end{figure}

\begin{figure}[htbp]
  \centering
  \includegraphics[width=\linewidth]{assets/figure_p57_12.png}
\end{figure}

\begin{figure}[htbp]
  \centering
  \includegraphics[width=\linewidth]{assets/figure_p57_13.png}
\end{figure}

\begin{figure}[htbp]
  \centering
  \includegraphics[width=\linewidth]{assets/figure_p57_14.png}
\end{figure}

\begin{figure}[htbp]
  \centering
  \includegraphics[width=\linewidth]{assets/figure_p57_15.png}
\end{figure}

\begin{figure}[htbp]
  \centering
  \includegraphics[width=\linewidth]{assets/figure_p57_16.png}
\end{figure}

\begin{figure}[htbp]
  \centering
  \includegraphics[width=\linewidth]{assets/figure_p57_17.png}
\end{figure}

设F为任意逻辑表达式,若将F中所有运算符、常量及

变量作如下变换：

$\textbf{\cdot + 0 1}$ 原变量 反变量

\textbf{+ $\cdot 1 0}$ 反变量 原变量
则所得新的逻辑式即为\textbf{F}的反函数，记为\textbf{F}。

例 已知 \textbf{F=A B + A B,} 根据上述规则可得：

\[ F=(A+B)(A+B) \]

\begin{figure}[htbp]
  \centering
  \includegraphics[width=\linewidth]{assets/figure_p58_1.png}
\end{figure}

\begin{figure}[htbp]
  \centering
  \includegraphics[width=\linewidth]{assets/figure_p58_2.png}
\end{figure}

\begin{figure}[htbp]
  \centering
  \includegraphics[width=\linewidth]{assets/figure_p58_3.png}
\end{figure}

\begin{figure}[htbp]
  \centering
  \includegraphics[width=\linewidth]{assets/figure_p58_4.png}
\end{figure}

\begin{figure}[htbp]
  \centering
  \includegraphics[width=\linewidth]{assets/figure_p58_5.png}
\end{figure}

\begin{figure}[htbp]
  \centering
  \includegraphics[width=\linewidth]{assets/figure_p58_6.png}
\end{figure}

\begin{figure}[htbp]
  \centering
  \includegraphics[width=\linewidth]{assets/figure_p58_7.png}
\end{figure}

\begin{figure}[htbp]
  \centering
  \includegraphics[width=\linewidth]{assets/figure_p58_8.png}
\end{figure}

\begin{figure}[htbp]
  \centering
  \includegraphics[width=\linewidth]{assets/figure_p58_9.png}
\end{figure}

\begin{figure}[htbp]
  \centering
  \includegraphics[width=\linewidth]{assets/figure_p58_10.png}
\end{figure}

\begin{figure}[htbp]
  \centering
  \includegraphics[width=\linewidth]{assets/figure_p58_11.png}
\end{figure}

\begin{figure}[htbp]
  \centering
  \includegraphics[width=\linewidth]{assets/figure_p58_12.png}
\end{figure}

\begin{figure}[htbp]
  \centering
  \includegraphics[width=\linewidth]{assets/figure_p58_13.png}
\end{figure}

\begin{figure}[htbp]
  \centering
  \includegraphics[width=\linewidth]{assets/figure_p58_14.png}
\end{figure}

\begin{figure}[htbp]
  \centering
  \includegraphics[width=\linewidth]{assets/figure_p58_15.png}
\end{figure}

\begin{figure}[htbp]
  \centering
  \includegraphics[width=\linewidth]{assets/figure_p58_16.png}
\end{figure}

\begin{figure}[htbp]
  \centering
  \includegraphics[width=\linewidth]{assets/figure_p58_17.png}
\end{figure}

\[ 例 已知 F=A+B+C+D+E, 则 \]

\textbf{F=A B C D E}

由F求反函数注意：

1）保持原式运算的优先次序；

2）原式中的不属于单变量上的非号不变；

\begin{figure}[htbp]
  \centering
  \includegraphics[width=\linewidth]{assets/figure_p59_1.png}
\end{figure}

\begin{figure}[htbp]
  \centering
  \includegraphics[width=\linewidth]{assets/figure_p59_2.png}
\end{figure}

\begin{figure}[htbp]
  \centering
  \includegraphics[width=\linewidth]{assets/figure_p59_3.png}
\end{figure}

\begin{figure}[htbp]
  \centering
  \includegraphics[width=\linewidth]{assets/figure_p59_4.png}
\end{figure}

\begin{figure}[htbp]
  \centering
  \includegraphics[width=\linewidth]{assets/figure_p59_5.png}
\end{figure}

\begin{figure}[htbp]
  \centering
  \includegraphics[width=\linewidth]{assets/figure_p59_6.png}
\end{figure}

\begin{figure}[htbp]
  \centering
  \includegraphics[width=\linewidth]{assets/figure_p59_7.png}
\end{figure}

\begin{figure}[htbp]
  \centering
  \includegraphics[width=\linewidth]{assets/figure_p59_8.png}
\end{figure}

\begin{figure}[htbp]
  \centering
  \includegraphics[width=\linewidth]{assets/figure_p59_9.png}
\end{figure}

\begin{figure}[htbp]
  \centering
  \includegraphics[width=\linewidth]{assets/figure_p59_10.png}
\end{figure}

\begin{figure}[htbp]
  \centering
  \includegraphics[width=\linewidth]{assets/figure_p59_11.png}
\end{figure}

\begin{figure}[htbp]
  \centering
  \includegraphics[width=\linewidth]{assets/figure_p59_12.png}
\end{figure}

\begin{figure}[htbp]
  \centering
  \includegraphics[width=\linewidth]{assets/figure_p59_13.png}
\end{figure}

\begin{figure}[htbp]
  \centering
  \includegraphics[width=\linewidth]{assets/figure_p59_14.png}
\end{figure}

\begin{figure}[htbp]
  \centering
  \includegraphics[width=\linewidth]{assets/figure_p59_15.png}
\end{figure}

\begin{figure}[htbp]
  \centering
  \includegraphics[width=\linewidth]{assets/figure_p59_16.png}
\end{figure}

\begin{figure}[htbp]
  \centering
  \includegraphics[width=\linewidth]{assets/figure_p59_17.png}
\end{figure}

(3) 对偶规则

设\textbf{F}为任意逻辑表达式,若将F中所有运算符和常量作

如下变换：

$\textbf{\cdot + 0 1}$

\textbf{+ $\cdot 1 0}$

则所得新的逻辑表达式即为F的对偶式，记为\textbf{F’.}
例 有\textbf{F=A B + C D F’=(A+B)(C+D)}
例 有 \textbf{F=A+B+C+D+E F’=A B C D E}

\begin{figure}[htbp]
  \centering
  \includegraphics[width=\linewidth]{assets/figure_p60_1.png}
\end{figure}

\begin{figure}[htbp]
  \centering
  \includegraphics[width=\linewidth]{assets/figure_p60_2.png}
\end{figure}

\begin{figure}[htbp]
  \centering
  \includegraphics[width=\linewidth]{assets/figure_p60_3.png}
\end{figure}

\begin{figure}[htbp]
  \centering
  \includegraphics[width=\linewidth]{assets/figure_p60_4.png}
\end{figure}

\begin{figure}[htbp]
  \centering
  \includegraphics[width=\linewidth]{assets/figure_p60_5.png}
\end{figure}

\begin{figure}[htbp]
  \centering
  \includegraphics[width=\linewidth]{assets/figure_p60_6.png}
\end{figure}

\begin{figure}[htbp]
  \centering
  \includegraphics[width=\linewidth]{assets/figure_p60_7.png}
\end{figure}

\begin{figure}[htbp]
  \centering
  \includegraphics[width=\linewidth]{assets/figure_p60_8.png}
\end{figure}

\begin{figure}[htbp]
  \centering
  \includegraphics[width=\linewidth]{assets/figure_p60_9.png}
\end{figure}

\begin{figure}[htbp]
  \centering
  \includegraphics[width=\linewidth]{assets/figure_p60_10.png}
\end{figure}

\begin{figure}[htbp]
  \centering
  \includegraphics[width=\linewidth]{assets/figure_p60_11.png}
\end{figure}

\begin{figure}[htbp]
  \centering
  \includegraphics[width=\linewidth]{assets/figure_p60_12.png}
\end{figure}

\begin{figure}[htbp]
  \centering
  \includegraphics[width=\linewidth]{assets/figure_p60_13.png}
\end{figure}

\begin{figure}[htbp]
  \centering
  \includegraphics[width=\linewidth]{assets/figure_p60_14.png}
\end{figure}

\begin{figure}[htbp]
  \centering
  \includegraphics[width=\linewidth]{assets/figure_p60_15.png}
\end{figure}

\begin{figure}[htbp]
  \centering
  \includegraphics[width=\linewidth]{assets/figure_p60_16.png}
\end{figure}

\begin{figure}[htbp]
  \centering
  \includegraphics[width=\linewidth]{assets/figure_p60_17.png}
\end{figure}

对偶是相互的,\textbf{F}和\textbf{F’}互为对偶式.求对偶式注意：
1） 保持原式运算的优先次序；
2）原式中的长短“非”号不变；
3）单变量的对偶式为自己。
对偶规则：若有两个逻辑表达式\textbf{F}和\textbf{G}相等，则各自的对

偶式\textbf{F’}和\textbf{G’}也相等。
使用对偶规则可使得某些表达式的证明更加方便。
例 ：
对偶关系
已知 \textbf{A(B+C)=AB+AC A+BC=(A+B)(A+C)}

\begin{figure}[htbp]
  \centering
  \includegraphics[width=\linewidth]{assets/figure_p61_1.png}
\end{figure}

\begin{figure}[htbp]
  \centering
  \includegraphics[width=\linewidth]{assets/figure_p61_2.png}
\end{figure}

\begin{figure}[htbp]
  \centering
  \includegraphics[width=\linewidth]{assets/figure_p61_3.png}
\end{figure}

\begin{figure}[htbp]
  \centering
  \includegraphics[width=\linewidth]{assets/figure_p61_4.png}
\end{figure}

\begin{figure}[htbp]
  \centering
  \includegraphics[width=\linewidth]{assets/figure_p61_5.png}
\end{figure}

\begin{figure}[htbp]
  \centering
  \includegraphics[width=\linewidth]{assets/figure_p61_6.png}
\end{figure}

\begin{figure}[htbp]
  \centering
  \includegraphics[width=\linewidth]{assets/figure_p61_7.png}
\end{figure}

\begin{figure}[htbp]
  \centering
  \includegraphics[width=\linewidth]{assets/figure_p61_8.png}
\end{figure}

\begin{figure}[htbp]
  \centering
  \includegraphics[width=\linewidth]{assets/figure_p61_9.png}
\end{figure}

\begin{figure}[htbp]
  \centering
  \includegraphics[width=\linewidth]{assets/figure_p61_10.png}
\end{figure}

\begin{figure}[htbp]
  \centering
  \includegraphics[width=\linewidth]{assets/figure_p61_11.png}
\end{figure}

\begin{figure}[htbp]
  \centering
  \includegraphics[width=\linewidth]{assets/figure_p61_12.png}
\end{figure}

\begin{figure}[htbp]
  \centering
  \includegraphics[width=\linewidth]{assets/figure_p61_13.png}
\end{figure}

\begin{figure}[htbp]
  \centering
  \includegraphics[width=\linewidth]{assets/figure_p61_14.png}
\end{figure}

\begin{figure}[htbp]
  \centering
  \includegraphics[width=\linewidth]{assets/figure_p61_15.png}
\end{figure}

\begin{figure}[htbp]
  \centering
  \includegraphics[width=\linewidth]{assets/figure_p61_16.png}
\end{figure}

\begin{figure}[htbp]
  \centering
  \includegraphics[width=\linewidth]{assets/figure_p61_17.png}
\end{figure}

4.常用公式
1）消去律 \textbf{AB+AB=A}

证明：

对偶关系

\[ AB+AB=A(B+B)=A•1=A (A+B)(A+B)=A \]

2) 吸收律1 \textbf{A+AB=A}

证明：

对偶关系

\[ A+AB=A(1+B)=A•1=A A(A+B)=A \]

\begin{figure}[htbp]
  \centering
  \includegraphics[width=\linewidth]{assets/figure_p62_1.png}
\end{figure}

\begin{figure}[htbp]
  \centering
  \includegraphics[width=\linewidth]{assets/figure_p62_2.png}
\end{figure}

\begin{figure}[htbp]
  \centering
  \includegraphics[width=\linewidth]{assets/figure_p62_3.png}
\end{figure}

\begin{figure}[htbp]
  \centering
  \includegraphics[width=\linewidth]{assets/figure_p62_4.png}
\end{figure}

\begin{figure}[htbp]
  \centering
  \includegraphics[width=\linewidth]{assets/figure_p62_5.png}
\end{figure}

\begin{figure}[htbp]
  \centering
  \includegraphics[width=\linewidth]{assets/figure_p62_6.png}
\end{figure}

\begin{figure}[htbp]
  \centering
  \includegraphics[width=\linewidth]{assets/figure_p62_7.png}
\end{figure}

\begin{figure}[htbp]
  \centering
  \includegraphics[width=\linewidth]{assets/figure_p62_8.png}
\end{figure}

\begin{figure}[htbp]
  \centering
  \includegraphics[width=\linewidth]{assets/figure_p62_9.png}
\end{figure}

\begin{figure}[htbp]
  \centering
  \includegraphics[width=\linewidth]{assets/figure_p62_10.png}
\end{figure}

\begin{figure}[htbp]
  \centering
  \includegraphics[width=\linewidth]{assets/figure_p62_11.png}
\end{figure}

\begin{figure}[htbp]
  \centering
  \includegraphics[width=\linewidth]{assets/figure_p62_12.png}
\end{figure}

\begin{figure}[htbp]
  \centering
  \includegraphics[width=\linewidth]{assets/figure_p62_13.png}
\end{figure}

\begin{figure}[htbp]
  \centering
  \includegraphics[width=\linewidth]{assets/figure_p62_14.png}
\end{figure}

\begin{figure}[htbp]
  \centering
  \includegraphics[width=\linewidth]{assets/figure_p62_15.png}
\end{figure}

\begin{figure}[htbp]
  \centering
  \includegraphics[width=\linewidth]{assets/figure_p62_16.png}
\end{figure}

\begin{figure}[htbp]
  \centering
  \includegraphics[width=\linewidth]{assets/figure_p62_17.png}
\end{figure}

\[ 3) 吸收律2 A+AB=A+B \]

证明：

对偶关系

\textbf{A+AB=(A+A)(A+B)=1•(A+B) A(A+B)=AB}
\textbf{=A+B}

\[ 4）包含律 AB+AC+BC=AB+AC \]

证明：

\begin{figure}[htbp]
  \centering
  \includegraphics[width=\linewidth]{assets/figure_p63_1.png}
\end{figure}

\begin{figure}[htbp]
  \centering
  \includegraphics[width=\linewidth]{assets/figure_p63_2.png}
\end{figure}

\begin{figure}[htbp]
  \centering
  \includegraphics[width=\linewidth]{assets/figure_p63_3.png}
\end{figure}

\begin{figure}[htbp]
  \centering
  \includegraphics[width=\linewidth]{assets/figure_p63_4.png}
\end{figure}

\begin{figure}[htbp]
  \centering
  \includegraphics[width=\linewidth]{assets/figure_p63_5.png}
\end{figure}

\begin{figure}[htbp]
  \centering
  \includegraphics[width=\linewidth]{assets/figure_p63_6.png}
\end{figure}

\begin{figure}[htbp]
  \centering
  \includegraphics[width=\linewidth]{assets/figure_p63_7.png}
\end{figure}

\begin{figure}[htbp]
  \centering
  \includegraphics[width=\linewidth]{assets/figure_p63_8.png}
\end{figure}

\begin{figure}[htbp]
  \centering
  \includegraphics[width=\linewidth]{assets/figure_p63_9.png}
\end{figure}

\begin{figure}[htbp]
  \centering
  \includegraphics[width=\linewidth]{assets/figure_p63_10.png}
\end{figure}

\begin{figure}[htbp]
  \centering
  \includegraphics[width=\linewidth]{assets/figure_p63_11.png}
\end{figure}

\begin{figure}[htbp]
  \centering
  \includegraphics[width=\linewidth]{assets/figure_p63_12.png}
\end{figure}

\begin{figure}[htbp]
  \centering
  \includegraphics[width=\linewidth]{assets/figure_p63_13.png}
\end{figure}

\begin{figure}[htbp]
  \centering
  \includegraphics[width=\linewidth]{assets/figure_p63_14.png}
\end{figure}

\begin{figure}[htbp]
  \centering
  \includegraphics[width=\linewidth]{assets/figure_p63_15.png}
\end{figure}

\begin{figure}[htbp]
  \centering
  \includegraphics[width=\linewidth]{assets/figure_p63_16.png}
\end{figure}

\begin{figure}[htbp]
  \centering
  \includegraphics[width=\linewidth]{assets/figure_p63_17.png}
\end{figure}

\textbf{AB+AC+BC}
\textbf{=AB+AC+(A+A)BC}

\[ 对偶关系 (A+B)(A+C)(B+C) \]

\textbf{=AB+AC+ABC+ABC}
\textbf{=(A+B)(A+C)}
\textbf{=AB(1+C)+AC(1+B)}
\textbf{=AB+AC}
5) 关于异或和同或运算

对偶数个变量而言，
有 \textbf{A}\textsubscript{\textbf{1}}\textbf{A}\textsubscript{\textbf{2}}\textbf{...}  \textbf{A}\textsubscript{\textbf{n}}\textbf{=A}\textsubscript{\textbf{1}} \textbf{A}\textsubscript{\textbf{2}} \textbf{...} \textbf{A}\textsubscript{\textbf{n}}
对奇数个变量而言，
有 \textbf{A}\textsubscript{\textbf{1}}\textbf{A}\textsubscript{\textbf{2}}\textbf{...}  \textbf{A}\textsubscript{\textbf{n}}\textbf{=A}\textsubscript{\textbf{1}} \textbf{A}\textsubscript{\textbf{2}} \textbf{...} \textbf{A}\textsubscript{\textbf{n}}

\begin{figure}[htbp]
  \centering
  \includegraphics[width=\linewidth]{assets/figure_p64_1.png}
\end{figure}

\begin{figure}[htbp]
  \centering
  \includegraphics[width=\linewidth]{assets/figure_p64_2.png}
\end{figure}

\begin{figure}[htbp]
  \centering
  \includegraphics[width=\linewidth]{assets/figure_p64_3.png}
\end{figure}

\begin{figure}[htbp]
  \centering
  \includegraphics[width=\linewidth]{assets/figure_p64_4.png}
\end{figure}

\begin{figure}[htbp]
  \centering
  \includegraphics[width=\linewidth]{assets/figure_p64_5.png}
\end{figure}

\begin{figure}[htbp]
  \centering
  \includegraphics[width=\linewidth]{assets/figure_p64_6.png}
\end{figure}

\begin{figure}[htbp]
  \centering
  \includegraphics[width=\linewidth]{assets/figure_p64_7.png}
\end{figure}

\begin{figure}[htbp]
  \centering
  \includegraphics[width=\linewidth]{assets/figure_p64_8.png}
\end{figure}

\begin{figure}[htbp]
  \centering
  \includegraphics[width=\linewidth]{assets/figure_p64_9.png}
\end{figure}

\begin{figure}[htbp]
  \centering
  \includegraphics[width=\linewidth]{assets/figure_p64_10.png}
\end{figure}

\begin{figure}[htbp]
  \centering
  \includegraphics[width=\linewidth]{assets/figure_p64_11.png}
\end{figure}

\begin{figure}[htbp]
  \centering
  \includegraphics[width=\linewidth]{assets/figure_p64_12.png}
\end{figure}

\begin{figure}[htbp]
  \centering
  \includegraphics[width=\linewidth]{assets/figure_p64_13.png}
\end{figure}

\begin{figure}[htbp]
  \centering
  \includegraphics[width=\linewidth]{assets/figure_p64_14.png}
\end{figure}

\begin{figure}[htbp]
  \centering
  \includegraphics[width=\linewidth]{assets/figure_p64_15.png}
\end{figure}

\begin{figure}[htbp]
  \centering
  \includegraphics[width=\linewidth]{assets/figure_p64_16.png}
\end{figure}

\begin{figure}[htbp]
  \centering
  \includegraphics[width=\linewidth]{assets/figure_p64_17.png}
\end{figure}

异或和同或的其他性质:
\textbf{A} \textbf{0=A A} \textbf{1=A}
\textbf{A} \textbf{1=A A} \textbf{0 =A}
\textbf{A}\textbf{A=0 A}\textbf{A= 1}
\textbf{A} \textbf{(B}  \textbf{C)=(A} \textbf{B )}  \textbf{C A} \textbf{(B}  \textbf{C)=(A} \textbf{B)}  \textbf{C}
\textbf{A (B}  \textbf{C)=AB} \textbf{AC A+(B}  \textbf{C )=(A+B)}  \textbf{(A+C)}

利用异或门可实现数字信号的极性控制.

同或功能由异或门实现.

\begin{figure}[htbp]
  \centering
  \includegraphics[width=\linewidth]{assets/figure_p65_1.png}
\end{figure}

\begin{figure}[htbp]
  \centering
  \includegraphics[width=\linewidth]{assets/figure_p65_2.png}
\end{figure}

\begin{figure}[htbp]
  \centering
  \includegraphics[width=\linewidth]{assets/figure_p65_3.png}
\end{figure}

\begin{figure}[htbp]
  \centering
  \includegraphics[width=\linewidth]{assets/figure_p65_4.png}
\end{figure}

\begin{figure}[htbp]
  \centering
  \includegraphics[width=\linewidth]{assets/figure_p65_5.png}
\end{figure}

\begin{figure}[htbp]
  \centering
  \includegraphics[width=\linewidth]{assets/figure_p65_6.png}
\end{figure}

\begin{figure}[htbp]
  \centering
  \includegraphics[width=\linewidth]{assets/figure_p65_7.png}
\end{figure}

\begin{figure}[htbp]
  \centering
  \includegraphics[width=\linewidth]{assets/figure_p65_8.png}
\end{figure}

\begin{figure}[htbp]
  \centering
  \includegraphics[width=\linewidth]{assets/figure_p65_9.png}
\end{figure}

\begin{figure}[htbp]
  \centering
  \includegraphics[width=\linewidth]{assets/figure_p65_10.png}
\end{figure}

\begin{figure}[htbp]
  \centering
  \includegraphics[width=\linewidth]{assets/figure_p65_11.png}
\end{figure}

\begin{figure}[htbp]
  \centering
  \includegraphics[width=\linewidth]{assets/figure_p65_12.png}
\end{figure}

\begin{figure}[htbp]
  \centering
  \includegraphics[width=\linewidth]{assets/figure_p65_13.png}
\end{figure}

\begin{figure}[htbp]
  \centering
  \includegraphics[width=\linewidth]{assets/figure_p65_14.png}
\end{figure}

\begin{figure}[htbp]
  \centering
  \includegraphics[width=\linewidth]{assets/figure_p65_15.png}
\end{figure}

\begin{figure}[htbp]
  \centering
  \includegraphics[width=\linewidth]{assets/figure_p65_16.png}
\end{figure}

\begin{figure}[htbp]
  \centering
  \includegraphics[width=\linewidth]{assets/figure_p65_17.png}
\end{figure}

8.2.5 逻辑函数的标准形式
1. 函数的“与–或”式和“或–与”式
“与–或”式，指一个函数表达式中包含若干个与”
项，这些“与”项的“或”表示这个函数。

\[ 例： F(A,B,C,D)=A+BC+ABCD \]

\textbf{“}或–与”式，指一个函数表达式中包含若干个“或”

项，这些“或”项的“与”表示这个函数。

\[ 例 ： F(A,B,C,D)=(A+C+D)(B+D)(A+B+D) \]

\begin{figure}[htbp]
  \centering
  \includegraphics[width=\linewidth]{assets/figure_p66_1.png}
\end{figure}

\begin{figure}[htbp]
  \centering
  \includegraphics[width=\linewidth]{assets/figure_p66_2.png}
\end{figure}

\begin{figure}[htbp]
  \centering
  \includegraphics[width=\linewidth]{assets/figure_p66_3.png}
\end{figure}

\begin{figure}[htbp]
  \centering
  \includegraphics[width=\linewidth]{assets/figure_p66_4.png}
\end{figure}

\begin{figure}[htbp]
  \centering
  \includegraphics[width=\linewidth]{assets/figure_p66_5.png}
\end{figure}

\begin{figure}[htbp]
  \centering
  \includegraphics[width=\linewidth]{assets/figure_p66_6.png}
\end{figure}

\begin{figure}[htbp]
  \centering
  \includegraphics[width=\linewidth]{assets/figure_p66_7.png}
\end{figure}

\begin{figure}[htbp]
  \centering
  \includegraphics[width=\linewidth]{assets/figure_p66_8.png}
\end{figure}

\begin{figure}[htbp]
  \centering
  \includegraphics[width=\linewidth]{assets/figure_p66_9.png}
\end{figure}

\begin{figure}[htbp]
  \centering
  \includegraphics[width=\linewidth]{assets/figure_p66_10.png}
\end{figure}

\begin{figure}[htbp]
  \centering
  \includegraphics[width=\linewidth]{assets/figure_p66_11.png}
\end{figure}

\begin{figure}[htbp]
  \centering
  \includegraphics[width=\linewidth]{assets/figure_p66_12.png}
\end{figure}

\begin{figure}[htbp]
  \centering
  \includegraphics[width=\linewidth]{assets/figure_p66_13.png}
\end{figure}

\begin{figure}[htbp]
  \centering
  \includegraphics[width=\linewidth]{assets/figure_p66_14.png}
\end{figure}

\begin{figure}[htbp]
  \centering
  \includegraphics[width=\linewidth]{assets/figure_p66_15.png}
\end{figure}

\begin{figure}[htbp]
  \centering
  \includegraphics[width=\linewidth]{assets/figure_p66_16.png}
\end{figure}

\begin{figure}[htbp]
  \centering
  \includegraphics[width=\linewidth]{assets/figure_p66_17.png}
\end{figure}

2. 逻辑函数的两种标准形式

1）最小项的概念

1）最小项特点 最小项是“与”项。

① n个变量构成的每个最小项，一定是包含n个因子
的乘积项；

② 在各个最小项中，每个变量必须以原变量或反变
量形式作为因子出现一次，而且仅出现一次。

\begin{figure}[htbp]
  \centering
  \includegraphics[width=\linewidth]{assets/figure_p67_1.png}
\end{figure}

\begin{figure}[htbp]
  \centering
  \includegraphics[width=\linewidth]{assets/figure_p67_2.png}
\end{figure}

\begin{figure}[htbp]
  \centering
  \includegraphics[width=\linewidth]{assets/figure_p67_3.png}
\end{figure}

\begin{figure}[htbp]
  \centering
  \includegraphics[width=\linewidth]{assets/figure_p67_4.png}
\end{figure}

\begin{figure}[htbp]
  \centering
  \includegraphics[width=\linewidth]{assets/figure_p67_5.png}
\end{figure}

\begin{figure}[htbp]
  \centering
  \includegraphics[width=\linewidth]{assets/figure_p67_6.png}
\end{figure}

\begin{figure}[htbp]
  \centering
  \includegraphics[width=\linewidth]{assets/figure_p67_7.png}
\end{figure}

\begin{figure}[htbp]
  \centering
  \includegraphics[width=\linewidth]{assets/figure_p67_8.png}
\end{figure}

\begin{figure}[htbp]
  \centering
  \includegraphics[width=\linewidth]{assets/figure_p67_9.png}
\end{figure}

\begin{figure}[htbp]
  \centering
  \includegraphics[width=\linewidth]{assets/figure_p67_10.png}
\end{figure}

\begin{figure}[htbp]
  \centering
  \includegraphics[width=\linewidth]{assets/figure_p67_11.png}
\end{figure}

\begin{figure}[htbp]
  \centering
  \includegraphics[width=\linewidth]{assets/figure_p67_12.png}
\end{figure}

\begin{figure}[htbp]
  \centering
  \includegraphics[width=\linewidth]{assets/figure_p67_13.png}
\end{figure}

\begin{figure}[htbp]
  \centering
  \includegraphics[width=\linewidth]{assets/figure_p67_14.png}
\end{figure}

\begin{figure}[htbp]
  \centering
  \includegraphics[width=\linewidth]{assets/figure_p67_15.png}
\end{figure}

\begin{figure}[htbp]
  \centering
  \includegraphics[width=\linewidth]{assets/figure_p67_16.png}
\end{figure}

\begin{figure}[htbp]
  \centering
  \includegraphics[width=\linewidth]{assets/figure_p67_17.png}
\end{figure}

例 有A、B两变量的最小项共有四项(2\textsuperscript{2})：

\textbf{A B A B A B A B}

例 有A、B、C三变量的最小项共有八项(2\textsuperscript{3})：

\textbf{ABC}、\textbf{ABC}、\textbf{ABC}、\textbf{ABC}、\textbf{ABC}、\textbf{ABC}、\textbf{ABC}、\textbf{ABC}

（2） 最小项编号

任一个最小项用 \textbf{m}\textsubscript{\textbf{i}} 表示，\textbf{m}表示最小项，下标 \textbf{i}

为使该最小项为\textbf{1}的变量取值所对应的等效十进制数。

\begin{figure}[htbp]
  \centering
  \includegraphics[width=\linewidth]{assets/figure_p68_1.png}
\end{figure}

\begin{figure}[htbp]
  \centering
  \includegraphics[width=\linewidth]{assets/figure_p68_2.png}
\end{figure}

\begin{figure}[htbp]
  \centering
  \includegraphics[width=\linewidth]{assets/figure_p68_3.png}
\end{figure}

\begin{figure}[htbp]
  \centering
  \includegraphics[width=\linewidth]{assets/figure_p68_4.png}
\end{figure}

\begin{figure}[htbp]
  \centering
  \includegraphics[width=\linewidth]{assets/figure_p68_5.png}
\end{figure}

\begin{figure}[htbp]
  \centering
  \includegraphics[width=\linewidth]{assets/figure_p68_6.png}
\end{figure}

\begin{figure}[htbp]
  \centering
  \includegraphics[width=\linewidth]{assets/figure_p68_7.png}
\end{figure}

\begin{figure}[htbp]
  \centering
  \includegraphics[width=\linewidth]{assets/figure_p68_8.png}
\end{figure}

\begin{figure}[htbp]
  \centering
  \includegraphics[width=\linewidth]{assets/figure_p68_9.png}
\end{figure}

\begin{figure}[htbp]
  \centering
  \includegraphics[width=\linewidth]{assets/figure_p68_10.png}
\end{figure}

\begin{figure}[htbp]
  \centering
  \includegraphics[width=\linewidth]{assets/figure_p68_11.png}
\end{figure}

\begin{figure}[htbp]
  \centering
  \includegraphics[width=\linewidth]{assets/figure_p68_12.png}
\end{figure}

\begin{figure}[htbp]
  \centering
  \includegraphics[width=\linewidth]{assets/figure_p68_13.png}
\end{figure}

\begin{figure}[htbp]
  \centering
  \includegraphics[width=\linewidth]{assets/figure_p68_14.png}
\end{figure}

\begin{figure}[htbp]
  \centering
  \includegraphics[width=\linewidth]{assets/figure_p68_15.png}
\end{figure}

\begin{figure}[htbp]
  \centering
  \includegraphics[width=\linewidth]{assets/figure_p68_16.png}
\end{figure}

\begin{figure}[htbp]
  \centering
  \includegraphics[width=\linewidth]{assets/figure_p68_17.png}
\end{figure}

例 ：有最小项 \textbf{A B C},要使该最小项为1，\textbf{A}、\textbf{B}、\textbf{C}的取

值应为1、0、1，二进制数 101所等效的十进制数为5，所

以 \textbf{ABC = m}

\textbf{5}

(3) 最小项的性质

① 变量任取一组值，仅有一个最小项为1，其他最小项为
零；

② n变量的全体最小项之和为1；

\begin{figure}[htbp]
  \centering
  \includegraphics[width=\linewidth]{assets/figure_p69_1.png}
\end{figure}

\begin{figure}[htbp]
  \centering
  \includegraphics[width=\linewidth]{assets/figure_p69_2.png}
\end{figure}

\begin{figure}[htbp]
  \centering
  \includegraphics[width=\linewidth]{assets/figure_p69_3.png}
\end{figure}

\begin{figure}[htbp]
  \centering
  \includegraphics[width=\linewidth]{assets/figure_p69_4.png}
\end{figure}

\begin{figure}[htbp]
  \centering
  \includegraphics[width=\linewidth]{assets/figure_p69_5.png}
\end{figure}

\begin{figure}[htbp]
  \centering
  \includegraphics[width=\linewidth]{assets/figure_p69_6.png}
\end{figure}

\begin{figure}[htbp]
  \centering
  \includegraphics[width=\linewidth]{assets/figure_p69_7.png}
\end{figure}

\begin{figure}[htbp]
  \centering
  \includegraphics[width=\linewidth]{assets/figure_p69_8.png}
\end{figure}

\begin{figure}[htbp]
  \centering
  \includegraphics[width=\linewidth]{assets/figure_p69_9.png}
\end{figure}

\begin{figure}[htbp]
  \centering
  \includegraphics[width=\linewidth]{assets/figure_p69_10.png}
\end{figure}

\begin{figure}[htbp]
  \centering
  \includegraphics[width=\linewidth]{assets/figure_p69_11.png}
\end{figure}

\begin{figure}[htbp]
  \centering
  \includegraphics[width=\linewidth]{assets/figure_p69_12.png}
\end{figure}

\begin{figure}[htbp]
  \centering
  \includegraphics[width=\linewidth]{assets/figure_p69_13.png}
\end{figure}

\begin{figure}[htbp]
  \centering
  \includegraphics[width=\linewidth]{assets/figure_p69_14.png}
\end{figure}

\begin{figure}[htbp]
  \centering
  \includegraphics[width=\linewidth]{assets/figure_p69_15.png}
\end{figure}

\begin{figure}[htbp]
  \centering
  \includegraphics[width=\linewidth]{assets/figure_p69_16.png}
\end{figure}

\begin{figure}[htbp]
  \centering
  \includegraphics[width=\linewidth]{assets/figure_p69_17.png}
\end{figure}

③ 不同的最小项相与，结果为0；

④ 两最小项相邻，相邻最小项相“或”，可以合并成一
项，并可以消去一个变量因子。

相邻的概念： 两最小项如仅有一个变量因子不同，其他

变量均相同，则称这两个最小项相邻.

相邻最小项相“或”的情况：

例： \textbf{A B C+A B C =A B}

\begin{figure}[htbp]
  \centering
  \includegraphics[width=\linewidth]{assets/figure_p70_1.png}
\end{figure}

\begin{figure}[htbp]
  \centering
  \includegraphics[width=\linewidth]{assets/figure_p70_2.png}
\end{figure}

\begin{figure}[htbp]
  \centering
  \includegraphics[width=\linewidth]{assets/figure_p70_3.png}
\end{figure}

\begin{figure}[htbp]
  \centering
  \includegraphics[width=\linewidth]{assets/figure_p70_4.png}
\end{figure}

\begin{figure}[htbp]
  \centering
  \includegraphics[width=\linewidth]{assets/figure_p70_5.png}
\end{figure}

\begin{figure}[htbp]
  \centering
  \includegraphics[width=\linewidth]{assets/figure_p70_6.png}
\end{figure}

\begin{figure}[htbp]
  \centering
  \includegraphics[width=\linewidth]{assets/figure_p70_7.png}
\end{figure}

\begin{figure}[htbp]
  \centering
  \includegraphics[width=\linewidth]{assets/figure_p70_8.png}
\end{figure}

\begin{figure}[htbp]
  \centering
  \includegraphics[width=\linewidth]{assets/figure_p70_9.png}
\end{figure}

\begin{figure}[htbp]
  \centering
  \includegraphics[width=\linewidth]{assets/figure_p70_10.png}
\end{figure}

\begin{figure}[htbp]
  \centering
  \includegraphics[width=\linewidth]{assets/figure_p70_11.png}
\end{figure}

\begin{figure}[htbp]
  \centering
  \includegraphics[width=\linewidth]{assets/figure_p70_12.png}
\end{figure}

\begin{figure}[htbp]
  \centering
  \includegraphics[width=\linewidth]{assets/figure_p70_13.png}
\end{figure}

\begin{figure}[htbp]
  \centering
  \includegraphics[width=\linewidth]{assets/figure_p70_14.png}
\end{figure}

\begin{figure}[htbp]
  \centering
  \includegraphics[width=\linewidth]{assets/figure_p70_15.png}
\end{figure}

\begin{figure}[htbp]
  \centering
  \includegraphics[width=\linewidth]{assets/figure_p70_16.png}
\end{figure}

\begin{figure}[htbp]
  \centering
  \includegraphics[width=\linewidth]{assets/figure_p70_17.png}
\end{figure}

任一 n 变量的最小项，必定和其他 n 个不同最小
项相邻。

2）最大项的概念
（1）最大项特点 最大项是“或”项。
① n个变量构成的每个最大项，一定是包含n个因子的
“或”项；

② 在各个最大项中，每个变量必须以原变量或反变量
形式作为因子出现一次，而且仅出现一次。

\begin{figure}[htbp]
  \centering
  \includegraphics[width=\linewidth]{assets/figure_p71_1.png}
\end{figure}

\begin{figure}[htbp]
  \centering
  \includegraphics[width=\linewidth]{assets/figure_p71_2.png}
\end{figure}

\begin{figure}[htbp]
  \centering
  \includegraphics[width=\linewidth]{assets/figure_p71_3.png}
\end{figure}

\begin{figure}[htbp]
  \centering
  \includegraphics[width=\linewidth]{assets/figure_p71_4.png}
\end{figure}

\begin{figure}[htbp]
  \centering
  \includegraphics[width=\linewidth]{assets/figure_p71_5.png}
\end{figure}

\begin{figure}[htbp]
  \centering
  \includegraphics[width=\linewidth]{assets/figure_p71_6.png}
\end{figure}

\begin{figure}[htbp]
  \centering
  \includegraphics[width=\linewidth]{assets/figure_p71_7.png}
\end{figure}

\begin{figure}[htbp]
  \centering
  \includegraphics[width=\linewidth]{assets/figure_p71_8.png}
\end{figure}

\begin{figure}[htbp]
  \centering
  \includegraphics[width=\linewidth]{assets/figure_p71_9.png}
\end{figure}

\begin{figure}[htbp]
  \centering
  \includegraphics[width=\linewidth]{assets/figure_p71_10.png}
\end{figure}

\begin{figure}[htbp]
  \centering
  \includegraphics[width=\linewidth]{assets/figure_p71_11.png}
\end{figure}

\begin{figure}[htbp]
  \centering
  \includegraphics[width=\linewidth]{assets/figure_p71_12.png}
\end{figure}

\begin{figure}[htbp]
  \centering
  \includegraphics[width=\linewidth]{assets/figure_p71_13.png}
\end{figure}

\begin{figure}[htbp]
  \centering
  \includegraphics[width=\linewidth]{assets/figure_p71_14.png}
\end{figure}

\begin{figure}[htbp]
  \centering
  \includegraphics[width=\linewidth]{assets/figure_p71_15.png}
\end{figure}

\begin{figure}[htbp]
  \centering
  \includegraphics[width=\linewidth]{assets/figure_p71_16.png}
\end{figure}

\begin{figure}[htbp]
  \centering
  \includegraphics[width=\linewidth]{assets/figure_p71_17.png}
\end{figure}

例 有A、B两变量的最大项共有四项：
\textbf{A+ B A+ B A+ B A+ B}
例 有A、B、C三变量的最大项共有八项：

\textbf{A+B+C}、\textbf{A+B+C}、\textbf{A+B+C}、\textbf{A+B+C}、
\textbf{A+B+C}、\textbf{A+B+C}、\textbf{A+B+C}、\textbf{A+B+C}
(2) 最大项编号

任一个最大项用 \textbf{M}\textsubscript{\textbf{i}} 表示，\textbf{M}表示最大项，下标 \textbf{i}

为使该最大项为0的变量取值所对应的等效十进制数。

\begin{figure}[htbp]
  \centering
  \includegraphics[width=\linewidth]{assets/figure_p72_1.png}
\end{figure}

\begin{figure}[htbp]
  \centering
  \includegraphics[width=\linewidth]{assets/figure_p72_2.png}
\end{figure}

\begin{figure}[htbp]
  \centering
  \includegraphics[width=\linewidth]{assets/figure_p72_3.png}
\end{figure}

\begin{figure}[htbp]
  \centering
  \includegraphics[width=\linewidth]{assets/figure_p72_4.png}
\end{figure}

\begin{figure}[htbp]
  \centering
  \includegraphics[width=\linewidth]{assets/figure_p72_5.png}
\end{figure}

\begin{figure}[htbp]
  \centering
  \includegraphics[width=\linewidth]{assets/figure_p72_6.png}
\end{figure}

\begin{figure}[htbp]
  \centering
  \includegraphics[width=\linewidth]{assets/figure_p72_7.png}
\end{figure}

\begin{figure}[htbp]
  \centering
  \includegraphics[width=\linewidth]{assets/figure_p72_8.png}
\end{figure}

\begin{figure}[htbp]
  \centering
  \includegraphics[width=\linewidth]{assets/figure_p72_9.png}
\end{figure}

\begin{figure}[htbp]
  \centering
  \includegraphics[width=\linewidth]{assets/figure_p72_10.png}
\end{figure}

\begin{figure}[htbp]
  \centering
  \includegraphics[width=\linewidth]{assets/figure_p72_11.png}
\end{figure}

\begin{figure}[htbp]
  \centering
  \includegraphics[width=\linewidth]{assets/figure_p72_12.png}
\end{figure}

\begin{figure}[htbp]
  \centering
  \includegraphics[width=\linewidth]{assets/figure_p72_13.png}
\end{figure}

\begin{figure}[htbp]
  \centering
  \includegraphics[width=\linewidth]{assets/figure_p72_14.png}
\end{figure}

\begin{figure}[htbp]
  \centering
  \includegraphics[width=\linewidth]{assets/figure_p72_15.png}
\end{figure}

\begin{figure}[htbp]
  \centering
  \includegraphics[width=\linewidth]{assets/figure_p72_16.png}
\end{figure}

\begin{figure}[htbp]
  \centering
  \includegraphics[width=\linewidth]{assets/figure_p72_17.png}
\end{figure}

例 ：有最大项 \textbf{A +B+ C},要使该最大项为0，\textbf{A}、\textbf{B}、\textbf{C}

的取值应为1、0、0，二进制数 100所等效的十进制数为

\[ 4，所以 A+B+C =M \]

\textbf{4}

(3) 最大项的性质
① 变量任取一组值，仅有一个最大项为0，其它最大项
为1；
② n变量的全体最大项之积为0；

③ 不同的最大项相或，结果为 1；

\begin{figure}[htbp]
  \centering
  \includegraphics[width=\linewidth]{assets/figure_p73_1.png}
\end{figure}

\begin{figure}[htbp]
  \centering
  \includegraphics[width=\linewidth]{assets/figure_p73_2.png}
\end{figure}

\begin{figure}[htbp]
  \centering
  \includegraphics[width=\linewidth]{assets/figure_p73_3.png}
\end{figure}

\begin{figure}[htbp]
  \centering
  \includegraphics[width=\linewidth]{assets/figure_p73_4.png}
\end{figure}

\begin{figure}[htbp]
  \centering
  \includegraphics[width=\linewidth]{assets/figure_p73_5.png}
\end{figure}

\begin{figure}[htbp]
  \centering
  \includegraphics[width=\linewidth]{assets/figure_p73_6.png}
\end{figure}

\begin{figure}[htbp]
  \centering
  \includegraphics[width=\linewidth]{assets/figure_p73_7.png}
\end{figure}

\begin{figure}[htbp]
  \centering
  \includegraphics[width=\linewidth]{assets/figure_p73_8.png}
\end{figure}

\begin{figure}[htbp]
  \centering
  \includegraphics[width=\linewidth]{assets/figure_p73_9.png}
\end{figure}

\begin{figure}[htbp]
  \centering
  \includegraphics[width=\linewidth]{assets/figure_p73_10.png}
\end{figure}

\begin{figure}[htbp]
  \centering
  \includegraphics[width=\linewidth]{assets/figure_p73_11.png}
\end{figure}

\begin{figure}[htbp]
  \centering
  \includegraphics[width=\linewidth]{assets/figure_p73_12.png}
\end{figure}

\begin{figure}[htbp]
  \centering
  \includegraphics[width=\linewidth]{assets/figure_p73_13.png}
\end{figure}

\begin{figure}[htbp]
  \centering
  \includegraphics[width=\linewidth]{assets/figure_p73_14.png}
\end{figure}

\begin{figure}[htbp]
  \centering
  \includegraphics[width=\linewidth]{assets/figure_p73_15.png}
\end{figure}

\begin{figure}[htbp]
  \centering
  \includegraphics[width=\linewidth]{assets/figure_p73_16.png}
\end{figure}

\begin{figure}[htbp]
  \centering
  \includegraphics[width=\linewidth]{assets/figure_p73_17.png}
\end{figure}

④ 两相邻的最大项相“与”，可以合并成一项，并可以
消去一个变量因子。

相邻的概念：两最大项如仅有一个变量因子不同，其他
变量均相同，则称这两个最大项相邻。

相邻最大项相“与”的情况：

\[ 例： (A+B+C)(A+B+C)=A+B \]

任一 n 变量的最大项，必定和其他 n 个不同最大项
相邻。

\begin{figure}[htbp]
  \centering
  \includegraphics[width=\linewidth]{assets/figure_p74_1.png}
\end{figure}

\begin{figure}[htbp]
  \centering
  \includegraphics[width=\linewidth]{assets/figure_p74_2.png}
\end{figure}

\begin{figure}[htbp]
  \centering
  \includegraphics[width=\linewidth]{assets/figure_p74_3.png}
\end{figure}

\begin{figure}[htbp]
  \centering
  \includegraphics[width=\linewidth]{assets/figure_p74_4.png}
\end{figure}

\begin{figure}[htbp]
  \centering
  \includegraphics[width=\linewidth]{assets/figure_p74_5.png}
\end{figure}

\begin{figure}[htbp]
  \centering
  \includegraphics[width=\linewidth]{assets/figure_p74_6.png}
\end{figure}

\begin{figure}[htbp]
  \centering
  \includegraphics[width=\linewidth]{assets/figure_p74_7.png}
\end{figure}

\begin{figure}[htbp]
  \centering
  \includegraphics[width=\linewidth]{assets/figure_p74_8.png}
\end{figure}

\begin{figure}[htbp]
  \centering
  \includegraphics[width=\linewidth]{assets/figure_p74_9.png}
\end{figure}

\begin{figure}[htbp]
  \centering
  \includegraphics[width=\linewidth]{assets/figure_p74_10.png}
\end{figure}

\begin{figure}[htbp]
  \centering
  \includegraphics[width=\linewidth]{assets/figure_p74_11.png}
\end{figure}

\begin{figure}[htbp]
  \centering
  \includegraphics[width=\linewidth]{assets/figure_p74_12.png}
\end{figure}

\begin{figure}[htbp]
  \centering
  \includegraphics[width=\linewidth]{assets/figure_p74_13.png}
\end{figure}

\begin{figure}[htbp]
  \centering
  \includegraphics[width=\linewidth]{assets/figure_p74_14.png}
\end{figure}

\begin{figure}[htbp]
  \centering
  \includegraphics[width=\linewidth]{assets/figure_p74_15.png}
\end{figure}

\begin{figure}[htbp]
  \centering
  \includegraphics[width=\linewidth]{assets/figure_p74_16.png}
\end{figure}

\begin{figure}[htbp]
  \centering
  \includegraphics[width=\linewidth]{assets/figure_p74_17.png}
\end{figure}

3) 最小项和最大项的关系

编号下标相同的最小项和最大项互为反函数， 即

\textbf{M}\textsubscript{\textbf{i}}\textbf{= m}\textsubscript{\textbf{i}} 或 \textbf{m = M}

\textbf{i i}

4) 逻辑函数的最小项之和形式

最小项之和式为“与或”式，例：

\textbf{F(A,B,C) = ABC + ABC +ABC}
$\textbf{=\Sigma m(2 , 4 , 6)}$
$\textbf{=\Sigma (2 , 4 , 6)}$

\begin{figure}[htbp]
  \centering
  \includegraphics[width=\linewidth]{assets/figure_p75_1.png}
\end{figure}

\begin{figure}[htbp]
  \centering
  \includegraphics[width=\linewidth]{assets/figure_p75_2.png}
\end{figure}

\begin{figure}[htbp]
  \centering
  \includegraphics[width=\linewidth]{assets/figure_p75_3.png}
\end{figure}

\begin{figure}[htbp]
  \centering
  \includegraphics[width=\linewidth]{assets/figure_p75_4.png}
\end{figure}

\begin{figure}[htbp]
  \centering
  \includegraphics[width=\linewidth]{assets/figure_p75_5.png}
\end{figure}

\begin{figure}[htbp]
  \centering
  \includegraphics[width=\linewidth]{assets/figure_p75_6.png}
\end{figure}

\begin{figure}[htbp]
  \centering
  \includegraphics[width=\linewidth]{assets/figure_p75_7.png}
\end{figure}

\begin{figure}[htbp]
  \centering
  \includegraphics[width=\linewidth]{assets/figure_p75_8.png}
\end{figure}

\begin{figure}[htbp]
  \centering
  \includegraphics[width=\linewidth]{assets/figure_p75_9.png}
\end{figure}

\begin{figure}[htbp]
  \centering
  \includegraphics[width=\linewidth]{assets/figure_p75_10.png}
\end{figure}

\begin{figure}[htbp]
  \centering
  \includegraphics[width=\linewidth]{assets/figure_p75_11.png}
\end{figure}

\begin{figure}[htbp]
  \centering
  \includegraphics[width=\linewidth]{assets/figure_p75_12.png}
\end{figure}

\begin{figure}[htbp]
  \centering
  \includegraphics[width=\linewidth]{assets/figure_p75_13.png}
\end{figure}

\begin{figure}[htbp]
  \centering
  \includegraphics[width=\linewidth]{assets/figure_p75_14.png}
\end{figure}

\begin{figure}[htbp]
  \centering
  \includegraphics[width=\linewidth]{assets/figure_p75_15.png}
\end{figure}

\begin{figure}[htbp]
  \centering
  \includegraphics[width=\linewidth]{assets/figure_p75_16.png}
\end{figure}

\begin{figure}[htbp]
  \centering
  \includegraphics[width=\linewidth]{assets/figure_p75_17.png}
\end{figure}

任一逻辑函数都可以表达为最小项之和的形式,而且
是唯一的.

例 : \textbf{F(A,B,C) = A B +A C} 该式不是最小项之和形式
\textbf{=AB}（\textbf{C+C}）\textbf{+AC}（\textbf{B+B}）
\textbf{=ABC+ABC+ABC+ABC}
$\textbf{=\Sigma m}（\textbf{1}，\textbf{3}，\textbf{6}，\textbf{7}）$

5）逻辑函数的最大项之积的形式

\begin{figure}[htbp]
  \centering
  \includegraphics[width=\linewidth]{assets/figure_p76_1.png}
\end{figure}

\begin{figure}[htbp]
  \centering
  \includegraphics[width=\linewidth]{assets/figure_p76_2.png}
\end{figure}

\begin{figure}[htbp]
  \centering
  \includegraphics[width=\linewidth]{assets/figure_p76_3.png}
\end{figure}

\begin{figure}[htbp]
  \centering
  \includegraphics[width=\linewidth]{assets/figure_p76_4.png}
\end{figure}

\begin{figure}[htbp]
  \centering
  \includegraphics[width=\linewidth]{assets/figure_p76_5.png}
\end{figure}

\begin{figure}[htbp]
  \centering
  \includegraphics[width=\linewidth]{assets/figure_p76_6.png}
\end{figure}

\begin{figure}[htbp]
  \centering
  \includegraphics[width=\linewidth]{assets/figure_p76_7.png}
\end{figure}

\begin{figure}[htbp]
  \centering
  \includegraphics[width=\linewidth]{assets/figure_p76_8.png}
\end{figure}

\begin{figure}[htbp]
  \centering
  \includegraphics[width=\linewidth]{assets/figure_p76_9.png}
\end{figure}

\begin{figure}[htbp]
  \centering
  \includegraphics[width=\linewidth]{assets/figure_p76_10.png}
\end{figure}

\begin{figure}[htbp]
  \centering
  \includegraphics[width=\linewidth]{assets/figure_p76_11.png}
\end{figure}

\begin{figure}[htbp]
  \centering
  \includegraphics[width=\linewidth]{assets/figure_p76_12.png}
\end{figure}

\begin{figure}[htbp]
  \centering
  \includegraphics[width=\linewidth]{assets/figure_p76_13.png}
\end{figure}

\begin{figure}[htbp]
  \centering
  \includegraphics[width=\linewidth]{assets/figure_p76_14.png}
\end{figure}

\begin{figure}[htbp]
  \centering
  \includegraphics[width=\linewidth]{assets/figure_p76_15.png}
\end{figure}

\begin{figure}[htbp]
  \centering
  \includegraphics[width=\linewidth]{assets/figure_p76_16.png}
\end{figure}

\begin{figure}[htbp]
  \centering
  \includegraphics[width=\linewidth]{assets/figure_p76_17.png}
\end{figure}

逻辑函数的最大项之积的形式为“或与”式，

\[ 例： F(A,B,C) = (A+B+C)(A+B+C)(A+B+C) \]

$\textbf{=\Pi M (0 , 2 , 4 )}$
\textbf{= $\Pi (0 , 2 , 4 )}$

任一逻辑函数都可以表达为最大项之积的形式,而且
是唯一的.

\begin{figure}[htbp]
  \centering
  \includegraphics[width=\linewidth]{assets/figure_p77_1.png}
\end{figure}

\begin{figure}[htbp]
  \centering
  \includegraphics[width=\linewidth]{assets/figure_p77_2.png}
\end{figure}

\begin{figure}[htbp]
  \centering
  \includegraphics[width=\linewidth]{assets/figure_p77_3.png}
\end{figure}

\begin{figure}[htbp]
  \centering
  \includegraphics[width=\linewidth]{assets/figure_p77_4.png}
\end{figure}

\begin{figure}[htbp]
  \centering
  \includegraphics[width=\linewidth]{assets/figure_p77_5.png}
\end{figure}

\begin{figure}[htbp]
  \centering
  \includegraphics[width=\linewidth]{assets/figure_p77_6.png}
\end{figure}

\begin{figure}[htbp]
  \centering
  \includegraphics[width=\linewidth]{assets/figure_p77_7.png}
\end{figure}

\begin{figure}[htbp]
  \centering
  \includegraphics[width=\linewidth]{assets/figure_p77_8.png}
\end{figure}

\begin{figure}[htbp]
  \centering
  \includegraphics[width=\linewidth]{assets/figure_p77_9.png}
\end{figure}

\begin{figure}[htbp]
  \centering
  \includegraphics[width=\linewidth]{assets/figure_p77_10.png}
\end{figure}

\begin{figure}[htbp]
  \centering
  \includegraphics[width=\linewidth]{assets/figure_p77_11.png}
\end{figure}

\begin{figure}[htbp]
  \centering
  \includegraphics[width=\linewidth]{assets/figure_p77_12.png}
\end{figure}

\begin{figure}[htbp]
  \centering
  \includegraphics[width=\linewidth]{assets/figure_p77_13.png}
\end{figure}

\begin{figure}[htbp]
  \centering
  \includegraphics[width=\linewidth]{assets/figure_p77_14.png}
\end{figure}

\begin{figure}[htbp]
  \centering
  \includegraphics[width=\linewidth]{assets/figure_p77_15.png}
\end{figure}

\begin{figure}[htbp]
  \centering
  \includegraphics[width=\linewidth]{assets/figure_p77_16.png}
\end{figure}

\begin{figure}[htbp]
  \centering
  \includegraphics[width=\linewidth]{assets/figure_p77_17.png}
\end{figure}

\[ 例 : F(A,B,C) = (A + C )(B + C) \]

\textbf{=(A+B $\cdot$ B+C)(A $\cdot$ A+B+C)}
\textbf{=(A+B+C)(A+B+C)(A+B+C)(A+B+C)}
$\textbf{=\Pi M (1 , 4 , 5 , 6 )}$

6) 最小项之和的形式和最大项之积的形式之间的关系

\[ 若 F = \Sigma m i 则 F = \Sigma m j \]

\textbf{j}  \textbf{i}

\[ F = \Sigma m j =\Pi m = \Pi M \]

\textbf{j j}

\textbf{j}  \textbf{i}
\textbf{j}  \textbf{i j}  \textbf{i}

\begin{figure}[htbp]
  \centering
  \includegraphics[width=\linewidth]{assets/figure_p78_1.png}
\end{figure}

\begin{figure}[htbp]
  \centering
  \includegraphics[width=\linewidth]{assets/figure_p78_2.png}
\end{figure}

\begin{figure}[htbp]
  \centering
  \includegraphics[width=\linewidth]{assets/figure_p78_3.png}
\end{figure}

\begin{figure}[htbp]
  \centering
  \includegraphics[width=\linewidth]{assets/figure_p78_4.png}
\end{figure}

\begin{figure}[htbp]
  \centering
  \includegraphics[width=\linewidth]{assets/figure_p78_5.png}
\end{figure}

\begin{figure}[htbp]
  \centering
  \includegraphics[width=\linewidth]{assets/figure_p78_6.png}
\end{figure}

\begin{figure}[htbp]
  \centering
  \includegraphics[width=\linewidth]{assets/figure_p78_7.png}
\end{figure}

\begin{figure}[htbp]
  \centering
  \includegraphics[width=\linewidth]{assets/figure_p78_8.png}
\end{figure}

\begin{figure}[htbp]
  \centering
  \includegraphics[width=\linewidth]{assets/figure_p78_9.png}
\end{figure}

\begin{figure}[htbp]
  \centering
  \includegraphics[width=\linewidth]{assets/figure_p78_10.png}
\end{figure}

\begin{figure}[htbp]
  \centering
  \includegraphics[width=\linewidth]{assets/figure_p78_11.png}
\end{figure}

\begin{figure}[htbp]
  \centering
  \includegraphics[width=\linewidth]{assets/figure_p78_12.png}
\end{figure}

\begin{figure}[htbp]
  \centering
  \includegraphics[width=\linewidth]{assets/figure_p78_13.png}
\end{figure}

\begin{figure}[htbp]
  \centering
  \includegraphics[width=\linewidth]{assets/figure_p78_14.png}
\end{figure}

\begin{figure}[htbp]
  \centering
  \includegraphics[width=\linewidth]{assets/figure_p78_15.png}
\end{figure}

\begin{figure}[htbp]
  \centering
  \includegraphics[width=\linewidth]{assets/figure_p78_16.png}
\end{figure}

\begin{figure}[htbp]
  \centering
  \includegraphics[width=\linewidth]{assets/figure_p78_17.png}
\end{figure}

例 : \textbf{F (A , B , C) $= \Sigma (1 , 3 , 4 , 6 , 7)}$

$\textbf{=\Pi (0 , 2 , 5 )}$

3. 真值表与逻辑表达式

真值表与逻辑表达式都是表示逻辑函数的方法。
（1） 由逻辑函数式列真值表

由逻辑函数式列真值表可采用三种方法，以例说明：

例： 试列出下列逻辑函数式的真值表。

\textbf{F}（\textbf{A}，\textbf{B}，\textbf{C}）\textbf{=AB+BC}

\begin{figure}[htbp]
  \centering
  \includegraphics[width=\linewidth]{assets/figure_p79_1.png}
\end{figure}

\begin{figure}[htbp]
  \centering
  \includegraphics[width=\linewidth]{assets/figure_p79_2.png}
\end{figure}

\begin{figure}[htbp]
  \centering
  \includegraphics[width=\linewidth]{assets/figure_p79_3.png}
\end{figure}

\begin{figure}[htbp]
  \centering
  \includegraphics[width=\linewidth]{assets/figure_p79_4.png}
\end{figure}

\begin{figure}[htbp]
  \centering
  \includegraphics[width=\linewidth]{assets/figure_p79_5.png}
\end{figure}

\begin{figure}[htbp]
  \centering
  \includegraphics[width=\linewidth]{assets/figure_p79_6.png}
\end{figure}

\begin{figure}[htbp]
  \centering
  \includegraphics[width=\linewidth]{assets/figure_p79_7.png}
\end{figure}

\begin{figure}[htbp]
  \centering
  \includegraphics[width=\linewidth]{assets/figure_p79_8.png}
\end{figure}

\begin{figure}[htbp]
  \centering
  \includegraphics[width=\linewidth]{assets/figure_p79_9.png}
\end{figure}

\begin{figure}[htbp]
  \centering
  \includegraphics[width=\linewidth]{assets/figure_p79_10.png}
\end{figure}

\begin{figure}[htbp]
  \centering
  \includegraphics[width=\linewidth]{assets/figure_p79_11.png}
\end{figure}

\begin{figure}[htbp]
  \centering
  \includegraphics[width=\linewidth]{assets/figure_p79_12.png}
\end{figure}

\begin{figure}[htbp]
  \centering
  \includegraphics[width=\linewidth]{assets/figure_p79_13.png}
\end{figure}

\begin{figure}[htbp]
  \centering
  \includegraphics[width=\linewidth]{assets/figure_p79_14.png}
\end{figure}

\begin{figure}[htbp]
  \centering
  \includegraphics[width=\linewidth]{assets/figure_p79_15.png}
\end{figure}

\begin{figure}[htbp]
  \centering
  \includegraphics[width=\linewidth]{assets/figure_p79_16.png}
\end{figure}

\begin{figure}[htbp]
  \centering
  \includegraphics[width=\linewidth]{assets/figure_p79_17.png}
\end{figure}

方法一：将\textbf{A}、\textbf{B}、\textbf{C}三变量的所有取值的组合（共八

种），分别代入函数式，逐一算出函数值，填入

真值表中。
方法二：先将函数式\textbf{F}表示为最小项之和的形式：

\textbf{F(A,B,C) =AB+BC}
\textbf{=AB}（\textbf{C+C}）\textbf{+BC}（\textbf{A+A}）
\textbf{=ABC+ABC+ABC}
$\textbf{=\Sigma m}（\textbf{3}，\textbf{6}，\textbf{7}）$

\begin{figure}[htbp]
  \centering
  \includegraphics[width=\linewidth]{assets/figure_p80_1.png}
\end{figure}

\begin{figure}[htbp]
  \centering
  \includegraphics[width=\linewidth]{assets/figure_p80_2.png}
\end{figure}

\begin{figure}[htbp]
  \centering
  \includegraphics[width=\linewidth]{assets/figure_p80_3.png}
\end{figure}

\begin{figure}[htbp]
  \centering
  \includegraphics[width=\linewidth]{assets/figure_p80_4.png}
\end{figure}

\begin{figure}[htbp]
  \centering
  \includegraphics[width=\linewidth]{assets/figure_p80_5.png}
\end{figure}

\begin{figure}[htbp]
  \centering
  \includegraphics[width=\linewidth]{assets/figure_p80_6.png}
\end{figure}

\begin{figure}[htbp]
  \centering
  \includegraphics[width=\linewidth]{assets/figure_p80_7.png}
\end{figure}

\begin{figure}[htbp]
  \centering
  \includegraphics[width=\linewidth]{assets/figure_p80_8.png}
\end{figure}

\begin{figure}[htbp]
  \centering
  \includegraphics[width=\linewidth]{assets/figure_p80_9.png}
\end{figure}

\begin{figure}[htbp]
  \centering
  \includegraphics[width=\linewidth]{assets/figure_p80_10.png}
\end{figure}

\begin{figure}[htbp]
  \centering
  \includegraphics[width=\linewidth]{assets/figure_p80_11.png}
\end{figure}

\begin{figure}[htbp]
  \centering
  \includegraphics[width=\linewidth]{assets/figure_p80_12.png}
\end{figure}

\begin{figure}[htbp]
  \centering
  \includegraphics[width=\linewidth]{assets/figure_p80_13.png}
\end{figure}

\begin{figure}[htbp]
  \centering
  \includegraphics[width=\linewidth]{assets/figure_p80_14.png}
\end{figure}

\begin{figure}[htbp]
  \centering
  \includegraphics[width=\linewidth]{assets/figure_p80_15.png}
\end{figure}

\begin{figure}[htbp]
  \centering
  \includegraphics[width=\linewidth]{assets/figure_p80_16.png}
\end{figure}

\begin{figure}[htbp]
  \centering
  \includegraphics[width=\linewidth]{assets/figure_p80_17.png}
\end{figure}

最后根据最小项的性质，在真值表中对应于\textbf{ABC}取值为

011、110、111处填“1”，其它位置填“0”。

\textbf{A B C F}
\textbf{0 0 0 0}
\textbf{0 0 1 0}
\textbf{0 1 0 0}
\textbf{0 1 1 1}
\textbf{1 0 0 0}
\textbf{1 0 1 0}
\textbf{1 1 0 1}
\textbf{1 1 1 1}

\begin{figure}[htbp]
  \centering
  \includegraphics[width=\linewidth]{assets/figure_p81_1.png}
\end{figure}

\begin{figure}[htbp]
  \centering
  \includegraphics[width=\linewidth]{assets/figure_p81_2.png}
\end{figure}

\begin{figure}[htbp]
  \centering
  \includegraphics[width=\linewidth]{assets/figure_p81_3.png}
\end{figure}

\begin{figure}[htbp]
  \centering
  \includegraphics[width=\linewidth]{assets/figure_p81_4.png}
\end{figure}

\begin{figure}[htbp]
  \centering
  \includegraphics[width=\linewidth]{assets/figure_p81_5.png}
\end{figure}

\begin{figure}[htbp]
  \centering
  \includegraphics[width=\linewidth]{assets/figure_p81_6.png}
\end{figure}

\begin{figure}[htbp]
  \centering
  \includegraphics[width=\linewidth]{assets/figure_p81_7.png}
\end{figure}

\begin{figure}[htbp]
  \centering
  \includegraphics[width=\linewidth]{assets/figure_p81_8.png}
\end{figure}

\begin{figure}[htbp]
  \centering
  \includegraphics[width=\linewidth]{assets/figure_p81_9.png}
\end{figure}

\begin{figure}[htbp]
  \centering
  \includegraphics[width=\linewidth]{assets/figure_p81_10.png}
\end{figure}

\begin{figure}[htbp]
  \centering
  \includegraphics[width=\linewidth]{assets/figure_p81_11.png}
\end{figure}

\begin{figure}[htbp]
  \centering
  \includegraphics[width=\linewidth]{assets/figure_p81_12.png}
\end{figure}

\begin{figure}[htbp]
  \centering
  \includegraphics[width=\linewidth]{assets/figure_p81_13.png}
\end{figure}

\begin{figure}[htbp]
  \centering
  \includegraphics[width=\linewidth]{assets/figure_p81_14.png}
\end{figure}

\begin{figure}[htbp]
  \centering
  \includegraphics[width=\linewidth]{assets/figure_p81_15.png}
\end{figure}

\begin{figure}[htbp]
  \centering
  \includegraphics[width=\linewidth]{assets/figure_p81_16.png}
\end{figure}

\begin{figure}[htbp]
  \centering
  \includegraphics[width=\linewidth]{assets/figure_p81_17.png}
\end{figure}

方法三：根据函数式F的含义，直接填表。
函数\textbf{F=AB+BC}表示的含义为：

1）当\textbf{A}和\textbf{B}同时为“1”（即\textbf{AB=1}）时，\textbf{F=1}

2）当\textbf{B}和\textbf{C}同时为“1”（即\textbf{BC=1}）时，\textbf{F=1}

3）当不满足上面两种情况时，\textbf{F=0}

\begin{figure}[htbp]
  \centering
  \includegraphics[width=\linewidth]{assets/figure_p82_1.png}
\end{figure}

\begin{figure}[htbp]
  \centering
  \includegraphics[width=\linewidth]{assets/figure_p82_2.png}
\end{figure}

\begin{figure}[htbp]
  \centering
  \includegraphics[width=\linewidth]{assets/figure_p82_3.png}
\end{figure}

\begin{figure}[htbp]
  \centering
  \includegraphics[width=\linewidth]{assets/figure_p82_4.png}
\end{figure}

\begin{figure}[htbp]
  \centering
  \includegraphics[width=\linewidth]{assets/figure_p82_5.png}
\end{figure}

\begin{figure}[htbp]
  \centering
  \includegraphics[width=\linewidth]{assets/figure_p82_6.png}
\end{figure}

\begin{figure}[htbp]
  \centering
  \includegraphics[width=\linewidth]{assets/figure_p82_7.png}
\end{figure}

\begin{figure}[htbp]
  \centering
  \includegraphics[width=\linewidth]{assets/figure_p82_8.png}
\end{figure}

\begin{figure}[htbp]
  \centering
  \includegraphics[width=\linewidth]{assets/figure_p82_9.png}
\end{figure}

\begin{figure}[htbp]
  \centering
  \includegraphics[width=\linewidth]{assets/figure_p82_10.png}
\end{figure}

\begin{figure}[htbp]
  \centering
  \includegraphics[width=\linewidth]{assets/figure_p82_11.png}
\end{figure}

\begin{figure}[htbp]
  \centering
  \includegraphics[width=\linewidth]{assets/figure_p82_12.png}
\end{figure}

\begin{figure}[htbp]
  \centering
  \includegraphics[width=\linewidth]{assets/figure_p82_13.png}
\end{figure}

\begin{figure}[htbp]
  \centering
  \includegraphics[width=\linewidth]{assets/figure_p82_14.png}
\end{figure}

\begin{figure}[htbp]
  \centering
  \includegraphics[width=\linewidth]{assets/figure_p82_15.png}
\end{figure}

\begin{figure}[htbp]
  \centering
  \includegraphics[width=\linewidth]{assets/figure_p82_16.png}
\end{figure}

\begin{figure}[htbp]
  \centering
  \includegraphics[width=\linewidth]{assets/figure_p82_17.png}
\end{figure}

\textbf{A B C F}
\textbf{0 0 0 0}
\textbf{0 0 1 0}
\textbf{0 1 0 0} 方法三是一种较好的
\textbf{0 1 1 1} 方法，要熟练掌握。
\textbf{1 0 0 0}
\textbf{1 0 1 0}
\textbf{1 1 0 1}
\textbf{1 1 1 1}

\begin{figure}[htbp]
  \centering
  \includegraphics[width=\linewidth]{assets/figure_p83_1.png}
\end{figure}

\begin{figure}[htbp]
  \centering
  \includegraphics[width=\linewidth]{assets/figure_p83_2.png}
\end{figure}

\begin{figure}[htbp]
  \centering
  \includegraphics[width=\linewidth]{assets/figure_p83_3.png}
\end{figure}

\begin{figure}[htbp]
  \centering
  \includegraphics[width=\linewidth]{assets/figure_p83_4.png}
\end{figure}

\begin{figure}[htbp]
  \centering
  \includegraphics[width=\linewidth]{assets/figure_p83_5.png}
\end{figure}

\begin{figure}[htbp]
  \centering
  \includegraphics[width=\linewidth]{assets/figure_p83_6.png}
\end{figure}

\begin{figure}[htbp]
  \centering
  \includegraphics[width=\linewidth]{assets/figure_p83_7.png}
\end{figure}

\begin{figure}[htbp]
  \centering
  \includegraphics[width=\linewidth]{assets/figure_p83_8.png}
\end{figure}

\begin{figure}[htbp]
  \centering
  \includegraphics[width=\linewidth]{assets/figure_p83_9.png}
\end{figure}

\begin{figure}[htbp]
  \centering
  \includegraphics[width=\linewidth]{assets/figure_p83_10.png}
\end{figure}

\begin{figure}[htbp]
  \centering
  \includegraphics[width=\linewidth]{assets/figure_p83_11.png}
\end{figure}

\begin{figure}[htbp]
  \centering
  \includegraphics[width=\linewidth]{assets/figure_p83_12.png}
\end{figure}

\begin{figure}[htbp]
  \centering
  \includegraphics[width=\linewidth]{assets/figure_p83_13.png}
\end{figure}

\begin{figure}[htbp]
  \centering
  \includegraphics[width=\linewidth]{assets/figure_p83_14.png}
\end{figure}

\begin{figure}[htbp]
  \centering
  \includegraphics[width=\linewidth]{assets/figure_p83_15.png}
\end{figure}

\begin{figure}[htbp]
  \centering
  \includegraphics[width=\linewidth]{assets/figure_p83_16.png}
\end{figure}

\begin{figure}[htbp]
  \centering
  \includegraphics[width=\linewidth]{assets/figure_p83_17.png}
\end{figure}

\textbf{A B C F}\textsubscript{\textbf{1}} \textbf{F}\textsubscript{\textbf{2}} \textbf{F F}
例: \textbf{F=(A}\textbf{B) (B}\textbf{C) 0 0 0 0 0 0 1}
\textbf{0 0 1 0 1 0 1}
令: \textbf{F1=(A}\textbf{B) ; F2=(B}\textbf{C) 0 1 0 1 1 1 0}
\textbf{F=F1F2 0 1 1 1 0 0 1}
\textbf{1 0 0 1 0 0 1}
\textbf{1 0 1 1 1 1 0}
\textbf{1 1 0 0 1 0 1}
\textbf{1 1 1 0 0 0 1}

（2）由真值表写逻辑函数式

\begin{figure}[htbp]
  \centering
  \includegraphics[width=\linewidth]{assets/figure_p84_1.png}
\end{figure}

\begin{figure}[htbp]
  \centering
  \includegraphics[width=\linewidth]{assets/figure_p84_2.png}
\end{figure}

\begin{figure}[htbp]
  \centering
  \includegraphics[width=\linewidth]{assets/figure_p84_3.png}
\end{figure}

\begin{figure}[htbp]
  \centering
  \includegraphics[width=\linewidth]{assets/figure_p84_4.png}
\end{figure}

\begin{figure}[htbp]
  \centering
  \includegraphics[width=\linewidth]{assets/figure_p84_5.png}
\end{figure}

\begin{figure}[htbp]
  \centering
  \includegraphics[width=\linewidth]{assets/figure_p84_6.png}
\end{figure}

\begin{figure}[htbp]
  \centering
  \includegraphics[width=\linewidth]{assets/figure_p84_7.png}
\end{figure}

\begin{figure}[htbp]
  \centering
  \includegraphics[width=\linewidth]{assets/figure_p84_8.png}
\end{figure}

\begin{figure}[htbp]
  \centering
  \includegraphics[width=\linewidth]{assets/figure_p84_9.png}
\end{figure}

\begin{figure}[htbp]
  \centering
  \includegraphics[width=\linewidth]{assets/figure_p84_10.png}
\end{figure}

\begin{figure}[htbp]
  \centering
  \includegraphics[width=\linewidth]{assets/figure_p84_11.png}
\end{figure}

\begin{figure}[htbp]
  \centering
  \includegraphics[width=\linewidth]{assets/figure_p84_12.png}
\end{figure}

\begin{figure}[htbp]
  \centering
  \includegraphics[width=\linewidth]{assets/figure_p84_13.png}
\end{figure}

\begin{figure}[htbp]
  \centering
  \includegraphics[width=\linewidth]{assets/figure_p84_14.png}
\end{figure}

\begin{figure}[htbp]
  \centering
  \includegraphics[width=\linewidth]{assets/figure_p84_15.png}
\end{figure}

\begin{figure}[htbp]
  \centering
  \includegraphics[width=\linewidth]{assets/figure_p84_16.png}
\end{figure}

\begin{figure}[htbp]
  \centering
  \includegraphics[width=\linewidth]{assets/figure_p84_17.png}
\end{figure}

根据最小项的性质，用观察法，可直接从真值表写出

函数的最小项之和表达式。

例：已知函数\textbf{F}的真值表如下，求逻辑函数表达式。

\textbf{A B C F}
\textbf{0 0 0 0}
\textbf{0 0 1 0}
\textbf{0 1 0 0}
\textbf{0 1 1 1}
\textbf{1 0 0 0}
\textbf{1 0 1 1}
\textbf{1 1 0 1}
\textbf{1 1 1 1}

\begin{figure}[htbp]
  \centering
  \includegraphics[width=\linewidth]{assets/figure_p85_1.png}
\end{figure}

\begin{figure}[htbp]
  \centering
  \includegraphics[width=\linewidth]{assets/figure_p85_2.png}
\end{figure}

\begin{figure}[htbp]
  \centering
  \includegraphics[width=\linewidth]{assets/figure_p85_3.png}
\end{figure}

\begin{figure}[htbp]
  \centering
  \includegraphics[width=\linewidth]{assets/figure_p85_4.png}
\end{figure}

\begin{figure}[htbp]
  \centering
  \includegraphics[width=\linewidth]{assets/figure_p85_5.png}
\end{figure}

\begin{figure}[htbp]
  \centering
  \includegraphics[width=\linewidth]{assets/figure_p85_6.png}
\end{figure}

\begin{figure}[htbp]
  \centering
  \includegraphics[width=\linewidth]{assets/figure_p85_7.png}
\end{figure}

\begin{figure}[htbp]
  \centering
  \includegraphics[width=\linewidth]{assets/figure_p85_8.png}
\end{figure}

\begin{figure}[htbp]
  \centering
  \includegraphics[width=\linewidth]{assets/figure_p85_9.png}
\end{figure}

\begin{figure}[htbp]
  \centering
  \includegraphics[width=\linewidth]{assets/figure_p85_10.png}
\end{figure}

\begin{figure}[htbp]
  \centering
  \includegraphics[width=\linewidth]{assets/figure_p85_11.png}
\end{figure}

\begin{figure}[htbp]
  \centering
  \includegraphics[width=\linewidth]{assets/figure_p85_12.png}
\end{figure}

\begin{figure}[htbp]
  \centering
  \includegraphics[width=\linewidth]{assets/figure_p85_13.png}
\end{figure}

\begin{figure}[htbp]
  \centering
  \includegraphics[width=\linewidth]{assets/figure_p85_14.png}
\end{figure}

\begin{figure}[htbp]
  \centering
  \includegraphics[width=\linewidth]{assets/figure_p85_15.png}
\end{figure}

\begin{figure}[htbp]
  \centering
  \includegraphics[width=\linewidth]{assets/figure_p85_16.png}
\end{figure}

\begin{figure}[htbp]
  \centering
  \includegraphics[width=\linewidth]{assets/figure_p85_17.png}
\end{figure}

\textbf{A B C F} 解：由真值表可见，当
\textbf{0 0 0 0}

ABC取011、101、

\textbf{0 0 1 0}

110、111时，F为

\textbf{0 1 0 0}
\textbf{0 1 1 1} “1”。
\textbf{1 0 0 0}
\textbf{1 0 1 1}
\textbf{1 1 0 1}
\textbf{1 1 1 1}

所以，\textbf{F}由4个最小项组成：

$\textbf{F}（\textbf{A}，\textbf{B}，\textbf{C}）\textbf{=\Sigma m}（\textbf{3}，\textbf{5,$ 6}，\textbf{7}）
\textbf{=ABC+ABC+ABC+ABC}

\begin{figure}[htbp]
  \centering
  \includegraphics[width=\linewidth]{assets/figure_p86_1.png}
\end{figure}

\begin{figure}[htbp]
  \centering
  \includegraphics[width=\linewidth]{assets/figure_p86_2.png}
\end{figure}

\begin{figure}[htbp]
  \centering
  \includegraphics[width=\linewidth]{assets/figure_p86_3.png}
\end{figure}

\begin{figure}[htbp]
  \centering
  \includegraphics[width=\linewidth]{assets/figure_p86_4.png}
\end{figure}

\begin{figure}[htbp]
  \centering
  \includegraphics[width=\linewidth]{assets/figure_p86_5.png}
\end{figure}

\begin{figure}[htbp]
  \centering
  \includegraphics[width=\linewidth]{assets/figure_p86_6.png}
\end{figure}

\begin{figure}[htbp]
  \centering
  \includegraphics[width=\linewidth]{assets/figure_p86_7.png}
\end{figure}

\begin{figure}[htbp]
  \centering
  \includegraphics[width=\linewidth]{assets/figure_p86_8.png}
\end{figure}

\begin{figure}[htbp]
  \centering
  \includegraphics[width=\linewidth]{assets/figure_p86_9.png}
\end{figure}

\begin{figure}[htbp]
  \centering
  \includegraphics[width=\linewidth]{assets/figure_p86_10.png}
\end{figure}

\begin{figure}[htbp]
  \centering
  \includegraphics[width=\linewidth]{assets/figure_p86_11.png}
\end{figure}

\begin{figure}[htbp]
  \centering
  \includegraphics[width=\linewidth]{assets/figure_p86_12.png}
\end{figure}

\begin{figure}[htbp]
  \centering
  \includegraphics[width=\linewidth]{assets/figure_p86_13.png}
\end{figure}

\begin{figure}[htbp]
  \centering
  \includegraphics[width=\linewidth]{assets/figure_p86_14.png}
\end{figure}

\begin{figure}[htbp]
  \centering
  \includegraphics[width=\linewidth]{assets/figure_p86_15.png}
\end{figure}

\begin{figure}[htbp]
  \centering
  \includegraphics[width=\linewidth]{assets/figure_p86_16.png}
\end{figure}

\begin{figure}[htbp]
  \centering
  \includegraphics[width=\linewidth]{assets/figure_p86_17.png}
\end{figure}

8.2.6 逻辑函数的化简

化简的意义: ①节省元器件,降低电路成本;
② 提高电路可靠性;

③ 减少连线,制作方便.

逻辑函数的几种常用表达式:

\begin{figure}[htbp]
  \centering
  \includegraphics[width=\linewidth]{assets/figure_p87_1.png}
\end{figure}

\begin{figure}[htbp]
  \centering
  \includegraphics[width=\linewidth]{assets/figure_p87_2.png}
\end{figure}

\begin{figure}[htbp]
  \centering
  \includegraphics[width=\linewidth]{assets/figure_p87_3.png}
\end{figure}

\begin{figure}[htbp]
  \centering
  \includegraphics[width=\linewidth]{assets/figure_p87_4.png}
\end{figure}

\begin{figure}[htbp]
  \centering
  \includegraphics[width=\linewidth]{assets/figure_p87_5.png}
\end{figure}

\begin{figure}[htbp]
  \centering
  \includegraphics[width=\linewidth]{assets/figure_p87_6.png}
\end{figure}

\begin{figure}[htbp]
  \centering
  \includegraphics[width=\linewidth]{assets/figure_p87_7.png}
\end{figure}

\begin{figure}[htbp]
  \centering
  \includegraphics[width=\linewidth]{assets/figure_p87_8.png}
\end{figure}

\begin{figure}[htbp]
  \centering
  \includegraphics[width=\linewidth]{assets/figure_p87_9.png}
\end{figure}

\begin{figure}[htbp]
  \centering
  \includegraphics[width=\linewidth]{assets/figure_p87_10.png}
\end{figure}

\begin{figure}[htbp]
  \centering
  \includegraphics[width=\linewidth]{assets/figure_p87_11.png}
\end{figure}

\begin{figure}[htbp]
  \centering
  \includegraphics[width=\linewidth]{assets/figure_p87_12.png}
\end{figure}

\begin{figure}[htbp]
  \centering
  \includegraphics[width=\linewidth]{assets/figure_p87_13.png}
\end{figure}

\begin{figure}[htbp]
  \centering
  \includegraphics[width=\linewidth]{assets/figure_p87_14.png}
\end{figure}

\begin{figure}[htbp]
  \centering
  \includegraphics[width=\linewidth]{assets/figure_p87_15.png}
\end{figure}

\begin{figure}[htbp]
  \centering
  \includegraphics[width=\linewidth]{assets/figure_p87_16.png}
\end{figure}

\begin{figure}[htbp]
  \centering
  \includegraphics[width=\linewidth]{assets/figure_p87_17.png}
\end{figure}

\textbf{F(A,B,C) =AB+AC} 与或式
\textbf{=(A+C)(A+B)} 或与式 （两次求对偶）
$\textbf{=AB\cdot AC}$ 与非－与非式（两次求反）
\textbf{=A+C+A+B} 或非－或非式
\textbf{=AB+AC} 与或非式

最简与或表达式的标准：

1） 所得与或表达式中，乘积项（与项）数目最少；

2） 每个乘积项中所含的变量数最少。

\begin{figure}[htbp]
  \centering
  \includegraphics[width=\linewidth]{assets/figure_p88_1.png}
\end{figure}

\begin{figure}[htbp]
  \centering
  \includegraphics[width=\linewidth]{assets/figure_p88_2.png}
\end{figure}

\begin{figure}[htbp]
  \centering
  \includegraphics[width=\linewidth]{assets/figure_p88_3.png}
\end{figure}

\begin{figure}[htbp]
  \centering
  \includegraphics[width=\linewidth]{assets/figure_p88_4.png}
\end{figure}

\begin{figure}[htbp]
  \centering
  \includegraphics[width=\linewidth]{assets/figure_p88_5.png}
\end{figure}

\begin{figure}[htbp]
  \centering
  \includegraphics[width=\linewidth]{assets/figure_p88_6.png}
\end{figure}

\begin{figure}[htbp]
  \centering
  \includegraphics[width=\linewidth]{assets/figure_p88_7.png}
\end{figure}

\begin{figure}[htbp]
  \centering
  \includegraphics[width=\linewidth]{assets/figure_p88_8.png}
\end{figure}

\begin{figure}[htbp]
  \centering
  \includegraphics[width=\linewidth]{assets/figure_p88_9.png}
\end{figure}

\begin{figure}[htbp]
  \centering
  \includegraphics[width=\linewidth]{assets/figure_p88_10.png}
\end{figure}

\begin{figure}[htbp]
  \centering
  \includegraphics[width=\linewidth]{assets/figure_p88_11.png}
\end{figure}

\begin{figure}[htbp]
  \centering
  \includegraphics[width=\linewidth]{assets/figure_p88_12.png}
\end{figure}

\begin{figure}[htbp]
  \centering
  \includegraphics[width=\linewidth]{assets/figure_p88_13.png}
\end{figure}

\begin{figure}[htbp]
  \centering
  \includegraphics[width=\linewidth]{assets/figure_p88_14.png}
\end{figure}

\begin{figure}[htbp]
  \centering
  \includegraphics[width=\linewidth]{assets/figure_p88_15.png}
\end{figure}

\begin{figure}[htbp]
  \centering
  \includegraphics[width=\linewidth]{assets/figure_p88_16.png}
\end{figure}

\begin{figure}[htbp]
  \centering
  \includegraphics[width=\linewidth]{assets/figure_p88_17.png}
\end{figure}

逻辑函数常用的化简方法有： 公式法、卡诺图法和列

表法。本课程要求掌握公式法和卡诺图法。

1. 公式化简法

针对某一逻辑式,反复运用逻辑代数公式消去多余的乘

积项和每个乘积项中多余的因子,使函数式符合最简标准.

化简中常用方法:

\begin{figure}[htbp]
  \centering
  \includegraphics[width=\linewidth]{assets/figure_p89_1.png}
\end{figure}

\begin{figure}[htbp]
  \centering
  \includegraphics[width=\linewidth]{assets/figure_p89_2.png}
\end{figure}

\begin{figure}[htbp]
  \centering
  \includegraphics[width=\linewidth]{assets/figure_p89_3.png}
\end{figure}

\begin{figure}[htbp]
  \centering
  \includegraphics[width=\linewidth]{assets/figure_p89_4.png}
\end{figure}

\begin{figure}[htbp]
  \centering
  \includegraphics[width=\linewidth]{assets/figure_p89_5.png}
\end{figure}

\begin{figure}[htbp]
  \centering
  \includegraphics[width=\linewidth]{assets/figure_p89_6.png}
\end{figure}

\begin{figure}[htbp]
  \centering
  \includegraphics[width=\linewidth]{assets/figure_p89_7.png}
\end{figure}

\begin{figure}[htbp]
  \centering
  \includegraphics[width=\linewidth]{assets/figure_p89_8.png}
\end{figure}

\begin{figure}[htbp]
  \centering
  \includegraphics[width=\linewidth]{assets/figure_p89_9.png}
\end{figure}

\begin{figure}[htbp]
  \centering
  \includegraphics[width=\linewidth]{assets/figure_p89_10.png}
\end{figure}

\begin{figure}[htbp]
  \centering
  \includegraphics[width=\linewidth]{assets/figure_p89_11.png}
\end{figure}

\begin{figure}[htbp]
  \centering
  \includegraphics[width=\linewidth]{assets/figure_p89_12.png}
\end{figure}

\begin{figure}[htbp]
  \centering
  \includegraphics[width=\linewidth]{assets/figure_p89_13.png}
\end{figure}

\begin{figure}[htbp]
  \centering
  \includegraphics[width=\linewidth]{assets/figure_p89_14.png}
\end{figure}

\begin{figure}[htbp]
  \centering
  \includegraphics[width=\linewidth]{assets/figure_p89_15.png}
\end{figure}

\begin{figure}[htbp]
  \centering
  \includegraphics[width=\linewidth]{assets/figure_p89_16.png}
\end{figure}

\begin{figure}[htbp]
  \centering
  \includegraphics[width=\linewidth]{assets/figure_p89_17.png}
\end{figure}

(1) 并项法 利用公式 \textbf{AB+AB=A}
例: \textbf{F=ABC+ABC+ABC+ABC}

\[ 在化简中 = （ AB+AB ） C+ （ AB+AB ） C \]

注意

代入规则 \textbf{=}（\textbf{A}\textbf{B}）\textbf{C+}（\textbf{A}\textbf{B}）\textbf{C}
的使用 \textbf{=}（\textbf{A} \textbf{B}）\textbf{C+}（\textbf{A} \textbf{B}）\textbf{C=C}

(2)吸收法 利用公式 \textbf{A+AB=A}

\[ 例： F=A+ABC B+AC+D+BC \]

\[ =(A+BC)+(A+BC)B+AC+D \]

反演律

\textbf{=A+BC}

\begin{figure}[htbp]
  \centering
  \includegraphics[width=\linewidth]{assets/figure_p90_1.png}
\end{figure}

\begin{figure}[htbp]
  \centering
  \includegraphics[width=\linewidth]{assets/figure_p90_2.png}
\end{figure}

\begin{figure}[htbp]
  \centering
  \includegraphics[width=\linewidth]{assets/figure_p90_3.png}
\end{figure}

\begin{figure}[htbp]
  \centering
  \includegraphics[width=\linewidth]{assets/figure_p90_4.png}
\end{figure}

\begin{figure}[htbp]
  \centering
  \includegraphics[width=\linewidth]{assets/figure_p90_5.png}
\end{figure}

\begin{figure}[htbp]
  \centering
  \includegraphics[width=\linewidth]{assets/figure_p90_6.png}
\end{figure}

\begin{figure}[htbp]
  \centering
  \includegraphics[width=\linewidth]{assets/figure_p90_7.png}
\end{figure}

\begin{figure}[htbp]
  \centering
  \includegraphics[width=\linewidth]{assets/figure_p90_8.png}
\end{figure}

\begin{figure}[htbp]
  \centering
  \includegraphics[width=\linewidth]{assets/figure_p90_9.png}
\end{figure}

\begin{figure}[htbp]
  \centering
  \includegraphics[width=\linewidth]{assets/figure_p90_10.png}
\end{figure}

\begin{figure}[htbp]
  \centering
  \includegraphics[width=\linewidth]{assets/figure_p90_11.png}
\end{figure}

\begin{figure}[htbp]
  \centering
  \includegraphics[width=\linewidth]{assets/figure_p90_12.png}
\end{figure}

\begin{figure}[htbp]
  \centering
  \includegraphics[width=\linewidth]{assets/figure_p90_13.png}
\end{figure}

\begin{figure}[htbp]
  \centering
  \includegraphics[width=\linewidth]{assets/figure_p90_14.png}
\end{figure}

\begin{figure}[htbp]
  \centering
  \includegraphics[width=\linewidth]{assets/figure_p90_15.png}
\end{figure}

\begin{figure}[htbp]
  \centering
  \includegraphics[width=\linewidth]{assets/figure_p90_16.png}
\end{figure}

\begin{figure}[htbp]
  \centering
  \includegraphics[width=\linewidth]{assets/figure_p90_17.png}
\end{figure}

\[ (3) 消项法 利用公式 AB+AC+BC=AB+AC \]

例 ： \textbf{F=ABCD+AE+BE+CDE}
\textbf{=ABCD+(A+B)E+CDE}
\textbf{=ABCD+ABE+CDE}
\textbf{=ABCD+(A+B)E}
\textbf{=ABCD+AE+BE}

\[ (4) 消因子法 利用公式 A+AB=A+B \]

\begin{figure}[htbp]
  \centering
  \includegraphics[width=\linewidth]{assets/figure_p91_1.png}
\end{figure}

\begin{figure}[htbp]
  \centering
  \includegraphics[width=\linewidth]{assets/figure_p91_2.png}
\end{figure}

\begin{figure}[htbp]
  \centering
  \includegraphics[width=\linewidth]{assets/figure_p91_3.png}
\end{figure}

\begin{figure}[htbp]
  \centering
  \includegraphics[width=\linewidth]{assets/figure_p91_4.png}
\end{figure}

\begin{figure}[htbp]
  \centering
  \includegraphics[width=\linewidth]{assets/figure_p91_5.png}
\end{figure}

\begin{figure}[htbp]
  \centering
  \includegraphics[width=\linewidth]{assets/figure_p91_6.png}
\end{figure}

\begin{figure}[htbp]
  \centering
  \includegraphics[width=\linewidth]{assets/figure_p91_7.png}
\end{figure}

\begin{figure}[htbp]
  \centering
  \includegraphics[width=\linewidth]{assets/figure_p91_8.png}
\end{figure}

\begin{figure}[htbp]
  \centering
  \includegraphics[width=\linewidth]{assets/figure_p91_9.png}
\end{figure}

\begin{figure}[htbp]
  \centering
  \includegraphics[width=\linewidth]{assets/figure_p91_10.png}
\end{figure}

\begin{figure}[htbp]
  \centering
  \includegraphics[width=\linewidth]{assets/figure_p91_11.png}
\end{figure}

\begin{figure}[htbp]
  \centering
  \includegraphics[width=\linewidth]{assets/figure_p91_12.png}
\end{figure}

\begin{figure}[htbp]
  \centering
  \includegraphics[width=\linewidth]{assets/figure_p91_13.png}
\end{figure}

\begin{figure}[htbp]
  \centering
  \includegraphics[width=\linewidth]{assets/figure_p91_14.png}
\end{figure}

\begin{figure}[htbp]
  \centering
  \includegraphics[width=\linewidth]{assets/figure_p91_15.png}
\end{figure}

\begin{figure}[htbp]
  \centering
  \includegraphics[width=\linewidth]{assets/figure_p91_16.png}
\end{figure}

\begin{figure}[htbp]
  \centering
  \includegraphics[width=\linewidth]{assets/figure_p91_17.png}
\end{figure}

\[ 例： F=AB+AC+BC \]

\textbf{=AB+}（\textbf{A+B}）\textbf{C}
\textbf{=AB+ABC}
\textbf{=AB+C}
(5) 配项法 利用公式 \textbf{A+A=1} ；\textbf{A • 1=A}等

\[ 例： F=AB+AC+BC \]

\textbf{=AB+AC+(A+A)BC}
\textbf{=AB+AC+ABC+ABC}
\textbf{=(AB+ABC)+(AC+ABC)}
\textbf{=AB+AC}

\begin{figure}[htbp]
  \centering
  \includegraphics[width=\linewidth]{assets/figure_p92_1.png}
\end{figure}

\begin{figure}[htbp]
  \centering
  \includegraphics[width=\linewidth]{assets/figure_p92_2.png}
\end{figure}

\begin{figure}[htbp]
  \centering
  \includegraphics[width=\linewidth]{assets/figure_p92_3.png}
\end{figure}

\begin{figure}[htbp]
  \centering
  \includegraphics[width=\linewidth]{assets/figure_p92_4.png}
\end{figure}

\begin{figure}[htbp]
  \centering
  \includegraphics[width=\linewidth]{assets/figure_p92_5.png}
\end{figure}

\begin{figure}[htbp]
  \centering
  \includegraphics[width=\linewidth]{assets/figure_p92_6.png}
\end{figure}

\begin{figure}[htbp]
  \centering
  \includegraphics[width=\linewidth]{assets/figure_p92_7.png}
\end{figure}

\begin{figure}[htbp]
  \centering
  \includegraphics[width=\linewidth]{assets/figure_p92_8.png}
\end{figure}

\begin{figure}[htbp]
  \centering
  \includegraphics[width=\linewidth]{assets/figure_p92_9.png}
\end{figure}

\begin{figure}[htbp]
  \centering
  \includegraphics[width=\linewidth]{assets/figure_p92_10.png}
\end{figure}

\begin{figure}[htbp]
  \centering
  \includegraphics[width=\linewidth]{assets/figure_p92_11.png}
\end{figure}

\begin{figure}[htbp]
  \centering
  \includegraphics[width=\linewidth]{assets/figure_p92_12.png}
\end{figure}

\begin{figure}[htbp]
  \centering
  \includegraphics[width=\linewidth]{assets/figure_p92_13.png}
\end{figure}

\begin{figure}[htbp]
  \centering
  \includegraphics[width=\linewidth]{assets/figure_p92_14.png}
\end{figure}

\begin{figure}[htbp]
  \centering
  \includegraphics[width=\linewidth]{assets/figure_p92_15.png}
\end{figure}

\begin{figure}[htbp]
  \centering
  \includegraphics[width=\linewidth]{assets/figure_p92_16.png}
\end{figure}

\begin{figure}[htbp]
  \centering
  \includegraphics[width=\linewidth]{assets/figure_p92_17.png}
\end{figure}

对比较复杂的函数式，要求熟练掌握上述方法，才能
把函数化成最简。
2. 卡诺图化简法
该方法是将逻辑函数用一种称为“卡诺图”的图形来

表示,然后在卡诺图上进行函数的化简的方法.

1)卡诺图的构成

\begin{figure}[htbp]
  \centering
  \includegraphics[width=\linewidth]{assets/figure_p93_1.png}
\end{figure}

\begin{figure}[htbp]
  \centering
  \includegraphics[width=\linewidth]{assets/figure_p93_2.png}
\end{figure}

\begin{figure}[htbp]
  \centering
  \includegraphics[width=\linewidth]{assets/figure_p93_3.png}
\end{figure}

\begin{figure}[htbp]
  \centering
  \includegraphics[width=\linewidth]{assets/figure_p93_4.png}
\end{figure}

\begin{figure}[htbp]
  \centering
  \includegraphics[width=\linewidth]{assets/figure_p93_5.png}
\end{figure}

\begin{figure}[htbp]
  \centering
  \includegraphics[width=\linewidth]{assets/figure_p93_6.png}
\end{figure}

\begin{figure}[htbp]
  \centering
  \includegraphics[width=\linewidth]{assets/figure_p93_7.png}
\end{figure}

\begin{figure}[htbp]
  \centering
  \includegraphics[width=\linewidth]{assets/figure_p93_8.png}
\end{figure}

\begin{figure}[htbp]
  \centering
  \includegraphics[width=\linewidth]{assets/figure_p93_9.png}
\end{figure}

\begin{figure}[htbp]
  \centering
  \includegraphics[width=\linewidth]{assets/figure_p93_10.png}
\end{figure}

\begin{figure}[htbp]
  \centering
  \includegraphics[width=\linewidth]{assets/figure_p93_11.png}
\end{figure}

\begin{figure}[htbp]
  \centering
  \includegraphics[width=\linewidth]{assets/figure_p93_12.png}
\end{figure}

\begin{figure}[htbp]
  \centering
  \includegraphics[width=\linewidth]{assets/figure_p93_13.png}
\end{figure}

\begin{figure}[htbp]
  \centering
  \includegraphics[width=\linewidth]{assets/figure_p93_14.png}
\end{figure}

\begin{figure}[htbp]
  \centering
  \includegraphics[width=\linewidth]{assets/figure_p93_15.png}
\end{figure}

\begin{figure}[htbp]
  \centering
  \includegraphics[width=\linewidth]{assets/figure_p93_16.png}
\end{figure}

\begin{figure}[htbp]
  \centering
  \includegraphics[width=\linewidth]{assets/figure_p93_17.png}
\end{figure}

卡诺图是一种包含一些小方块的几何图形,图中每个小

方块称为一个单元,每个单元对应一个最小项.两个相邻的
最小项在卡诺图中也必须是相邻的.卡诺图中相邻的含义:

① 几何相邻性,即几何位置上相邻,也就是左右
紧挨着或者上下相接;

② 对称相邻性,即图形中对称位置的单元是相
邻的.

\begin{figure}[htbp]
  \centering
  \includegraphics[width=\linewidth]{assets/figure_p94_1.png}
\end{figure}

\begin{figure}[htbp]
  \centering
  \includegraphics[width=\linewidth]{assets/figure_p94_2.png}
\end{figure}

\begin{figure}[htbp]
  \centering
  \includegraphics[width=\linewidth]{assets/figure_p94_3.png}
\end{figure}

\begin{figure}[htbp]
  \centering
  \includegraphics[width=\linewidth]{assets/figure_p94_4.png}
\end{figure}

\begin{figure}[htbp]
  \centering
  \includegraphics[width=\linewidth]{assets/figure_p94_5.png}
\end{figure}

\begin{figure}[htbp]
  \centering
  \includegraphics[width=\linewidth]{assets/figure_p94_6.png}
\end{figure}

\begin{figure}[htbp]
  \centering
  \includegraphics[width=\linewidth]{assets/figure_p94_7.png}
\end{figure}

\begin{figure}[htbp]
  \centering
  \includegraphics[width=\linewidth]{assets/figure_p94_8.png}
\end{figure}

\begin{figure}[htbp]
  \centering
  \includegraphics[width=\linewidth]{assets/figure_p94_9.png}
\end{figure}

\begin{figure}[htbp]
  \centering
  \includegraphics[width=\linewidth]{assets/figure_p94_10.png}
\end{figure}

\begin{figure}[htbp]
  \centering
  \includegraphics[width=\linewidth]{assets/figure_p94_11.png}
\end{figure}

\begin{figure}[htbp]
  \centering
  \includegraphics[width=\linewidth]{assets/figure_p94_12.png}
\end{figure}

\begin{figure}[htbp]
  \centering
  \includegraphics[width=\linewidth]{assets/figure_p94_13.png}
\end{figure}

\begin{figure}[htbp]
  \centering
  \includegraphics[width=\linewidth]{assets/figure_p94_14.png}
\end{figure}

\begin{figure}[htbp]
  \centering
  \includegraphics[width=\linewidth]{assets/figure_p94_15.png}
\end{figure}

\begin{figure}[htbp]
  \centering
  \includegraphics[width=\linewidth]{assets/figure_p94_16.png}
\end{figure}

\begin{figure}[htbp]
  \centering
  \includegraphics[width=\linewidth]{assets/figure_p94_17.png}
\end{figure}

例 三变量卡诺图 循环码

\textbf{BC} 相邻性规则
\textbf{A 00 01 11 10}
\textbf{ABC ABC ABC ABC m m m}

\textbf{0 1 3 2}

\textbf{m}\textsubscript{\textbf{0}} \textbf{m m m}\textsubscript{\textbf{2}}

\textbf{1 3}

\textbf{ABC ABC ABC ABC}
\textbf{1 m}

\textbf{m4 m5 7 m6 m7}

相邻性规则
\textbf{m}\textsubscript{\textbf{2}} \textbf{m}\textsubscript{\textbf{0}} \textbf{m}\textsubscript{\textbf{1}}

（对称）

\textbf{m}\textsubscript{\textbf{4}}

\begin{figure}[htbp]
  \centering
  \includegraphics[width=\linewidth]{assets/figure_p95_1.png}
\end{figure}

\begin{figure}[htbp]
  \centering
  \includegraphics[width=\linewidth]{assets/figure_p95_2.png}
\end{figure}

\begin{figure}[htbp]
  \centering
  \includegraphics[width=\linewidth]{assets/figure_p95_3.png}
\end{figure}

\begin{figure}[htbp]
  \centering
  \includegraphics[width=\linewidth]{assets/figure_p95_4.png}
\end{figure}

\begin{figure}[htbp]
  \centering
  \includegraphics[width=\linewidth]{assets/figure_p95_5.png}
\end{figure}

\begin{figure}[htbp]
  \centering
  \includegraphics[width=\linewidth]{assets/figure_p95_6.png}
\end{figure}

\begin{figure}[htbp]
  \centering
  \includegraphics[width=\linewidth]{assets/figure_p95_7.png}
\end{figure}

\begin{figure}[htbp]
  \centering
  \includegraphics[width=\linewidth]{assets/figure_p95_8.png}
\end{figure}

\begin{figure}[htbp]
  \centering
  \includegraphics[width=\linewidth]{assets/figure_p95_9.png}
\end{figure}

\begin{figure}[htbp]
  \centering
  \includegraphics[width=\linewidth]{assets/figure_p95_10.png}
\end{figure}

\begin{figure}[htbp]
  \centering
  \includegraphics[width=\linewidth]{assets/figure_p95_11.png}
\end{figure}

\begin{figure}[htbp]
  \centering
  \includegraphics[width=\linewidth]{assets/figure_p95_12.png}
\end{figure}

\begin{figure}[htbp]
  \centering
  \includegraphics[width=\linewidth]{assets/figure_p95_13.png}
\end{figure}

\begin{figure}[htbp]
  \centering
  \includegraphics[width=\linewidth]{assets/figure_p95_14.png}
\end{figure}

\begin{figure}[htbp]
  \centering
  \includegraphics[width=\linewidth]{assets/figure_p95_15.png}
\end{figure}

\begin{figure}[htbp]
  \centering
  \includegraphics[width=\linewidth]{assets/figure_p95_16.png}
\end{figure}

\begin{figure}[htbp]
  \centering
  \includegraphics[width=\linewidth]{assets/figure_p95_17.png}
\end{figure}

二、四、五变量卡诺图

\textbf{B}
\textbf{A 0 1}
\textbf{CD}
\textbf{0 0 1}
\textbf{AB 00 01 11 10}

\textbf{1 2 3 00}
\textbf{0 1 3 2}
\textbf{01 4 5 7 6}

相邻性规则

m3 \textbf{11 12 13 15 14}

\textbf{10 8 9 11 10}

m5 m7 m6

m15

\begin{figure}[htbp]
  \centering
  \includegraphics[width=\linewidth]{assets/figure_p96_1.png}
\end{figure}

\begin{figure}[htbp]
  \centering
  \includegraphics[width=\linewidth]{assets/figure_p96_2.png}
\end{figure}

\begin{figure}[htbp]
  \centering
  \includegraphics[width=\linewidth]{assets/figure_p96_3.png}
\end{figure}

\begin{figure}[htbp]
  \centering
  \includegraphics[width=\linewidth]{assets/figure_p96_4.png}
\end{figure}

\begin{figure}[htbp]
  \centering
  \includegraphics[width=\linewidth]{assets/figure_p96_5.png}
\end{figure}

\begin{figure}[htbp]
  \centering
  \includegraphics[width=\linewidth]{assets/figure_p96_6.png}
\end{figure}

\begin{figure}[htbp]
  \centering
  \includegraphics[width=\linewidth]{assets/figure_p96_7.png}
\end{figure}

\begin{figure}[htbp]
  \centering
  \includegraphics[width=\linewidth]{assets/figure_p96_8.png}
\end{figure}

\begin{figure}[htbp]
  \centering
  \includegraphics[width=\linewidth]{assets/figure_p96_9.png}
\end{figure}

\begin{figure}[htbp]
  \centering
  \includegraphics[width=\linewidth]{assets/figure_p96_10.png}
\end{figure}

\begin{figure}[htbp]
  \centering
  \includegraphics[width=\linewidth]{assets/figure_p96_11.png}
\end{figure}

\begin{figure}[htbp]
  \centering
  \includegraphics[width=\linewidth]{assets/figure_p96_12.png}
\end{figure}

\begin{figure}[htbp]
  \centering
  \includegraphics[width=\linewidth]{assets/figure_p96_13.png}
\end{figure}

\begin{figure}[htbp]
  \centering
  \includegraphics[width=\linewidth]{assets/figure_p96_14.png}
\end{figure}

\begin{figure}[htbp]
  \centering
  \includegraphics[width=\linewidth]{assets/figure_p96_15.png}
\end{figure}

\begin{figure}[htbp]
  \centering
  \includegraphics[width=\linewidth]{assets/figure_p96_16.png}
\end{figure}

\begin{figure}[htbp]
  \centering
  \includegraphics[width=\linewidth]{assets/figure_p96_17.png}
\end{figure}

\textbf{CDE}
\textbf{AB 000 001 011 010 110 111 101 100}
\textbf{00}
\textbf{0 1 3 2 6 7 5 4}

\textbf{01 8 9 11 10 14 15 13 12}

\textbf{11 24 25 27 26 30 31 29 28}

\textbf{10 22 23 21 20}
\textbf{16 17 19 18}

2）逻辑函数的卡诺图表示法

\begin{figure}[htbp]
  \centering
  \includegraphics[width=\linewidth]{assets/figure_p97_1.png}
\end{figure}

\begin{figure}[htbp]
  \centering
  \includegraphics[width=\linewidth]{assets/figure_p97_2.png}
\end{figure}

\begin{figure}[htbp]
  \centering
  \includegraphics[width=\linewidth]{assets/figure_p97_3.png}
\end{figure}

\begin{figure}[htbp]
  \centering
  \includegraphics[width=\linewidth]{assets/figure_p97_4.png}
\end{figure}

\begin{figure}[htbp]
  \centering
  \includegraphics[width=\linewidth]{assets/figure_p97_5.png}
\end{figure}

\begin{figure}[htbp]
  \centering
  \includegraphics[width=\linewidth]{assets/figure_p97_6.png}
\end{figure}

\begin{figure}[htbp]
  \centering
  \includegraphics[width=\linewidth]{assets/figure_p97_7.png}
\end{figure}

\begin{figure}[htbp]
  \centering
  \includegraphics[width=\linewidth]{assets/figure_p97_8.png}
\end{figure}

\begin{figure}[htbp]
  \centering
  \includegraphics[width=\linewidth]{assets/figure_p97_9.png}
\end{figure}

\begin{figure}[htbp]
  \centering
  \includegraphics[width=\linewidth]{assets/figure_p97_10.png}
\end{figure}

\begin{figure}[htbp]
  \centering
  \includegraphics[width=\linewidth]{assets/figure_p97_11.png}
\end{figure}

\begin{figure}[htbp]
  \centering
  \includegraphics[width=\linewidth]{assets/figure_p97_12.png}
\end{figure}

\begin{figure}[htbp]
  \centering
  \includegraphics[width=\linewidth]{assets/figure_p97_13.png}
\end{figure}

\begin{figure}[htbp]
  \centering
  \includegraphics[width=\linewidth]{assets/figure_p97_14.png}
\end{figure}

\begin{figure}[htbp]
  \centering
  \includegraphics[width=\linewidth]{assets/figure_p97_15.png}
\end{figure}

\begin{figure}[htbp]
  \centering
  \includegraphics[width=\linewidth]{assets/figure_p97_16.png}
\end{figure}

\begin{figure}[htbp]
  \centering
  \includegraphics[width=\linewidth]{assets/figure_p97_17.png}
\end{figure}

用卡诺图表示逻辑函数，只是把各组变量值所对应的

逻辑函数F的值，填在对应的小方格中。
（其实卡诺图是真值表的另一种画法）

\[ 例： F （ A ， B ， C ） =ABC+ABC+ABC \]

用卡诺图表示为：

\textbf{BC}
\textbf{A 00 01 11 10}
\textbf{0 0 0 1 0}
\textbf{m}\textsubscript{\textbf{3}}
\textbf{1 0 1 1 0}
\textbf{m m}\textsubscript{\textbf{7}}

\textbf{5}

\begin{figure}[htbp]
  \centering
  \includegraphics[width=\linewidth]{assets/figure_p98_1.png}
\end{figure}

\begin{figure}[htbp]
  \centering
  \includegraphics[width=\linewidth]{assets/figure_p98_2.png}
\end{figure}

\begin{figure}[htbp]
  \centering
  \includegraphics[width=\linewidth]{assets/figure_p98_3.png}
\end{figure}

\begin{figure}[htbp]
  \centering
  \includegraphics[width=\linewidth]{assets/figure_p98_4.png}
\end{figure}

\begin{figure}[htbp]
  \centering
  \includegraphics[width=\linewidth]{assets/figure_p98_5.png}
\end{figure}

\begin{figure}[htbp]
  \centering
  \includegraphics[width=\linewidth]{assets/figure_p98_6.png}
\end{figure}

\begin{figure}[htbp]
  \centering
  \includegraphics[width=\linewidth]{assets/figure_p98_7.png}
\end{figure}

\begin{figure}[htbp]
  \centering
  \includegraphics[width=\linewidth]{assets/figure_p98_8.png}
\end{figure}

\begin{figure}[htbp]
  \centering
  \includegraphics[width=\linewidth]{assets/figure_p98_9.png}
\end{figure}

\begin{figure}[htbp]
  \centering
  \includegraphics[width=\linewidth]{assets/figure_p98_10.png}
\end{figure}

\begin{figure}[htbp]
  \centering
  \includegraphics[width=\linewidth]{assets/figure_p98_11.png}
\end{figure}

\begin{figure}[htbp]
  \centering
  \includegraphics[width=\linewidth]{assets/figure_p98_12.png}
\end{figure}

\begin{figure}[htbp]
  \centering
  \includegraphics[width=\linewidth]{assets/figure_p98_13.png}
\end{figure}

\begin{figure}[htbp]
  \centering
  \includegraphics[width=\linewidth]{assets/figure_p98_14.png}
\end{figure}

\begin{figure}[htbp]
  \centering
  \includegraphics[width=\linewidth]{assets/figure_p98_15.png}
\end{figure}

\begin{figure}[htbp]
  \centering
  \includegraphics[width=\linewidth]{assets/figure_p98_16.png}
\end{figure}

\begin{figure}[htbp]
  \centering
  \includegraphics[width=\linewidth]{assets/figure_p98_17.png}
\end{figure}

3) 在卡诺图上合并最小项的规则

当卡诺图中有最小项相邻时（即：有标1的方格相邻)，

可利用最小项相邻的性质，对最小项合并。
规则为：

（1） 卡诺图上任何两个标1的方格相邻，可以合为1
项，并可消去1个变量。

\begin{figure}[htbp]
  \centering
  \includegraphics[width=\linewidth]{assets/figure_p99_1.png}
\end{figure}

\begin{figure}[htbp]
  \centering
  \includegraphics[width=\linewidth]{assets/figure_p99_2.png}
\end{figure}

\begin{figure}[htbp]
  \centering
  \includegraphics[width=\linewidth]{assets/figure_p99_3.png}
\end{figure}

\begin{figure}[htbp]
  \centering
  \includegraphics[width=\linewidth]{assets/figure_p99_4.png}
\end{figure}

\begin{figure}[htbp]
  \centering
  \includegraphics[width=\linewidth]{assets/figure_p99_5.png}
\end{figure}

\begin{figure}[htbp]
  \centering
  \includegraphics[width=\linewidth]{assets/figure_p99_6.png}
\end{figure}

\begin{figure}[htbp]
  \centering
  \includegraphics[width=\linewidth]{assets/figure_p99_7.png}
\end{figure}

\begin{figure}[htbp]
  \centering
  \includegraphics[width=\linewidth]{assets/figure_p99_8.png}
\end{figure}

\begin{figure}[htbp]
  \centering
  \includegraphics[width=\linewidth]{assets/figure_p99_9.png}
\end{figure}

\begin{figure}[htbp]
  \centering
  \includegraphics[width=\linewidth]{assets/figure_p99_10.png}
\end{figure}

\begin{figure}[htbp]
  \centering
  \includegraphics[width=\linewidth]{assets/figure_p99_11.png}
\end{figure}

\begin{figure}[htbp]
  \centering
  \includegraphics[width=\linewidth]{assets/figure_p99_12.png}
\end{figure}

\begin{figure}[htbp]
  \centering
  \includegraphics[width=\linewidth]{assets/figure_p99_13.png}
\end{figure}

\begin{figure}[htbp]
  \centering
  \includegraphics[width=\linewidth]{assets/figure_p99_14.png}
\end{figure}

\begin{figure}[htbp]
  \centering
  \includegraphics[width=\linewidth]{assets/figure_p99_15.png}
\end{figure}

\begin{figure}[htbp]
  \centering
  \includegraphics[width=\linewidth]{assets/figure_p99_16.png}
\end{figure}

\begin{figure}[htbp]
  \centering
  \includegraphics[width=\linewidth]{assets/figure_p99_17.png}
\end{figure}

例： \textbf{BC}

\textbf{A 00 01 11 10}
\textbf{0 0 0 1 0 ABC+ABC}

\textbf{=BC}

\textbf{1 0 1 1 0}

\textbf{ABC+ABC=AC}

\begin{figure}[htbp]
  \centering
  \includegraphics[width=\linewidth]{assets/figure_p100_1.png}
\end{figure}

\begin{figure}[htbp]
  \centering
  \includegraphics[width=\linewidth]{assets/figure_p100_2.png}
\end{figure}

\begin{figure}[htbp]
  \centering
  \includegraphics[width=\linewidth]{assets/figure_p100_3.png}
\end{figure}

\begin{figure}[htbp]
  \centering
  \includegraphics[width=\linewidth]{assets/figure_p100_4.png}
\end{figure}

\begin{figure}[htbp]
  \centering
  \includegraphics[width=\linewidth]{assets/figure_p100_5.png}
\end{figure}

\begin{figure}[htbp]
  \centering
  \includegraphics[width=\linewidth]{assets/figure_p100_6.png}
\end{figure}

\begin{figure}[htbp]
  \centering
  \includegraphics[width=\linewidth]{assets/figure_p100_7.png}
\end{figure}

\begin{figure}[htbp]
  \centering
  \includegraphics[width=\linewidth]{assets/figure_p100_8.png}
\end{figure}

\begin{figure}[htbp]
  \centering
  \includegraphics[width=\linewidth]{assets/figure_p100_9.png}
\end{figure}

\begin{figure}[htbp]
  \centering
  \includegraphics[width=\linewidth]{assets/figure_p100_10.png}
\end{figure}

\begin{figure}[htbp]
  \centering
  \includegraphics[width=\linewidth]{assets/figure_p100_11.png}
\end{figure}

\begin{figure}[htbp]
  \centering
  \includegraphics[width=\linewidth]{assets/figure_p100_12.png}
\end{figure}

\begin{figure}[htbp]
  \centering
  \includegraphics[width=\linewidth]{assets/figure_p100_13.png}
\end{figure}

\begin{figure}[htbp]
  \centering
  \includegraphics[width=\linewidth]{assets/figure_p100_14.png}
\end{figure}

\begin{figure}[htbp]
  \centering
  \includegraphics[width=\linewidth]{assets/figure_p100_15.png}
\end{figure}

\begin{figure}[htbp]
  \centering
  \includegraphics[width=\linewidth]{assets/figure_p100_16.png}
\end{figure}

\begin{figure}[htbp]
  \centering
  \includegraphics[width=\linewidth]{assets/figure_p100_17.png}
\end{figure}

\textbf{CD}
\textbf{AB 00 01 11 10 ABD}
\textbf{00 1 1}

\textbf{01}
\textbf{ABD}
\textbf{11 1 1}

\textbf{10}

（2）卡诺图上任何四个标1方格相邻，可合并为一项，并
可消去两个变量。
四个标1方格相邻的特点：
①同在一行或一列； ②同在一田字格中。

\begin{figure}[htbp]
  \centering
  \includegraphics[width=\linewidth]{assets/figure_p101_1.png}
\end{figure}

\begin{figure}[htbp]
  \centering
  \includegraphics[width=\linewidth]{assets/figure_p101_2.png}
\end{figure}

\begin{figure}[htbp]
  \centering
  \includegraphics[width=\linewidth]{assets/figure_p101_3.png}
\end{figure}

\begin{figure}[htbp]
  \centering
  \includegraphics[width=\linewidth]{assets/figure_p101_4.png}
\end{figure}

\begin{figure}[htbp]
  \centering
  \includegraphics[width=\linewidth]{assets/figure_p101_5.png}
\end{figure}

\begin{figure}[htbp]
  \centering
  \includegraphics[width=\linewidth]{assets/figure_p101_6.png}
\end{figure}

\begin{figure}[htbp]
  \centering
  \includegraphics[width=\linewidth]{assets/figure_p101_7.png}
\end{figure}

\begin{figure}[htbp]
  \centering
  \includegraphics[width=\linewidth]{assets/figure_p101_8.png}
\end{figure}

\begin{figure}[htbp]
  \centering
  \includegraphics[width=\linewidth]{assets/figure_p101_9.png}
\end{figure}

\begin{figure}[htbp]
  \centering
  \includegraphics[width=\linewidth]{assets/figure_p101_10.png}
\end{figure}

\begin{figure}[htbp]
  \centering
  \includegraphics[width=\linewidth]{assets/figure_p101_11.png}
\end{figure}

\begin{figure}[htbp]
  \centering
  \includegraphics[width=\linewidth]{assets/figure_p101_12.png}
\end{figure}

\begin{figure}[htbp]
  \centering
  \includegraphics[width=\linewidth]{assets/figure_p101_13.png}
\end{figure}

\begin{figure}[htbp]
  \centering
  \includegraphics[width=\linewidth]{assets/figure_p101_14.png}
\end{figure}

\begin{figure}[htbp]
  \centering
  \includegraphics[width=\linewidth]{assets/figure_p101_15.png}
\end{figure}

\begin{figure}[htbp]
  \centering
  \includegraphics[width=\linewidth]{assets/figure_p101_16.png}
\end{figure}

\begin{figure}[htbp]
  \centering
  \includegraphics[width=\linewidth]{assets/figure_p101_17.png}
\end{figure}

同在一个田字格中
例：

\textbf{CD AB CD}
\textbf{AB 00 01 11 10 AB 00 01 11 10}

\textbf{00 1 1 1 1 00 11 1}

\textbf{01 1 01 1 1}

\textbf{11 1 11 1 1}

\textbf{10 1 10 1 1}

\textbf{CD BD}
\textbf{BD}

同在一行或一列

\begin{figure}[htbp]
  \centering
  \includegraphics[width=\linewidth]{assets/figure_p102_1.png}
\end{figure}

\begin{figure}[htbp]
  \centering
  \includegraphics[width=\linewidth]{assets/figure_p102_2.png}
\end{figure}

\begin{figure}[htbp]
  \centering
  \includegraphics[width=\linewidth]{assets/figure_p102_3.png}
\end{figure}

\begin{figure}[htbp]
  \centering
  \includegraphics[width=\linewidth]{assets/figure_p102_4.png}
\end{figure}

\begin{figure}[htbp]
  \centering
  \includegraphics[width=\linewidth]{assets/figure_p102_5.png}
\end{figure}

\begin{figure}[htbp]
  \centering
  \includegraphics[width=\linewidth]{assets/figure_p102_6.png}
\end{figure}

\begin{figure}[htbp]
  \centering
  \includegraphics[width=\linewidth]{assets/figure_p102_7.png}
\end{figure}

\begin{figure}[htbp]
  \centering
  \includegraphics[width=\linewidth]{assets/figure_p102_8.png}
\end{figure}

\begin{figure}[htbp]
  \centering
  \includegraphics[width=\linewidth]{assets/figure_p102_9.png}
\end{figure}

\begin{figure}[htbp]
  \centering
  \includegraphics[width=\linewidth]{assets/figure_p102_10.png}
\end{figure}

\begin{figure}[htbp]
  \centering
  \includegraphics[width=\linewidth]{assets/figure_p102_11.png}
\end{figure}

\begin{figure}[htbp]
  \centering
  \includegraphics[width=\linewidth]{assets/figure_p102_12.png}
\end{figure}

\begin{figure}[htbp]
  \centering
  \includegraphics[width=\linewidth]{assets/figure_p102_13.png}
\end{figure}

\begin{figure}[htbp]
  \centering
  \includegraphics[width=\linewidth]{assets/figure_p102_14.png}
\end{figure}

\begin{figure}[htbp]
  \centering
  \includegraphics[width=\linewidth]{assets/figure_p102_15.png}
\end{figure}

\begin{figure}[htbp]
  \centering
  \includegraphics[width=\linewidth]{assets/figure_p102_16.png}
\end{figure}

\begin{figure}[htbp]
  \centering
  \includegraphics[width=\linewidth]{assets/figure_p102_17.png}
\end{figure}

思考题

\textbf{CD CD}
\textbf{AB 00 01 11 10 AB 00 01 11 10}

\textbf{00 1 1 1 1 00 1 1}

\textbf{01 1} 1 \textbf{01} 1 1

\textbf{11 1} 1 \textbf{11} 1 1

\textbf{10 10} 1 \textbf{1}

\begin{figure}[htbp]
  \centering
  \includegraphics[width=\linewidth]{assets/figure_p103_1.png}
\end{figure}

\begin{figure}[htbp]
  \centering
  \includegraphics[width=\linewidth]{assets/figure_p103_2.png}
\end{figure}

\begin{figure}[htbp]
  \centering
  \includegraphics[width=\linewidth]{assets/figure_p103_3.png}
\end{figure}

\begin{figure}[htbp]
  \centering
  \includegraphics[width=\linewidth]{assets/figure_p103_4.png}
\end{figure}

\begin{figure}[htbp]
  \centering
  \includegraphics[width=\linewidth]{assets/figure_p103_5.png}
\end{figure}

\begin{figure}[htbp]
  \centering
  \includegraphics[width=\linewidth]{assets/figure_p103_6.png}
\end{figure}

\begin{figure}[htbp]
  \centering
  \includegraphics[width=\linewidth]{assets/figure_p103_7.png}
\end{figure}

\begin{figure}[htbp]
  \centering
  \includegraphics[width=\linewidth]{assets/figure_p103_8.png}
\end{figure}

\begin{figure}[htbp]
  \centering
  \includegraphics[width=\linewidth]{assets/figure_p103_9.png}
\end{figure}

\begin{figure}[htbp]
  \centering
  \includegraphics[width=\linewidth]{assets/figure_p103_10.png}
\end{figure}

\begin{figure}[htbp]
  \centering
  \includegraphics[width=\linewidth]{assets/figure_p103_11.png}
\end{figure}

\begin{figure}[htbp]
  \centering
  \includegraphics[width=\linewidth]{assets/figure_p103_12.png}
\end{figure}

\begin{figure}[htbp]
  \centering
  \includegraphics[width=\linewidth]{assets/figure_p103_13.png}
\end{figure}

\begin{figure}[htbp]
  \centering
  \includegraphics[width=\linewidth]{assets/figure_p103_14.png}
\end{figure}

\begin{figure}[htbp]
  \centering
  \includegraphics[width=\linewidth]{assets/figure_p103_15.png}
\end{figure}

\begin{figure}[htbp]
  \centering
  \includegraphics[width=\linewidth]{assets/figure_p103_16.png}
\end{figure}

\begin{figure}[htbp]
  \centering
  \includegraphics[width=\linewidth]{assets/figure_p103_17.png}
\end{figure}

（3）卡诺图上任何八个标1的方格相邻，可以并为一
项，并可消去三个变量。例：

\textbf{CD CD}
\textbf{AB 00 01 AB 00 01 11 10}
\textbf{11 10}
\textbf{00 1 00 1 1 1 1}
\textbf{1 1 1}

\textbf{01 1 1 1 1 01}

\textbf{A}

\textbf{11 11}
\textbf{10 10 1 1 1 1}

\textbf{B}

\begin{figure}[htbp]
  \centering
  \includegraphics[width=\linewidth]{assets/figure_p104_1.png}
\end{figure}

\begin{figure}[htbp]
  \centering
  \includegraphics[width=\linewidth]{assets/figure_p104_2.png}
\end{figure}

\begin{figure}[htbp]
  \centering
  \includegraphics[width=\linewidth]{assets/figure_p104_3.png}
\end{figure}

\begin{figure}[htbp]
  \centering
  \includegraphics[width=\linewidth]{assets/figure_p104_4.png}
\end{figure}

\begin{figure}[htbp]
  \centering
  \includegraphics[width=\linewidth]{assets/figure_p104_5.png}
\end{figure}

\begin{figure}[htbp]
  \centering
  \includegraphics[width=\linewidth]{assets/figure_p104_6.png}
\end{figure}

\begin{figure}[htbp]
  \centering
  \includegraphics[width=\linewidth]{assets/figure_p104_7.png}
\end{figure}

\begin{figure}[htbp]
  \centering
  \includegraphics[width=\linewidth]{assets/figure_p104_8.png}
\end{figure}

\begin{figure}[htbp]
  \centering
  \includegraphics[width=\linewidth]{assets/figure_p104_9.png}
\end{figure}

\begin{figure}[htbp]
  \centering
  \includegraphics[width=\linewidth]{assets/figure_p104_10.png}
\end{figure}

\begin{figure}[htbp]
  \centering
  \includegraphics[width=\linewidth]{assets/figure_p104_11.png}
\end{figure}

\begin{figure}[htbp]
  \centering
  \includegraphics[width=\linewidth]{assets/figure_p104_12.png}
\end{figure}

\begin{figure}[htbp]
  \centering
  \includegraphics[width=\linewidth]{assets/figure_p104_13.png}
\end{figure}

\begin{figure}[htbp]
  \centering
  \includegraphics[width=\linewidth]{assets/figure_p104_14.png}
\end{figure}

\begin{figure}[htbp]
  \centering
  \includegraphics[width=\linewidth]{assets/figure_p104_15.png}
\end{figure}

\begin{figure}[htbp]
  \centering
  \includegraphics[width=\linewidth]{assets/figure_p104_16.png}
\end{figure}

\begin{figure}[htbp]
  \centering
  \includegraphics[width=\linewidth]{assets/figure_p104_17.png}
\end{figure}

\textbf{CD}
思考题: \textbf{AB 00 01 11 10}

\textbf{00 1 1}

\textbf{01 1 1}

\textbf{11 1} 1

\textbf{10 1 1}

综上所述，在\textbf{n}个变量的卡诺图中，只有\textbf{2}的\textbf{i}次方个相邻的标\textbf{1}方
格（必须排列成方形格或矩形格的形状）才能圈在一起，合并为
一项，该项保留了原来各项中\textbf{n-i}个相同的变量，消去\textbf{i}个不同变
量。

\begin{figure}[htbp]
  \centering
  \includegraphics[width=\linewidth]{assets/figure_p105_1.png}
\end{figure}

\begin{figure}[htbp]
  \centering
  \includegraphics[width=\linewidth]{assets/figure_p105_2.png}
\end{figure}

\begin{figure}[htbp]
  \centering
  \includegraphics[width=\linewidth]{assets/figure_p105_3.png}
\end{figure}

\begin{figure}[htbp]
  \centering
  \includegraphics[width=\linewidth]{assets/figure_p105_4.png}
\end{figure}

\begin{figure}[htbp]
  \centering
  \includegraphics[width=\linewidth]{assets/figure_p105_5.png}
\end{figure}

\begin{figure}[htbp]
  \centering
  \includegraphics[width=\linewidth]{assets/figure_p105_6.png}
\end{figure}

\begin{figure}[htbp]
  \centering
  \includegraphics[width=\linewidth]{assets/figure_p105_7.png}
\end{figure}

\begin{figure}[htbp]
  \centering
  \includegraphics[width=\linewidth]{assets/figure_p105_8.png}
\end{figure}

\begin{figure}[htbp]
  \centering
  \includegraphics[width=\linewidth]{assets/figure_p105_9.png}
\end{figure}

\begin{figure}[htbp]
  \centering
  \includegraphics[width=\linewidth]{assets/figure_p105_10.png}
\end{figure}

\begin{figure}[htbp]
  \centering
  \includegraphics[width=\linewidth]{assets/figure_p105_11.png}
\end{figure}

\begin{figure}[htbp]
  \centering
  \includegraphics[width=\linewidth]{assets/figure_p105_12.png}
\end{figure}

\begin{figure}[htbp]
  \centering
  \includegraphics[width=\linewidth]{assets/figure_p105_13.png}
\end{figure}

\begin{figure}[htbp]
  \centering
  \includegraphics[width=\linewidth]{assets/figure_p105_14.png}
\end{figure}

\begin{figure}[htbp]
  \centering
  \includegraphics[width=\linewidth]{assets/figure_p105_15.png}
\end{figure}

\begin{figure}[htbp]
  \centering
  \includegraphics[width=\linewidth]{assets/figure_p105_16.png}
\end{figure}

\begin{figure}[htbp]
  \centering
  \includegraphics[width=\linewidth]{assets/figure_p105_17.png}
\end{figure}

4） 用卡诺图化简逻辑函数（化为最简与或式）

①最简标准：项数最少，意味着卡诺图中圈数最少；
每项中的变量数最少，意味着卡诺图中
的圈尽可能大。

$例 将\textbf{F}（\textbf{A}，\textbf{B}，\textbf{C}）\textbf{=\Sigma m}（\textbf{3}，\textbf{4}，\textbf{5}，\textbf{6}，\textbf{7}）$

化为最简与或式。

\begin{figure}[htbp]
  \centering
  \includegraphics[width=\linewidth]{assets/figure_p106_1.png}
\end{figure}

\begin{figure}[htbp]
  \centering
  \includegraphics[width=\linewidth]{assets/figure_p106_2.png}
\end{figure}

\begin{figure}[htbp]
  \centering
  \includegraphics[width=\linewidth]{assets/figure_p106_3.png}
\end{figure}

\begin{figure}[htbp]
  \centering
  \includegraphics[width=\linewidth]{assets/figure_p106_4.png}
\end{figure}

\begin{figure}[htbp]
  \centering
  \includegraphics[width=\linewidth]{assets/figure_p106_5.png}
\end{figure}

\begin{figure}[htbp]
  \centering
  \includegraphics[width=\linewidth]{assets/figure_p106_6.png}
\end{figure}

\begin{figure}[htbp]
  \centering
  \includegraphics[width=\linewidth]{assets/figure_p106_7.png}
\end{figure}

\begin{figure}[htbp]
  \centering
  \includegraphics[width=\linewidth]{assets/figure_p106_8.png}
\end{figure}

\begin{figure}[htbp]
  \centering
  \includegraphics[width=\linewidth]{assets/figure_p106_9.png}
\end{figure}

\begin{figure}[htbp]
  \centering
  \includegraphics[width=\linewidth]{assets/figure_p106_10.png}
\end{figure}

\begin{figure}[htbp]
  \centering
  \includegraphics[width=\linewidth]{assets/figure_p106_11.png}
\end{figure}

\begin{figure}[htbp]
  \centering
  \includegraphics[width=\linewidth]{assets/figure_p106_12.png}
\end{figure}

\begin{figure}[htbp]
  \centering
  \includegraphics[width=\linewidth]{assets/figure_p106_13.png}
\end{figure}

\begin{figure}[htbp]
  \centering
  \includegraphics[width=\linewidth]{assets/figure_p106_14.png}
\end{figure}

\begin{figure}[htbp]
  \centering
  \includegraphics[width=\linewidth]{assets/figure_p106_15.png}
\end{figure}

\begin{figure}[htbp]
  \centering
  \includegraphics[width=\linewidth]{assets/figure_p106_16.png}
\end{figure}

\begin{figure}[htbp]
  \centering
  \includegraphics[width=\linewidth]{assets/figure_p106_17.png}
\end{figure}

\textbf{BC BC}
\textbf{A 00 01 11 10 A 00 01 11 10}
\textbf{0 1 0 1}

\textbf{1 1 1 1 1 1 1 1 1 1}

\textbf{F=A+BC}（最简） \textbf{F=AB+BC+ABC}（非最简）
\textbf{\{=A(B+BC)+BC}
\textbf{=A(B+C)+BC}
\textbf{=ABC+BC}
\textbf{=A+BC\}}
② 化简步骤（结合举例说明）

\begin{figure}[htbp]
  \centering
  \includegraphics[width=\linewidth]{assets/figure_p107_1.png}
\end{figure}

\begin{figure}[htbp]
  \centering
  \includegraphics[width=\linewidth]{assets/figure_p107_2.png}
\end{figure}

\begin{figure}[htbp]
  \centering
  \includegraphics[width=\linewidth]{assets/figure_p107_3.png}
\end{figure}

\begin{figure}[htbp]
  \centering
  \includegraphics[width=\linewidth]{assets/figure_p107_4.png}
\end{figure}

\begin{figure}[htbp]
  \centering
  \includegraphics[width=\linewidth]{assets/figure_p107_5.png}
\end{figure}

\begin{figure}[htbp]
  \centering
  \includegraphics[width=\linewidth]{assets/figure_p107_6.png}
\end{figure}

\begin{figure}[htbp]
  \centering
  \includegraphics[width=\linewidth]{assets/figure_p107_7.png}
\end{figure}

\begin{figure}[htbp]
  \centering
  \includegraphics[width=\linewidth]{assets/figure_p107_8.png}
\end{figure}

\begin{figure}[htbp]
  \centering
  \includegraphics[width=\linewidth]{assets/figure_p107_9.png}
\end{figure}

\begin{figure}[htbp]
  \centering
  \includegraphics[width=\linewidth]{assets/figure_p107_10.png}
\end{figure}

\begin{figure}[htbp]
  \centering
  \includegraphics[width=\linewidth]{assets/figure_p107_11.png}
\end{figure}

\begin{figure}[htbp]
  \centering
  \includegraphics[width=\linewidth]{assets/figure_p107_12.png}
\end{figure}

\begin{figure}[htbp]
  \centering
  \includegraphics[width=\linewidth]{assets/figure_p107_13.png}
\end{figure}

\begin{figure}[htbp]
  \centering
  \includegraphics[width=\linewidth]{assets/figure_p107_14.png}
\end{figure}

\begin{figure}[htbp]
  \centering
  \includegraphics[width=\linewidth]{assets/figure_p107_15.png}
\end{figure}

\begin{figure}[htbp]
  \centering
  \includegraphics[width=\linewidth]{assets/figure_p107_16.png}
\end{figure}

\begin{figure}[htbp]
  \centering
  \includegraphics[width=\linewidth]{assets/figure_p107_17.png}
\end{figure}

$例 将\textbf{F}（\textbf{A,B,C,D}）\textbf{=\Sigma m}（\textbf{0,1,3,7,8,10,13}）化为最简与$

或式。
解： (1) 由表达式填卡诺图; \textbf{CD}

\textbf{AB 00 01 11 10}

\textbf{00 1 1 1}

(2) 圈出孤立的标1方格;

\textbf{01 1}

\textbf{m}13

\textbf{11 1}

\textbf{10 1 1}

\begin{figure}[htbp]
  \centering
  \includegraphics[width=\linewidth]{assets/figure_p108_1.png}
\end{figure}

\begin{figure}[htbp]
  \centering
  \includegraphics[width=\linewidth]{assets/figure_p108_2.png}
\end{figure}

\begin{figure}[htbp]
  \centering
  \includegraphics[width=\linewidth]{assets/figure_p108_3.png}
\end{figure}

\begin{figure}[htbp]
  \centering
  \includegraphics[width=\linewidth]{assets/figure_p108_4.png}
\end{figure}

\begin{figure}[htbp]
  \centering
  \includegraphics[width=\linewidth]{assets/figure_p108_5.png}
\end{figure}

\begin{figure}[htbp]
  \centering
  \includegraphics[width=\linewidth]{assets/figure_p108_6.png}
\end{figure}

\begin{figure}[htbp]
  \centering
  \includegraphics[width=\linewidth]{assets/figure_p108_7.png}
\end{figure}

\begin{figure}[htbp]
  \centering
  \includegraphics[width=\linewidth]{assets/figure_p108_8.png}
\end{figure}

\begin{figure}[htbp]
  \centering
  \includegraphics[width=\linewidth]{assets/figure_p108_9.png}
\end{figure}

\begin{figure}[htbp]
  \centering
  \includegraphics[width=\linewidth]{assets/figure_p108_10.png}
\end{figure}

\begin{figure}[htbp]
  \centering
  \includegraphics[width=\linewidth]{assets/figure_p108_11.png}
\end{figure}

\begin{figure}[htbp]
  \centering
  \includegraphics[width=\linewidth]{assets/figure_p108_12.png}
\end{figure}

\begin{figure}[htbp]
  \centering
  \includegraphics[width=\linewidth]{assets/figure_p108_13.png}
\end{figure}

\begin{figure}[htbp]
  \centering
  \includegraphics[width=\linewidth]{assets/figure_p108_14.png}
\end{figure}

\begin{figure}[htbp]
  \centering
  \includegraphics[width=\linewidth]{assets/figure_p108_15.png}
\end{figure}

\begin{figure}[htbp]
  \centering
  \includegraphics[width=\linewidth]{assets/figure_p108_16.png}
\end{figure}

\begin{figure}[htbp]
  \centering
  \includegraphics[width=\linewidth]{assets/figure_p108_17.png}
\end{figure}

\textbf{CD ABC}
\textbf{AB 00 01 11 10}

\textbf{00 1 1 1 ACD}

\textbf{01 1}

\textbf{11 1 ABCD}

\textbf{10 1 1}

\textbf{ABD}

(3) 找出只被一个最大的圈所覆盖的标1方格,并

圈出覆盖该标1方格的最大圈; \textbf{m7,m10}

(4) 将剩余的相邻标1方格,圈成尽可能少,而且

尽可能大的圈. \textbf{m0,m1}

\begin{figure}[htbp]
  \centering
  \includegraphics[width=\linewidth]{assets/figure_p109_1.png}
\end{figure}

\begin{figure}[htbp]
  \centering
  \includegraphics[width=\linewidth]{assets/figure_p109_2.png}
\end{figure}

\begin{figure}[htbp]
  \centering
  \includegraphics[width=\linewidth]{assets/figure_p109_3.png}
\end{figure}

\begin{figure}[htbp]
  \centering
  \includegraphics[width=\linewidth]{assets/figure_p109_4.png}
\end{figure}

\begin{figure}[htbp]
  \centering
  \includegraphics[width=\linewidth]{assets/figure_p109_5.png}
\end{figure}

\begin{figure}[htbp]
  \centering
  \includegraphics[width=\linewidth]{assets/figure_p109_6.png}
\end{figure}

\begin{figure}[htbp]
  \centering
  \includegraphics[width=\linewidth]{assets/figure_p109_7.png}
\end{figure}

\begin{figure}[htbp]
  \centering
  \includegraphics[width=\linewidth]{assets/figure_p109_8.png}
\end{figure}

\begin{figure}[htbp]
  \centering
  \includegraphics[width=\linewidth]{assets/figure_p109_9.png}
\end{figure}

\begin{figure}[htbp]
  \centering
  \includegraphics[width=\linewidth]{assets/figure_p109_10.png}
\end{figure}

\begin{figure}[htbp]
  \centering
  \includegraphics[width=\linewidth]{assets/figure_p109_11.png}
\end{figure}

\begin{figure}[htbp]
  \centering
  \includegraphics[width=\linewidth]{assets/figure_p109_12.png}
\end{figure}

\begin{figure}[htbp]
  \centering
  \includegraphics[width=\linewidth]{assets/figure_p109_13.png}
\end{figure}

\begin{figure}[htbp]
  \centering
  \includegraphics[width=\linewidth]{assets/figure_p109_14.png}
\end{figure}

\begin{figure}[htbp]
  \centering
  \includegraphics[width=\linewidth]{assets/figure_p109_15.png}
\end{figure}

\begin{figure}[htbp]
  \centering
  \includegraphics[width=\linewidth]{assets/figure_p109_16.png}
\end{figure}

\begin{figure}[htbp]
  \centering
  \includegraphics[width=\linewidth]{assets/figure_p109_17.png}
\end{figure}

(5) 将各个对应的乘积项相加,写出最简与或式.

\[ F(A,B,C,D)=ABCD+ACD+ABD+ABC \]

\[ 例： F(A,B,C,D)=AC+ACD+ABD+BC+BCD \]

\textbf{CD}
\textbf{AB 00 01 11 10}

\textbf{00 1 1 1}

\textbf{01 1 1 1}
\textbf{1 1 F(A,B,C,D)=ABD+BD}
\textbf{11}

\textbf{+AD+CD}

\textbf{10 1 1 1}

\begin{figure}[htbp]
  \centering
  \includegraphics[width=\linewidth]{assets/figure_p110_1.png}
\end{figure}

\begin{figure}[htbp]
  \centering
  \includegraphics[width=\linewidth]{assets/figure_p110_2.png}
\end{figure}

\begin{figure}[htbp]
  \centering
  \includegraphics[width=\linewidth]{assets/figure_p110_3.png}
\end{figure}

\begin{figure}[htbp]
  \centering
  \includegraphics[width=\linewidth]{assets/figure_p110_4.png}
\end{figure}

\begin{figure}[htbp]
  \centering
  \includegraphics[width=\linewidth]{assets/figure_p110_5.png}
\end{figure}

\begin{figure}[htbp]
  \centering
  \includegraphics[width=\linewidth]{assets/figure_p110_6.png}
\end{figure}

\begin{figure}[htbp]
  \centering
  \includegraphics[width=\linewidth]{assets/figure_p110_7.png}
\end{figure}

\begin{figure}[htbp]
  \centering
  \includegraphics[width=\linewidth]{assets/figure_p110_8.png}
\end{figure}

\begin{figure}[htbp]
  \centering
  \includegraphics[width=\linewidth]{assets/figure_p110_9.png}
\end{figure}

\begin{figure}[htbp]
  \centering
  \includegraphics[width=\linewidth]{assets/figure_p110_10.png}
\end{figure}

\begin{figure}[htbp]
  \centering
  \includegraphics[width=\linewidth]{assets/figure_p110_11.png}
\end{figure}

\begin{figure}[htbp]
  \centering
  \includegraphics[width=\linewidth]{assets/figure_p110_12.png}
\end{figure}

\begin{figure}[htbp]
  \centering
  \includegraphics[width=\linewidth]{assets/figure_p110_13.png}
\end{figure}

\begin{figure}[htbp]
  \centering
  \includegraphics[width=\linewidth]{assets/figure_p110_14.png}
\end{figure}

\begin{figure}[htbp]
  \centering
  \includegraphics[width=\linewidth]{assets/figure_p110_15.png}
\end{figure}

\begin{figure}[htbp]
  \centering
  \includegraphics[width=\linewidth]{assets/figure_p110_16.png}
\end{figure}

\begin{figure}[htbp]
  \centering
  \includegraphics[width=\linewidth]{assets/figure_p110_17.png}
\end{figure}

一种特殊情况：

\textbf{C C}
\textbf{AB 0 1 AB 0 1}
\textbf{00 1 00 1}
\textbf{01 1 1 01 1 1}
\textbf{11 1 11 1}
\textbf{10 1 1 10 1 1}

\[ F=AB+BC+AC F=AB+BC+AC \]

得到两种化简结果，也都是最简的。

\begin{figure}[htbp]
  \centering
  \includegraphics[width=\linewidth]{assets/figure_p111_1.png}
\end{figure}

\begin{figure}[htbp]
  \centering
  \includegraphics[width=\linewidth]{assets/figure_p111_2.png}
\end{figure}

\begin{figure}[htbp]
  \centering
  \includegraphics[width=\linewidth]{assets/figure_p111_3.png}
\end{figure}

\begin{figure}[htbp]
  \centering
  \includegraphics[width=\linewidth]{assets/figure_p111_4.png}
\end{figure}

\begin{figure}[htbp]
  \centering
  \includegraphics[width=\linewidth]{assets/figure_p111_5.png}
\end{figure}

\begin{figure}[htbp]
  \centering
  \includegraphics[width=\linewidth]{assets/figure_p111_6.png}
\end{figure}

\begin{figure}[htbp]
  \centering
  \includegraphics[width=\linewidth]{assets/figure_p111_7.png}
\end{figure}

\begin{figure}[htbp]
  \centering
  \includegraphics[width=\linewidth]{assets/figure_p111_8.png}
\end{figure}

\begin{figure}[htbp]
  \centering
  \includegraphics[width=\linewidth]{assets/figure_p111_9.png}
\end{figure}

\begin{figure}[htbp]
  \centering
  \includegraphics[width=\linewidth]{assets/figure_p111_10.png}
\end{figure}

\begin{figure}[htbp]
  \centering
  \includegraphics[width=\linewidth]{assets/figure_p111_11.png}
\end{figure}

\begin{figure}[htbp]
  \centering
  \includegraphics[width=\linewidth]{assets/figure_p111_12.png}
\end{figure}

\begin{figure}[htbp]
  \centering
  \includegraphics[width=\linewidth]{assets/figure_p111_13.png}
\end{figure}

\begin{figure}[htbp]
  \centering
  \includegraphics[width=\linewidth]{assets/figure_p111_14.png}
\end{figure}

\begin{figure}[htbp]
  \centering
  \includegraphics[width=\linewidth]{assets/figure_p111_15.png}
\end{figure}

\begin{figure}[htbp]
  \centering
  \includegraphics[width=\linewidth]{assets/figure_p111_16.png}
\end{figure}

\begin{figure}[htbp]
  \centering
  \includegraphics[width=\linewidth]{assets/figure_p111_17.png}
\end{figure}

③ 化简中注意的问题
(1) 每一个标1的方格必须至少被圈一次;
(2) 每个圈中包含的相邻小方格数,必须为2的整数次幂;
(3) 为了得到尽可能大的圈,圈与圈之间可以重叠;

\textbf{CD}
\textbf{AB 00 01 11 10}

\textbf{00 1 1 1}

\textbf{01 1 1 1}
\textbf{11 1 1}

\textbf{10 1 1 1}

\begin{figure}[htbp]
  \centering
  \includegraphics[width=\linewidth]{assets/figure_p112_1.png}
\end{figure}

\begin{figure}[htbp]
  \centering
  \includegraphics[width=\linewidth]{assets/figure_p112_2.png}
\end{figure}

\begin{figure}[htbp]
  \centering
  \includegraphics[width=\linewidth]{assets/figure_p112_3.png}
\end{figure}

\begin{figure}[htbp]
  \centering
  \includegraphics[width=\linewidth]{assets/figure_p112_4.png}
\end{figure}

\begin{figure}[htbp]
  \centering
  \includegraphics[width=\linewidth]{assets/figure_p112_5.png}
\end{figure}

\begin{figure}[htbp]
  \centering
  \includegraphics[width=\linewidth]{assets/figure_p112_6.png}
\end{figure}

\begin{figure}[htbp]
  \centering
  \includegraphics[width=\linewidth]{assets/figure_p112_7.png}
\end{figure}

\begin{figure}[htbp]
  \centering
  \includegraphics[width=\linewidth]{assets/figure_p112_8.png}
\end{figure}

\begin{figure}[htbp]
  \centering
  \includegraphics[width=\linewidth]{assets/figure_p112_9.png}
\end{figure}

\begin{figure}[htbp]
  \centering
  \includegraphics[width=\linewidth]{assets/figure_p112_10.png}
\end{figure}

\begin{figure}[htbp]
  \centering
  \includegraphics[width=\linewidth]{assets/figure_p112_11.png}
\end{figure}

\begin{figure}[htbp]
  \centering
  \includegraphics[width=\linewidth]{assets/figure_p112_12.png}
\end{figure}

\begin{figure}[htbp]
  \centering
  \includegraphics[width=\linewidth]{assets/figure_p112_13.png}
\end{figure}

\begin{figure}[htbp]
  \centering
  \includegraphics[width=\linewidth]{assets/figure_p112_14.png}
\end{figure}

\begin{figure}[htbp]
  \centering
  \includegraphics[width=\linewidth]{assets/figure_p112_15.png}
\end{figure}

\begin{figure}[htbp]
  \centering
  \includegraphics[width=\linewidth]{assets/figure_p112_16.png}
\end{figure}

\begin{figure}[htbp]
  \centering
  \includegraphics[width=\linewidth]{assets/figure_p112_17.png}
\end{figure}

(4) 若某个圈中的所有标1方格,已经完全被其它圈所
覆盖,则该圈为多余的.
例如: \textbf{CD}

\textbf{AB 00 01 11 10}
蓝色的圈为多余的. \textbf{00 1}

\textbf{01 1 1 1}

\[ F=ABC+ACD+ACD+ABC \]

\textbf{11 1 1 1}

\textbf{+ (BD)}

\textbf{10 1}

\begin{figure}[htbp]
  \centering
  \includegraphics[width=\linewidth]{assets/figure_p113_1.png}
\end{figure}

\begin{figure}[htbp]
  \centering
  \includegraphics[width=\linewidth]{assets/figure_p113_2.png}
\end{figure}

\begin{figure}[htbp]
  \centering
  \includegraphics[width=\linewidth]{assets/figure_p113_3.png}
\end{figure}

\begin{figure}[htbp]
  \centering
  \includegraphics[width=\linewidth]{assets/figure_p113_4.png}
\end{figure}

\begin{figure}[htbp]
  \centering
  \includegraphics[width=\linewidth]{assets/figure_p113_5.png}
\end{figure}

\begin{figure}[htbp]
  \centering
  \includegraphics[width=\linewidth]{assets/figure_p113_6.png}
\end{figure}

\begin{figure}[htbp]
  \centering
  \includegraphics[width=\linewidth]{assets/figure_p113_7.png}
\end{figure}

\begin{figure}[htbp]
  \centering
  \includegraphics[width=\linewidth]{assets/figure_p113_8.png}
\end{figure}

\begin{figure}[htbp]
  \centering
  \includegraphics[width=\linewidth]{assets/figure_p113_9.png}
\end{figure}

\begin{figure}[htbp]
  \centering
  \includegraphics[width=\linewidth]{assets/figure_p113_10.png}
\end{figure}

\begin{figure}[htbp]
  \centering
  \includegraphics[width=\linewidth]{assets/figure_p113_11.png}
\end{figure}

\begin{figure}[htbp]
  \centering
  \includegraphics[width=\linewidth]{assets/figure_p113_12.png}
\end{figure}

\begin{figure}[htbp]
  \centering
  \includegraphics[width=\linewidth]{assets/figure_p113_13.png}
\end{figure}

\begin{figure}[htbp]
  \centering
  \includegraphics[width=\linewidth]{assets/figure_p113_14.png}
\end{figure}

\begin{figure}[htbp]
  \centering
  \includegraphics[width=\linewidth]{assets/figure_p113_15.png}
\end{figure}

\begin{figure}[htbp]
  \centering
  \includegraphics[width=\linewidth]{assets/figure_p113_16.png}
\end{figure}

\begin{figure}[htbp]
  \centering
  \includegraphics[width=\linewidth]{assets/figure_p113_17.png}
\end{figure}

④ 用卡诺图求反函数的最简与或式

方法:在卡诺图中合并标 0 方格,可得到反函数的最简

与或式.

例\textbf{:}

\textbf{BC}
\textbf{A 00 01 11 10}
\textbf{0 0 0 1 0}

\[ F=AB+BC+AC \]

\textbf{1 0 1 1 1}

\begin{figure}[htbp]
  \centering
  \includegraphics[width=\linewidth]{assets/figure_p114_1.png}
\end{figure}

\begin{figure}[htbp]
  \centering
  \includegraphics[width=\linewidth]{assets/figure_p114_2.png}
\end{figure}

\begin{figure}[htbp]
  \centering
  \includegraphics[width=\linewidth]{assets/figure_p114_3.png}
\end{figure}

\begin{figure}[htbp]
  \centering
  \includegraphics[width=\linewidth]{assets/figure_p114_4.png}
\end{figure}

\begin{figure}[htbp]
  \centering
  \includegraphics[width=\linewidth]{assets/figure_p114_5.png}
\end{figure}

\begin{figure}[htbp]
  \centering
  \includegraphics[width=\linewidth]{assets/figure_p114_6.png}
\end{figure}

\begin{figure}[htbp]
  \centering
  \includegraphics[width=\linewidth]{assets/figure_p114_7.png}
\end{figure}

\begin{figure}[htbp]
  \centering
  \includegraphics[width=\linewidth]{assets/figure_p114_8.png}
\end{figure}

\begin{figure}[htbp]
  \centering
  \includegraphics[width=\linewidth]{assets/figure_p114_9.png}
\end{figure}

\begin{figure}[htbp]
  \centering
  \includegraphics[width=\linewidth]{assets/figure_p114_10.png}
\end{figure}

\begin{figure}[htbp]
  \centering
  \includegraphics[width=\linewidth]{assets/figure_p114_11.png}
\end{figure}

\begin{figure}[htbp]
  \centering
  \includegraphics[width=\linewidth]{assets/figure_p114_12.png}
\end{figure}

\begin{figure}[htbp]
  \centering
  \includegraphics[width=\linewidth]{assets/figure_p114_13.png}
\end{figure}

\begin{figure}[htbp]
  \centering
  \includegraphics[width=\linewidth]{assets/figure_p114_14.png}
\end{figure}

\begin{figure}[htbp]
  \centering
  \includegraphics[width=\linewidth]{assets/figure_p114_15.png}
\end{figure}

\begin{figure}[htbp]
  \centering
  \includegraphics[width=\linewidth]{assets/figure_p114_16.png}
\end{figure}

\begin{figure}[htbp]
  \centering
  \includegraphics[width=\linewidth]{assets/figure_p114_17.png}
\end{figure}

•常利用该方法来求逻辑函数F的最简与或非式, 例如将上
式F上 的非号移到右边,就得到F的最简与或非表达式.

\[ F=AB+BC+AC \]

\begin{figure}[htbp]
  \centering
  \includegraphics[width=\linewidth]{assets/figure_p115_1.png}
\end{figure}

\begin{figure}[htbp]
  \centering
  \includegraphics[width=\linewidth]{assets/figure_p115_2.png}
\end{figure}

\begin{figure}[htbp]
  \centering
  \includegraphics[width=\linewidth]{assets/figure_p115_3.png}
\end{figure}

\begin{figure}[htbp]
  \centering
  \includegraphics[width=\linewidth]{assets/figure_p115_4.png}
\end{figure}

\begin{figure}[htbp]
  \centering
  \includegraphics[width=\linewidth]{assets/figure_p115_5.png}
\end{figure}

\begin{figure}[htbp]
  \centering
  \includegraphics[width=\linewidth]{assets/figure_p115_6.png}
\end{figure}

\begin{figure}[htbp]
  \centering
  \includegraphics[width=\linewidth]{assets/figure_p115_7.png}
\end{figure}

\begin{figure}[htbp]
  \centering
  \includegraphics[width=\linewidth]{assets/figure_p115_8.png}
\end{figure}

\begin{figure}[htbp]
  \centering
  \includegraphics[width=\linewidth]{assets/figure_p115_9.png}
\end{figure}

\begin{figure}[htbp]
  \centering
  \includegraphics[width=\linewidth]{assets/figure_p115_10.png}
\end{figure}

\begin{figure}[htbp]
  \centering
  \includegraphics[width=\linewidth]{assets/figure_p115_11.png}
\end{figure}

\begin{figure}[htbp]
  \centering
  \includegraphics[width=\linewidth]{assets/figure_p115_12.png}
\end{figure}

\begin{figure}[htbp]
  \centering
  \includegraphics[width=\linewidth]{assets/figure_p115_13.png}
\end{figure}

\begin{figure}[htbp]
  \centering
  \includegraphics[width=\linewidth]{assets/figure_p115_14.png}
\end{figure}

\begin{figure}[htbp]
  \centering
  \includegraphics[width=\linewidth]{assets/figure_p115_15.png}
\end{figure}

\begin{figure}[htbp]
  \centering
  \includegraphics[width=\linewidth]{assets/figure_p115_16.png}
\end{figure}

\begin{figure}[htbp]
  \centering
  \includegraphics[width=\linewidth]{assets/figure_p115_17.png}
\end{figure}

• 对较为复杂的逻辑函数，可将函数分解成
多个部分，先将每个部分分别填入各自的
卡诺图中，然后通过卡诺图对应方格的运
算，求出函数的卡诺图。
• 对卡诺图进行化简。

\begin{figure}[htbp]
  \centering
  \includegraphics[width=\linewidth]{assets/figure_p116_1.png}
\end{figure}

\begin{figure}[htbp]
  \centering
  \includegraphics[width=\linewidth]{assets/figure_p116_2.png}
\end{figure}

\begin{figure}[htbp]
  \centering
  \includegraphics[width=\linewidth]{assets/figure_p116_3.png}
\end{figure}

\begin{figure}[htbp]
  \centering
  \includegraphics[width=\linewidth]{assets/figure_p116_4.png}
\end{figure}

\begin{figure}[htbp]
  \centering
  \includegraphics[width=\linewidth]{assets/figure_p116_5.png}
\end{figure}

\begin{figure}[htbp]
  \centering
  \includegraphics[width=\linewidth]{assets/figure_p116_6.png}
\end{figure}

\begin{figure}[htbp]
  \centering
  \includegraphics[width=\linewidth]{assets/figure_p116_7.png}
\end{figure}

\begin{figure}[htbp]
  \centering
  \includegraphics[width=\linewidth]{assets/figure_p116_8.png}
\end{figure}

\begin{figure}[htbp]
  \centering
  \includegraphics[width=\linewidth]{assets/figure_p116_9.png}
\end{figure}

\begin{figure}[htbp]
  \centering
  \includegraphics[width=\linewidth]{assets/figure_p116_10.png}
\end{figure}

\begin{figure}[htbp]
  \centering
  \includegraphics[width=\linewidth]{assets/figure_p116_11.png}
\end{figure}

\begin{figure}[htbp]
  \centering
  \includegraphics[width=\linewidth]{assets/figure_p116_12.png}
\end{figure}

\begin{figure}[htbp]
  \centering
  \includegraphics[width=\linewidth]{assets/figure_p116_13.png}
\end{figure}

\begin{figure}[htbp]
  \centering
  \includegraphics[width=\linewidth]{assets/figure_p116_14.png}
\end{figure}

\begin{figure}[htbp]
  \centering
  \includegraphics[width=\linewidth]{assets/figure_p116_15.png}
\end{figure}

\begin{figure}[htbp]
  \centering
  \includegraphics[width=\linewidth]{assets/figure_p116_16.png}
\end{figure}

\begin{figure}[htbp]
  \centering
  \includegraphics[width=\linewidth]{assets/figure_p116_17.png}
\end{figure}

\begin{figure}[htbp]
  \centering
  \includegraphics[width=\linewidth]{assets/figure_p116_18.png}
\end{figure}

例：化简逻辑函数
F=(AB+AC+BD)(ABCD+ACD+BCD+BC)

CD CD
CD
AB 00 01 11 10 AB 00 01 11 10
AB 00 01 11 10
00 1 1 1 00 1 1
00 1
01 1 1 01 1
= 01 1



11 1 1 1 1 11
1 11 1 1 1
10 1 1 10 1 1 1
10 1

\[ F=ABCD+ABC+BCD+ACD \]

\begin{figure}[htbp]
  \centering
  \includegraphics[width=\linewidth]{assets/figure_p117_1.png}
\end{figure}

\begin{figure}[htbp]
  \centering
  \includegraphics[width=\linewidth]{assets/figure_p117_2.png}
\end{figure}

\begin{figure}[htbp]
  \centering
  \includegraphics[width=\linewidth]{assets/figure_p117_3.png}
\end{figure}

\begin{figure}[htbp]
  \centering
  \includegraphics[width=\linewidth]{assets/figure_p117_4.png}
\end{figure}

\begin{figure}[htbp]
  \centering
  \includegraphics[width=\linewidth]{assets/figure_p117_5.png}
\end{figure}

\begin{figure}[htbp]
  \centering
  \includegraphics[width=\linewidth]{assets/figure_p117_6.png}
\end{figure}

\begin{figure}[htbp]
  \centering
  \includegraphics[width=\linewidth]{assets/figure_p117_7.png}
\end{figure}

\begin{figure}[htbp]
  \centering
  \includegraphics[width=\linewidth]{assets/figure_p117_8.png}
\end{figure}

\begin{figure}[htbp]
  \centering
  \includegraphics[width=\linewidth]{assets/figure_p117_9.png}
\end{figure}

\begin{figure}[htbp]
  \centering
  \includegraphics[width=\linewidth]{assets/figure_p117_10.png}
\end{figure}

\begin{figure}[htbp]
  \centering
  \includegraphics[width=\linewidth]{assets/figure_p117_11.png}
\end{figure}

\begin{figure}[htbp]
  \centering
  \includegraphics[width=\linewidth]{assets/figure_p117_12.png}
\end{figure}

\begin{figure}[htbp]
  \centering
  \includegraphics[width=\linewidth]{assets/figure_p117_13.png}
\end{figure}

\begin{figure}[htbp]
  \centering
  \includegraphics[width=\linewidth]{assets/figure_p117_14.png}
\end{figure}

\begin{figure}[htbp]
  \centering
  \includegraphics[width=\linewidth]{assets/figure_p117_15.png}
\end{figure}

\begin{figure}[htbp]
  \centering
  \includegraphics[width=\linewidth]{assets/figure_p117_16.png}
\end{figure}

\begin{figure}[htbp]
  \centering
  \includegraphics[width=\linewidth]{assets/figure_p117_17.png}
\end{figure}

5) 不完全确定的逻辑函数及其化简

在某些实际数字电路中,逻辑函数的输出只和一部分

最小项有确定对应关系,而和余下的最小项无关.余下的
最小项无论写入逻辑函数式还是不写入逻辑函数式,都不
影响电路的逻辑功能.把这些最小项称为无关项.用英文
字母d(don\textbf{’}t $care)表示，对应的函数值记为“\times \textbf{”} 。$

包含无关项的逻辑函数称为不完全确定的逻辑函数.

\begin{figure}[htbp]
  \centering
  \includegraphics[width=\linewidth]{assets/figure_p118_1.png}
\end{figure}

\begin{figure}[htbp]
  \centering
  \includegraphics[width=\linewidth]{assets/figure_p118_2.png}
\end{figure}

\begin{figure}[htbp]
  \centering
  \includegraphics[width=\linewidth]{assets/figure_p118_3.png}
\end{figure}

\begin{figure}[htbp]
  \centering
  \includegraphics[width=\linewidth]{assets/figure_p118_4.png}
\end{figure}

\begin{figure}[htbp]
  \centering
  \includegraphics[width=\linewidth]{assets/figure_p118_5.png}
\end{figure}

\begin{figure}[htbp]
  \centering
  \includegraphics[width=\linewidth]{assets/figure_p118_6.png}
\end{figure}

\begin{figure}[htbp]
  \centering
  \includegraphics[width=\linewidth]{assets/figure_p118_7.png}
\end{figure}

\begin{figure}[htbp]
  \centering
  \includegraphics[width=\linewidth]{assets/figure_p118_8.png}
\end{figure}

\begin{figure}[htbp]
  \centering
  \includegraphics[width=\linewidth]{assets/figure_p118_9.png}
\end{figure}

\begin{figure}[htbp]
  \centering
  \includegraphics[width=\linewidth]{assets/figure_p118_10.png}
\end{figure}

\begin{figure}[htbp]
  \centering
  \includegraphics[width=\linewidth]{assets/figure_p118_11.png}
\end{figure}

\begin{figure}[htbp]
  \centering
  \includegraphics[width=\linewidth]{assets/figure_p118_12.png}
\end{figure}

\begin{figure}[htbp]
  \centering
  \includegraphics[width=\linewidth]{assets/figure_p118_13.png}
\end{figure}

\begin{figure}[htbp]
  \centering
  \includegraphics[width=\linewidth]{assets/figure_p118_14.png}
\end{figure}

\begin{figure}[htbp]
  \centering
  \includegraphics[width=\linewidth]{assets/figure_p118_15.png}
\end{figure}

\begin{figure}[htbp]
  \centering
  \includegraphics[width=\linewidth]{assets/figure_p118_16.png}
\end{figure}

\begin{figure}[htbp]
  \centering
  \includegraphics[width=\linewidth]{assets/figure_p118_17.png}
\end{figure}

利用不完全确定的逻辑函数中的无关项往往可以将函

数化得更简单.
例: 设计一个奇偶判别电路.电路输入为8421\textbf{BCD}码,当

输入为偶数时,输出为 0 ;当电路输入为奇数时,输出为1 .

由于8421\textbf{BCD}码中无\textbf{1010\textasciitilde{}1111}这6个码,电路禁止输

入这6个码.这6个码对应的最小项为无关项.

\begin{figure}[htbp]
  \centering
  \includegraphics[width=\linewidth]{assets/figure_p119_1.png}
\end{figure}

\begin{figure}[htbp]
  \centering
  \includegraphics[width=\linewidth]{assets/figure_p119_2.png}
\end{figure}

\begin{figure}[htbp]
  \centering
  \includegraphics[width=\linewidth]{assets/figure_p119_3.png}
\end{figure}

\begin{figure}[htbp]
  \centering
  \includegraphics[width=\linewidth]{assets/figure_p119_4.png}
\end{figure}

\begin{figure}[htbp]
  \centering
  \includegraphics[width=\linewidth]{assets/figure_p119_5.png}
\end{figure}

\begin{figure}[htbp]
  \centering
  \includegraphics[width=\linewidth]{assets/figure_p119_6.png}
\end{figure}

\begin{figure}[htbp]
  \centering
  \includegraphics[width=\linewidth]{assets/figure_p119_7.png}
\end{figure}

\begin{figure}[htbp]
  \centering
  \includegraphics[width=\linewidth]{assets/figure_p119_8.png}
\end{figure}

\begin{figure}[htbp]
  \centering
  \includegraphics[width=\linewidth]{assets/figure_p119_9.png}
\end{figure}

\begin{figure}[htbp]
  \centering
  \includegraphics[width=\linewidth]{assets/figure_p119_10.png}
\end{figure}

\begin{figure}[htbp]
  \centering
  \includegraphics[width=\linewidth]{assets/figure_p119_11.png}
\end{figure}

\begin{figure}[htbp]
  \centering
  \includegraphics[width=\linewidth]{assets/figure_p119_12.png}
\end{figure}

\begin{figure}[htbp]
  \centering
  \includegraphics[width=\linewidth]{assets/figure_p119_13.png}
\end{figure}

\begin{figure}[htbp]
  \centering
  \includegraphics[width=\linewidth]{assets/figure_p119_14.png}
\end{figure}

\begin{figure}[htbp]
  \centering
  \includegraphics[width=\linewidth]{assets/figure_p119_15.png}
\end{figure}

\begin{figure}[htbp]
  \centering
  \includegraphics[width=\linewidth]{assets/figure_p119_16.png}
\end{figure}

\begin{figure}[htbp]
  \centering
  \includegraphics[width=\linewidth]{assets/figure_p119_17.png}
\end{figure}

真值表

\textbf{A B C D F A B C D F}
\textbf{A 0 0 0 0 0 1 0 0 0 0}
奇偶
\textbf{B F 0 0 0 1 1 1 0 0 1 1}
判别
\textbf{C $0 0 1 0 0 1 0 1 0} \times$
电路
\textbf{D $0 0 1 1 1 1 0 1 1} \times$
\textbf{0 $1 0 0 0 1 1 0 0} \times$
\textbf{0 $1 0 1 1 1 1 0 1} \times$
\textbf{0 $1 1 0 0 1 1 1 0} \times$
\textbf{0 $1 1 1 1 1 1 1 1} \times$

$\textbf{F(A,B,C,D)=\Sigma m(1,3,5,7,9)+\Sigma d(10$ \textasciitilde{} 15)}

\begin{figure}[htbp]
  \centering
  \includegraphics[width=\linewidth]{assets/figure_p120_1.png}
\end{figure}

\begin{figure}[htbp]
  \centering
  \includegraphics[width=\linewidth]{assets/figure_p120_2.png}
\end{figure}

\begin{figure}[htbp]
  \centering
  \includegraphics[width=\linewidth]{assets/figure_p120_3.png}
\end{figure}

\begin{figure}[htbp]
  \centering
  \includegraphics[width=\linewidth]{assets/figure_p120_4.png}
\end{figure}

\begin{figure}[htbp]
  \centering
  \includegraphics[width=\linewidth]{assets/figure_p120_5.png}
\end{figure}

\begin{figure}[htbp]
  \centering
  \includegraphics[width=\linewidth]{assets/figure_p120_6.png}
\end{figure}

\begin{figure}[htbp]
  \centering
  \includegraphics[width=\linewidth]{assets/figure_p120_7.png}
\end{figure}

\begin{figure}[htbp]
  \centering
  \includegraphics[width=\linewidth]{assets/figure_p120_8.png}
\end{figure}

\begin{figure}[htbp]
  \centering
  \includegraphics[width=\linewidth]{assets/figure_p120_9.png}
\end{figure}

\begin{figure}[htbp]
  \centering
  \includegraphics[width=\linewidth]{assets/figure_p120_10.png}
\end{figure}

\begin{figure}[htbp]
  \centering
  \includegraphics[width=\linewidth]{assets/figure_p120_11.png}
\end{figure}

\begin{figure}[htbp]
  \centering
  \includegraphics[width=\linewidth]{assets/figure_p120_12.png}
\end{figure}

\begin{figure}[htbp]
  \centering
  \includegraphics[width=\linewidth]{assets/figure_p120_13.png}
\end{figure}

\begin{figure}[htbp]
  \centering
  \includegraphics[width=\linewidth]{assets/figure_p120_14.png}
\end{figure}

\begin{figure}[htbp]
  \centering
  \includegraphics[width=\linewidth]{assets/figure_p120_15.png}
\end{figure}

\begin{figure}[htbp]
  \centering
  \includegraphics[width=\linewidth]{assets/figure_p120_16.png}
\end{figure}

\begin{figure}[htbp]
  \centering
  \includegraphics[width=\linewidth]{assets/figure_p120_17.png}
\end{figure}

$\textbf{F(A,B,C,D)=\Sigma m(1,3,5,7,9)+\Sigma d(10$ \textasciitilde{} 15)}

\textbf{CD}
若不利用无关项（即将卡诺 \textbf{AB 00 01 11 10}
$图中的\times 均作0处理）,则化$ \textbf{00 0 1 1 0}
简结果为: \textbf{01 0 1 1 0}
\textbf{11} $\times \times \times \times$
\textbf{F(A,B,C,D)=AD+BCD $0 1} \times \times$
\textbf{10}

$若利用无关项(即将卡诺图中的\times 按化简的需要任意$
$处理,将有些\times 当作0,有些\times 当作1),则化简结果为:$

\textbf{F(A,B,C,D)=D}

\begin{figure}[htbp]
  \centering
  \includegraphics[width=\linewidth]{assets/figure_p121_1.png}
\end{figure}

\begin{figure}[htbp]
  \centering
  \includegraphics[width=\linewidth]{assets/figure_p121_2.png}
\end{figure}

\begin{figure}[htbp]
  \centering
  \includegraphics[width=\linewidth]{assets/figure_p121_3.png}
\end{figure}

\begin{figure}[htbp]
  \centering
  \includegraphics[width=\linewidth]{assets/figure_p121_4.png}
\end{figure}

\begin{figure}[htbp]
  \centering
  \includegraphics[width=\linewidth]{assets/figure_p121_5.png}
\end{figure}

\begin{figure}[htbp]
  \centering
  \includegraphics[width=\linewidth]{assets/figure_p121_6.png}
\end{figure}

\begin{figure}[htbp]
  \centering
  \includegraphics[width=\linewidth]{assets/figure_p121_7.png}
\end{figure}

\begin{figure}[htbp]
  \centering
  \includegraphics[width=\linewidth]{assets/figure_p121_8.png}
\end{figure}

\begin{figure}[htbp]
  \centering
  \includegraphics[width=\linewidth]{assets/figure_p121_9.png}
\end{figure}

\begin{figure}[htbp]
  \centering
  \includegraphics[width=\linewidth]{assets/figure_p121_10.png}
\end{figure}

\begin{figure}[htbp]
  \centering
  \includegraphics[width=\linewidth]{assets/figure_p121_11.png}
\end{figure}

\begin{figure}[htbp]
  \centering
  \includegraphics[width=\linewidth]{assets/figure_p121_12.png}
\end{figure}

\begin{figure}[htbp]
  \centering
  \includegraphics[width=\linewidth]{assets/figure_p121_13.png}
\end{figure}

\begin{figure}[htbp]
  \centering
  \includegraphics[width=\linewidth]{assets/figure_p121_14.png}
\end{figure}

\begin{figure}[htbp]
  \centering
  \includegraphics[width=\linewidth]{assets/figure_p121_15.png}
\end{figure}

\begin{figure}[htbp]
  \centering
  \includegraphics[width=\linewidth]{assets/figure_p121_16.png}
\end{figure}

\begin{figure}[htbp]
  \centering
  \includegraphics[width=\linewidth]{assets/figure_p121_17.png}
\end{figure}

完整地将函数写为：

\textbf{F(A,B,C,D)=D}
$\textbf{\Sigma d(10$ \textasciitilde{} 15)=0}

再例：

\[ F=(A  B)CD+ABC+ACD 且 AB+CD=0 \]

\textbf{CD}
\textbf{AB 00 01 11 10}
\textbf{00 $0 1} \times$ \textbf{0}

\[ F(A,B,C,D)=B+AD+AC \]

\textbf{01 $1 1} \times$ \textbf{1}
\textbf{11} $\times \times \times \times$
\textbf{10 $0 0} \times$ \textbf{1}

\begin{figure}[htbp]
  \centering
  \includegraphics[width=\linewidth]{assets/figure_p122_1.png}
\end{figure}

\begin{figure}[htbp]
  \centering
  \includegraphics[width=\linewidth]{assets/figure_p122_2.png}
\end{figure}

\begin{figure}[htbp]
  \centering
  \includegraphics[width=\linewidth]{assets/figure_p122_3.png}
\end{figure}

\begin{figure}[htbp]
  \centering
  \includegraphics[width=\linewidth]{assets/figure_p122_4.png}
\end{figure}

\begin{figure}[htbp]
  \centering
  \includegraphics[width=\linewidth]{assets/figure_p122_5.png}
\end{figure}

\begin{figure}[htbp]
  \centering
  \includegraphics[width=\linewidth]{assets/figure_p122_6.png}
\end{figure}

\begin{figure}[htbp]
  \centering
  \includegraphics[width=\linewidth]{assets/figure_p122_7.png}
\end{figure}

\begin{figure}[htbp]
  \centering
  \includegraphics[width=\linewidth]{assets/figure_p122_8.png}
\end{figure}

\begin{figure}[htbp]
  \centering
  \includegraphics[width=\linewidth]{assets/figure_p122_9.png}
\end{figure}

\begin{figure}[htbp]
  \centering
  \includegraphics[width=\linewidth]{assets/figure_p122_10.png}
\end{figure}

\begin{figure}[htbp]
  \centering
  \includegraphics[width=\linewidth]{assets/figure_p122_11.png}
\end{figure}

\begin{figure}[htbp]
  \centering
  \includegraphics[width=\linewidth]{assets/figure_p122_12.png}
\end{figure}

\begin{figure}[htbp]
  \centering
  \includegraphics[width=\linewidth]{assets/figure_p122_13.png}
\end{figure}

\begin{figure}[htbp]
  \centering
  \includegraphics[width=\linewidth]{assets/figure_p122_14.png}
\end{figure}

\begin{figure}[htbp]
  \centering
  \includegraphics[width=\linewidth]{assets/figure_p122_15.png}
\end{figure}

\begin{figure}[htbp]
  \centering
  \includegraphics[width=\linewidth]{assets/figure_p122_16.png}
\end{figure}

\begin{figure}[htbp]
  \centering
  \includegraphics[width=\linewidth]{assets/figure_p122_17.png}
\end{figure}

3. Q-M法（奎恩-麦克拉斯基）

Q-M法又称系统简化法，适用于任意多变量的函数，有
较严格的算法，可借助于计算机解题。

系统简化法分三个步骤进行：
(1) 求出函数的全部主要项；
(2) 选出函数的必要项；
(3) 选择主要项，使它与必要项一起包含给定函数的
全部最小项，建立函数的最简与－或式。

\begin{figure}[htbp]
  \centering
  \includegraphics[width=\linewidth]{assets/figure_p123_1.png}
\end{figure}

\begin{figure}[htbp]
  \centering
  \includegraphics[width=\linewidth]{assets/figure_p123_2.png}
\end{figure}

\begin{figure}[htbp]
  \centering
  \includegraphics[width=\linewidth]{assets/figure_p123_3.png}
\end{figure}

\begin{figure}[htbp]
  \centering
  \includegraphics[width=\linewidth]{assets/figure_p123_4.png}
\end{figure}

\begin{figure}[htbp]
  \centering
  \includegraphics[width=\linewidth]{assets/figure_p123_5.png}
\end{figure}

\begin{figure}[htbp]
  \centering
  \includegraphics[width=\linewidth]{assets/figure_p123_6.png}
\end{figure}

\begin{figure}[htbp]
  \centering
  \includegraphics[width=\linewidth]{assets/figure_p123_7.png}
\end{figure}

\begin{figure}[htbp]
  \centering
  \includegraphics[width=\linewidth]{assets/figure_p123_8.png}
\end{figure}

\begin{figure}[htbp]
  \centering
  \includegraphics[width=\linewidth]{assets/figure_p123_9.png}
\end{figure}

\begin{figure}[htbp]
  \centering
  \includegraphics[width=\linewidth]{assets/figure_p123_10.png}
\end{figure}

\begin{figure}[htbp]
  \centering
  \includegraphics[width=\linewidth]{assets/figure_p123_11.png}
\end{figure}

\begin{figure}[htbp]
  \centering
  \includegraphics[width=\linewidth]{assets/figure_p123_12.png}
\end{figure}

\begin{figure}[htbp]
  \centering
  \includegraphics[width=\linewidth]{assets/figure_p123_13.png}
\end{figure}

\begin{figure}[htbp]
  \centering
  \includegraphics[width=\linewidth]{assets/figure_p123_14.png}
\end{figure}

\begin{figure}[htbp]
  \centering
  \includegraphics[width=\linewidth]{assets/figure_p123_15.png}
\end{figure}

\begin{figure}[htbp]
  \centering
  \includegraphics[width=\linewidth]{assets/figure_p123_16.png}
\end{figure}

\begin{figure}[htbp]
  \centering
  \includegraphics[width=\linewidth]{assets/figure_p123_17.png}
\end{figure}

注意:在无特殊说明的情况下,为使逻辑函数化的更简单,
均应按上述第二种方法（卡诺图化简）处理最小项.

\begin{figure}[htbp]
  \centering
  \includegraphics[width=\linewidth]{assets/figure_p124_1.png}
\end{figure}

\begin{figure}[htbp]
  \centering
  \includegraphics[width=\linewidth]{assets/figure_p124_2.png}
\end{figure}

\begin{figure}[htbp]
  \centering
  \includegraphics[width=\linewidth]{assets/figure_p124_3.png}
\end{figure}

\begin{figure}[htbp]
  \centering
  \includegraphics[width=\linewidth]{assets/figure_p124_4.png}
\end{figure}

\begin{figure}[htbp]
  \centering
  \includegraphics[width=\linewidth]{assets/figure_p124_5.png}
\end{figure}

\begin{figure}[htbp]
  \centering
  \includegraphics[width=\linewidth]{assets/figure_p124_6.png}
\end{figure}

\begin{figure}[htbp]
  \centering
  \includegraphics[width=\linewidth]{assets/figure_p124_7.png}
\end{figure}

\begin{figure}[htbp]
  \centering
  \includegraphics[width=\linewidth]{assets/figure_p124_8.png}
\end{figure}

\begin{figure}[htbp]
  \centering
  \includegraphics[width=\linewidth]{assets/figure_p124_9.png}
\end{figure}

\begin{figure}[htbp]
  \centering
  \includegraphics[width=\linewidth]{assets/figure_p124_10.png}
\end{figure}

\begin{figure}[htbp]
  \centering
  \includegraphics[width=\linewidth]{assets/figure_p124_11.png}
\end{figure}

\begin{figure}[htbp]
  \centering
  \includegraphics[width=\linewidth]{assets/figure_p124_12.png}
\end{figure}

\begin{figure}[htbp]
  \centering
  \includegraphics[width=\linewidth]{assets/figure_p124_13.png}
\end{figure}

\begin{figure}[htbp]
  \centering
  \includegraphics[width=\linewidth]{assets/figure_p124_14.png}
\end{figure}

\begin{figure}[htbp]
  \centering
  \includegraphics[width=\linewidth]{assets/figure_p124_15.png}
\end{figure}

\begin{figure}[htbp]
  \centering
  \includegraphics[width=\linewidth]{assets/figure_p124_16.png}
\end{figure}

\begin{figure}[htbp]
  \centering
  \includegraphics[width=\linewidth]{assets/figure_p124_17.png}
\end{figure}

\begin{figure}[htbp]
  \centering
  \includegraphics[width=\linewidth]{assets/figure_p37_1.png}
\end{figure}

\begin{figure}[htbp]
  \centering
  \includegraphics[width=\linewidth]{assets/figure_p37_2.png}
\end{figure}

\begin{figure}[htbp]
  \centering
  \includegraphics[width=\linewidth]{assets/figure_p37_3.png}
\end{figure}

\begin{figure}[htbp]
  \centering
  \includegraphics[width=\linewidth]{assets/figure_p37_4.png}
\end{figure}

\begin{figure}[htbp]
  \centering
  \includegraphics[width=\linewidth]{assets/figure_p37_5.png}
\end{figure}

\begin{figure}[htbp]
  \centering
  \includegraphics[width=\linewidth]{assets/figure_p37_6.png}
\end{figure}

\begin{figure}[htbp]
  \centering
  \includegraphics[width=\linewidth]{assets/figure_p37_7.png}
\end{figure}

\begin{figure}[htbp]
  \centering
  \includegraphics[width=\linewidth]{assets/figure_p37_8.png}
\end{figure}

\begin{figure}[htbp]
  \centering
  \includegraphics[width=\linewidth]{assets/figure_p37_9.png}
\end{figure}

\begin{figure}[htbp]
  \centering
  \includegraphics[width=\linewidth]{assets/figure_p37_10.png}
\end{figure}

\begin{figure}[htbp]
  \centering
  \includegraphics[width=\linewidth]{assets/figure_p37_11.png}
\end{figure}

\begin{figure}[htbp]
  \centering
  \includegraphics[width=\linewidth]{assets/figure_p37_12.png}
\end{figure}

\begin{figure}[htbp]
  \centering
  \includegraphics[width=\linewidth]{assets/figure_p37_13.png}
\end{figure}

\begin{figure}[htbp]
  \centering
  \includegraphics[width=\linewidth]{assets/figure_p37_14.png}
\end{figure}

\begin{figure}[htbp]
  \centering
  \includegraphics[width=\linewidth]{assets/figure_p37_15.png}
\end{figure}

\begin{figure}[htbp]
  \centering
  \includegraphics[width=\linewidth]{assets/figure_p37_16.png}
\end{figure}

\begin{figure}[htbp]
  \centering
  \includegraphics[width=\linewidth]{assets/figure_p37_17.png}
\end{figure}

\begin{figure}[htbp]
  \centering
  \includegraphics[width=\linewidth]{assets/figure_p37_18.png}
\end{figure}

\end{document}
